\quellseite{102}
\chapter[C. Religious Instruction During the Reformation Period]{C. Religious Instruction During the Reformation Period and the Following Centuries}

J. W. Hoefling ii. * G. v. Zezschwitz i, ch. 24–30; ii, 1, ch. 12. 18. 25. 33. 34; ii, 2¹, ch. 14–19; ii, 2², ch. 31–47. Th. Harnack, pp. 66–105. Schumann and Sperber, \emph{Geschichte des Religionsunterrichts in der evgl. Volksschule}, pp. 23–155, 1890. E. Sachsse, pp. 186–300. E. Chr. Achelis, pp. 303–332 and 383–420. J. Steinbeck, pp. 26–61. Fr. Ehrenfeuchter, Zur \emph{Gesch. des Katechismus}, 1857. H. F. Th. L. Ernesti, \emph{Zur Orientierung ueber die Katechismusliteratur}, 1859. H. Heppe, \emph{Geschichte des deutschen Volksschulwesens}, 5 vols., 1858–1860. K. A. Schmid, \emph{Gesch. der Erziehung} ii 2 and the following volumes, 1889 ff. K. Hartfelder, \emph{Ph. Melanchthon als Praeceptor Germaniae}, 1889. *Fr. Paulsen, \emph{Gesch. des gelehrten Unterrichts vom Ausgang des Mittelalters bis zur Gegenwart}, 2 vols., ²1896. F. V. N. Painter, \emph{Luther on Education}. A. F. Leach, \emph{English Schools at the Reformation}, 1896. J. W. Richard, \emph{Melanchthon, the Protestant Preceptor of Germany}, 1898. Fr. Paulsen, \emph{German Education}, 1908. *F. P. Graves, \emph{History of Education during the Middle Ages and the Transition to Modern Times}, 1910. *F. W. Graves, \emph{History of Education in Modern Times}, 1913. Compare literature for § 13–15.

\quellseite{103}
\section[§ 13. The Work of the Reformatory Church]{§ 13. The Work of the Reformatory Church from 1517 to 1529}

G. Langemarck, \emph{Historia catechetica}, ii and iii, 1733, 1740. J. B. Riederer, \emph{Nachrichten zur Kirchen-, Gelehrten- und Buecher\-geschichte}, 1764–1768. G. Veesemeyer, \emph{Literarisch-bibliographische Nachrichten von einigen ev. Schriften u. Katechismen vor und nach Luthers Katechismen}, 1830. G. v. Zezschwitz ii, 1, ch. 36. G. v. Zezschwitz, \emph{Luthers Kleiner Katechismus, s. Bedeutung und s. Urgestalt}, 1880. G. Buchwald, \emph{Zur Wittenberger Stadt- und Universitaetsgeschichte in der Reformationszeit}, 1893. G. Buchwald, \emph{Stadtschreiber M. Steph. Roth in Zwickau} in Archiv f. Gesch. des deutschen Buchhandels, 1893. *G. Buchwald, \emph{Die Entstehung der Katechismen Luthers} 1894. M. Reu, \emph{Zur Entstehungsgesch. d. kl. Katech. Luthers} in Kirchl. Zeitchr., 1894. G. Rietschel, \emph{Luthers kl. Katech. in Tafelform} in Theol. Studien u. Kritiken, 1898. R. Neumann, \emph{Der evang. Religionsunterricht im Zeitalter der Reformation}, 1899. F. Cohrs, \emph{Katechismen Luthers in PRE³}, 1901. *Cohrs, \emph{Die evangelischen Katechismusversuche vor Luthers Enchiridion}, 5 vols., 1900–1904. *M. Reu, \emph{Quellen z. Geschichte d. kirchl. Unterrichts im evangel. Deutschland ii: Der biblische Unterricht}, 1906. O. Albrecht, Zur Bibliographie und Textkritik d. kl. Katech. in Archiv f. Ref. Gesch., 1903–1905. E. Sehling, \emph{Die ev. Kirchenordnungen d. 16. Jahrh.}, 5 vols., 1903 ff. O. Albrecht, Katechismusstudien in Theol. Studien u. Kritiken, 1907. 1908. 1909. Th. Kolde, \emph{Historische Einleitung in die Symbolischen Buecher d. luth. Kirche}, 1907. M. Reu, \emph{Quellen etc.}, vol. iii, 1911. F. Cohrs, \emph{Phil. Melanchthon in s. Bedeutung f. d. religioesen Jugendunterricht} in Theol. Studien u. Kritiken, 1912. J. Meyer, \emph{Luthers Kleiner Katechismus. Der deutsche Text i. s. geschichtl. Entwicklung}, 1912. A. Hardeland, \emph{Luthers Katechismusgedanken in ihrer Entwicklung bis z. Jahre 1529}, 1913. M. Reu Religious Instruction during the 16th Century in Luth. Church Review, 1916. *\emph{Luther's Works}, critical edition, 1883 ff.: Decem Praecepta Wittenbergensi praedicata populo, 1518, vol. i, pp. 394 ff.; Kurze Erklaerung d. zehn Gebote, 1518, vol. i, pp. 247 ff.; Auslegung deutsch d. Vaterunsers fuer d. einfaeltigen Laien, 1519, vol. ii, pp. 74 ff.; \emph{Eine kurze Form, das Paternoster zu verstehen u. zu beten}, 1519, vol. vi, pp. 9 ff.; Eine kurze u. gute Auslegung d. Vaterunsers fuer sich u. hinter sich, 1519, vol. vi, pp. 20 ff.; *Eine kurze Form d. zehn Gebote, eine kurze
\quellseite{104}
Form des Glaubens, eine kurze Form d. Vaterunsers, 1520, vol. vii, pp. 194 ff., Betbüchlein, 1522, vol. x, pp. 331 ff.; *\emph{An die Rathserrn aller Staede deutschen Landes, dass sie christl. Schulen aufrichten und erhalten sollen}, 1524, vol. xv, pp. 9 ff.; \emph{Deutsche Messe u. Ordnung d. Gottesdienstes}, 1526, vol. xix, pp. 53 ff.; \emph{Unterrichtung, wie man moege die Kinder fuehren zu Gottes Wort u. Dienst}, ibid., pp. 52, 61 f.; \emph{Vermahnung u. kurze Deutung d. Vaterunsers}, ibid., pp. 52 ff.; Was dem gemeinen Volk nach der Predigt vorzulesen, ibid., p. 52, 62 f.; *\emph{Katechismuspredigten}, 1528, vol. xxx, part i, pp. 1 ff.; *\emph{Deutsch Katechismus}, 1529, Text and various readings, vol. xxx, i, pp. 123 ff.; *\emph{\emph{Der kleine Katechismus}}, 1529, Text and various readings, vol. xxx, i, pp. 239 ff.; of great value is O. Albrecht's Introduction to the Two Catechisms, ibid., pp. 426 ff.; *\emph{Passionale}, 1529, vol. x, part ii, pp. 458 ff. (reprinted for the first time in Reu, Quellen ii, pp. 32 ff.). Luther's Small Catechism: C. Moenckeberg, \emph{Die erste Ausgabe von Luthers kleinem Katechismus in niedersaechsischer Uebersetzung} ²1868. K. F. Th. Schneider, \emph{Dr. M. \emph{Luthers kleiner Katechismus} nach den Originalausgaben kritisch bearbeitet} (edition of 1531), 1853. *Th. Harnack, \emph{Der kl. Katechismus Luthers in seiner Urgestalt} (Erfurt reprint of the original Wittenberg edition, 1529; enlarged Wittenberg edition, 1529; Wittenberg edition, 1539), 1856. H. Hartung, \emph{\emph{Der kleine Katechismus}} Luthers (Erfurt reprint, 1529), 1860. H. J. R. Calinich, Dr. M. \emph{Luthers kleiner Katechismus} (Wittenberg print, 1542), 1883. F. Goepfert, \emph{Woerterbuch zum kleinen Katechismus Luthers} (Wittenberg edition, 1537), 1889. A. Ebeling, \emph{Historisch-kritische Ausgabe von Luthers kleinem Katechismus}, 1901. K. Knoke, \emph{Ausgaben des Lutherschen Enchiridions bis zu Luthers Tod} und Neudruck d. (Wittenb. Ausg. von 1535, 1903. O. Albrecht, \emph{Luthers kleiner Katechismus} nach der Wittenberger Ausg. von 1540 in Jahrbuch der Akademie der Wissenschaften in Erfurt, 1904. K. Knoke, *M. \emph{Luthers kleiner Katechismus} nach den aeltesten Ausgaben in hoch-, niederdeutscher und lateinischer Sprache (6 different editions and translations, e. g., Marburg 1529 and Wittenberg 1543), 1904. O. Albrecht, Der kleine Kat. Luthers nach der Ausgabe von 1536, 1906. J. Meyer, \emph{Luthers grosser Katechismus. Textausgabe mit Kennzeichnung der Predigtgrundlagen}, 1914. *O. Albrecht, \emph{Luthers Katechismen}, 1915. F. Cohrs, \emph{Supplementa Melanchthoniana. Katechetische Schriften}, 1915. Catechisms antedating Luther's Catechisms are reprinted in Cohrs, op. cit., and Reu, Quellen. G. Buchwald and O. Albrecht, Bugenhagen's \emph{Katechismuspredigten} von 1525 und 1532, 1909. Instructions for Visitors, reprinted in Sehling, op cit., vol. i. *R. Vormbaum
\quellseite{105}
Die Schulordnungen des 16. Jahrhunderts, 1860. W. H. Frere and W. H. Kennedy, Visitation Articles and Injunctions of the Period of the Reformation, 1910. Koestlin-Kawerau, \emph{Martin Luther}, 5¹903. E. P. Enders, \emph{Luthers Briefwechsel}, 1884 ff. Smith-Jacobs, \emph{Luther’s Correspondence}, 2 vols., 1913–1918.

We must turn to the Church of the Reformation if we would find the most profound appreciation of the Church’s important teaching function and unexcelled zeal in its exercise. The Reformers examined and sifted the mass of material transmitted by the Medieval Church; they fixed a much more correct aim of instruction; they also improved the method of instruction; and above all, they understood evangelically and interpreted biblically the traditional material of instruction. As in nearly all other phases of the Reformation, so also here it was Luther himself who did most of the work and gained most of the achievements. His Small Catechism grew forth from intense pastoral occupation with the catechetical material; it is the result of thirteen years of strenuous labor.

When in the summer of 1516 Luther began to supply the pulpit of Simon Heinz at Wittenberg, he preached on the Ten Commandments, following the example of medieval preachers. Having finished this series in February 1517, he preached during lent 1517 a series of sermons on the Lord’s Prayer. In the same year he wrote a very brief exposition of the Ten Commandments in order to help the parishioners examine and prepare themselves for confession (cf. 1 Cor. 11:28). It is quite significant that then already he began to sift the material by calling attention to the fact that all customary catalogues of sins and vices are really superfluous because \emph{virtually contained in the Decalogue}. \emph{It is true}, he still enumerates the individual sins and virtues in the form of a \emph{tabula confessionis}, but he also points out that all sins flow from one source, selfishness, and all virtues from the love of God and fellowman. In 1518 Luther had this explanation
\quellseite{106}
of the Ten Commandments printed in German and Latin under the title \emph{Kurze Erklearung der 10 Gebote}. Its purpose as an aid in confession suggested the publication in the form of charts or placards rather than in bookform. In the same year he also revised for publication his sermons on the Decalogue of 1516/7 and published them under the title \emph{Decem praecepta Wittenbergensi praedicata populo}. Publication of the simple but edifying sermons on the Lord’s Prayer of 1517 followed in 1519; the title was \emph{Auslegung des Vaterunsers fuer die einfaeltigen Laien}. Two additional tracts on the Lord’s Prayer were printed in the same year, \emph{Eine kurze Form, das Paternoster zu verstehen und zu beten}, and the briefer \emph{Auslegung des Vaterunsers fuer sich und hinter sich}, both of them results of practical labors: in evening services Luther had explained the catechism parts \emph{pueris et rudibus}.

Luther’s explanation of the Decalogue received unstinted praise; it was declared that “he had taken the veil away from the face of Moses”; the explanation of the Lord’s Prayer was acclaimed even more highly. Beatus Rhenanus the humanist wrote to Zwingli that this treatise should be spread throughout Switzerland in every city, hamlet, and village, yes, it should find a place in every home. And the censor at Venice wrote: “Blessed are the hands which wrote this; blessed the eyes that read it; blessed the hearts that believe the book and so earnestly call upon God.”

In 1520 the Reformator\footnote{We use this term advisedly. The word “Reformer” fails utterly to do justice to Luther’s singular importance.} collected the results of his catechetical labors and added a third part to the Decalogue and the Lord’s Prayer. He published, in the form of a confessional mirror, a treatise for the common people, bearing the title Eine kurze Form der zehn Gebote; eine kurze Form des Glaubens; eine kurze Form des Vaterunsers. We find here
\quellseite{107}
Luther’s first explanation of the Creed, a truly classical model of vigorous diction and evangelical grasp. This work is an important landmark in Luther’s catechetical labors; it is the precursor of the Small Catechism. Three points deserve special attention: (1) The whole mass of auxiliary material handed down from the Middle Ages is swept aside; Luther says: “God so ordered it that the ordinary Christian who is unable to read the whole Bible, should be taught the Ten Commandments, the Creed, and the Lord’s Prayer, for these three parts comprehend all the essentials of Christian knowledge.” (2) Decalogue, Creed, and Lord’s Prayer had indeed been grouped in this order before Luther in a number of instances, but always more or less accidentally. Luther, however, consciously so arranged the material stating as his reason that these three parts so correlated reflect the way of salvation (cf. § 14, p. 111). (3) Instead of adhering to the traditional division of the Creed into twelve parts, Luther divided it into three parts with conscious reference to the saving deeds of the three persons of the Trinity. During the following decades the \emph{Short Form} was frequently used as basis and source for catechetical helps. In the form of the \emph{Booklet of Prayers} (1522) it was translated into English (1534?) and formed England’s first Evangelical catechism (Marshall’s Primer, ²1535).

Luther explains the first article in the \emph{Short Form} as follows:
\begin{quote}
“I renounce the evil spirit, all idolatry, all witchcraft and false belief. I put my trust in no man on earth, nor in myself, my strength, knowledge, property, piety, or anything I may possess. I put no trust in any creature in heaven or on earth. Deliberately I put all my trust in the one invisible and incomprehensible God, the creator of heaven and earth, the Lord of all the universe. Again, I have no dread of all the malice of the devil and his host; my God is higher than they all. I will believe in God though all men forsake or persecute me. I will still believe though I be poor, dull, unlearned, despised, or penniless. Sinner that I am, I will believe none the less. My faith shall soar above all things, above things that abide and things
\quellseite{108}
that fail, above sin and virtue, that it may rest, pure and true, in God alone, as the first commandment requires. I seek no sign from heaven to tempt God. And though he tarry, I will cling to him with constancy. I will fix for him no limit or time, I will bind him to no measure or mode, but in faith free and true I will commit all to his divine will. He is almighty: whatever I lack he can supply. He is the creator of heaven and earth and Lord of all things: who can harm or hurt me? Yea, all things must work together for good if he is on my side whom all things serve and obey. He is my God, omniscient and mighty to help in every time of trouble. He is my Father, with loving heart supplies my every need. Nothing doubting, I will put my trust in him and so remain his child, servant and heir forever. And as I trust, so shall it be.”
\end{quote}

Also after 1520, catechetical labors occupied Luther’s attention time and again. Perhaps the disturbances raised by Carlstadt at Wittenberg moved him to preach on catechetical material more frequently than before. In 1522 he preached on the trilogy (Decalogue, Creed, Lord’s Prayer) and, in order to suppress mariolatry, also on the Ave Maria. In the same year he also learned to know the German edition of the \emph{Questions for Children}, the Catechism of the Moravian Brethren (see p. 79); possibly this work exerted some influence upon Luther’s catechetical activity. The Short Form of 1520 was published in the summer of 1522 as the \emph{Booklet of Prayers}; in this form it went through many editions and met the requirements of school and home much more effectively than the numerous editions and revisions of the \emph{Questions for Children} just mentioned. In 1523 Luther again preached sermons on the trilogy; as a new feature we find the \emph{Five Questions concerning the Lord’s Supper} taken from a sermon by Luther on Holy Communion. These questions were propounded to prospective communicants who since Easter 1523 were examined as to their knowledge of the nature and meaning of the Lord’s Supper. Owing to Luther’s efforts Wittenberg received in 1521 a regular catechist for its youth in the person of Agricola.

\quellseite{109}
In the first part of 1524 the Reformer issued his proclamation “\emph{To the Councilors of all Cities in the German Empire on the Establishment and Maintenance of Christian Schools}” arousing all Germany with flaming words to attend to the duty of Christian training. “When the Turk threatens,” he says, “or when you are in danger of war or perils of water, you know without exception what to do; you decide how much money a city shall spend on muskets, roads, and dams. Much money has been squandered on indulgences, masses, pilgrimages, and the like. Oh, that but a part of this money might be spent upon the training of the children! If you give one guilder for the war against the Turk, you would not do too much by giving a hundred of them for the proper training of one single boy.” While chiefly interested, at that time, in the maintenance of the Latin schools, Luther hints at the building up of a Christian public school system, at least in the cities, when voicing the demand that boys should be given a minimum of two hours of daily instruction and girls one hour.

More and more Luther was impressed with the necessity of publishing in bookform a brief exposition of the chief parts, in order to facilitate Christian training of the young and immature. The Booklet of Prayers though used occasionally for this purpose, had not been written specifically with this aim. In letters to his friend Nicolaus Hausmann of Zwickau who had strongly emphasized the duty of Christian training in two official opinions (1523 and 1525), Luther repeatedly discussed this subject. Himself occupied too much with other matters, he commissioned, in February 1525, Justus Jonas and Agricola with the preparation of a \emph{catechismus puerorum}, this being the first time that the word catechism was applied to a book. When this project came to nought, Luther decided to write a catechism himself; to what extent he was acquainted with catechetical efforts put forth by other
\quellseite{110}
Evangelicals, we have no means of knowing. Before he was able to carry out his plan, however, the \emph{Booklet for Laymen and Children} appeared, in the fall of 1525. If this booklet was composed by Bugenhagen which is quite probable, Luther must have been cognizant of its composition; in fact, it presents in part the explanations of the Reformator and in other respects approximates his thoughts. Three factors render the booklet particularly noteworthy: for the first time Baptism and Communion appear as catechetical parts beside the trilogy; the text of the Catechism is largely couched in the wording with which we are familiar; and several prayers are added which we find later in Luther’s Small Catechism.

In the year 1526 Luther brought forth the “\emph{German Mass}”, another important landmark in his catechetical labors. Far and wide in evangelical lands it served as a model for divine worship. Everywhere it fired the consciences in regard to religious instruction with words like these: “On in the name of God! What we need most in our German worship, is a plain, simple, clear and succinct catechism (i.e., oral instruction).” As subject matter he had in mind the trilogy: “I do not know how to make such instruction better or simpler; it is best to leave it as it was established at the beginning of the Church and has remained ever since, and to teach these three parts: the Ten Commandments, the Creed, and the Lord’s Prayer.” The father shall diligently teach these parts to his children and servants in the \emph{home}; but such instruction is also to be part and parcel of the divine \emph{service}: “These lessons shall be read from the pulpit at stated seasons, or daily, according to circumstances, and in the homes they shall be read or recited to the children in the morning and in the evening.” In addition he requires that on Mondays and Tuesdays at the matin service a German lesson (i.e., a catechetical \emph{sermon}) shall be given concerning Decalogue, Creed, Lord’s Prayer, Baptism, and Communion,
\quellseite{111}
‘in order that these two days serve to maintain and to explain the catechism.” Furthermore it is noteworthy that Luther seriously insists that the children really learn to \emph{understand the catechetical parts}, and that he shows by practical examples how one should teach a child. He revives the important pedagogical principle of Chrysostom and Augustine: “Christ when he came to train men, had to become man; if we are to train children, we must become children with them.” Herewith a new educational aim was established. He also demands that parents teach their children a number of Bible passages and unfold to them their meaning. Finally, Luther suggests the preaching of catechetical sermons on Baptism and Communion; he himself preached a sermon on Communion annually from 1523 to 1528, and on Baptism from 1526 to 1528. Thus it is seen that the addition of the Fourth and Fifth Chief Parts of the Catechism is heralded.

The German Mass stimulated a considerable number of Evangelicals to write catechetical text books for the youth. The total number of catechisms issued between 1522 and 1529 is about thirty, and some of them were published in many editions. Besides the Booklet of Prayers and the Booklet for the Laity and Children, the most important are the books by Capito, Lachmann or Graeter, Agricola, Sam, Althamer, and Brenz. Althamer’s book was the first to bear the name “catechism” in its title; and Brenz’s is worthy of notice because it contained two graded sets of questions prepared respectively for younger and older children, a feature which later became important, and because it begins with Baptism and by means of appropriate questions unites the several parts of the catechism into a system. How Luther evaluated these attempts if he was acquainted with them at all, we do not know. At any rate, he did not abandon his plan of composing a catechism of his own.

\quellseite{112}
In 1527 Luther preached once more on the Ten Commandments while Melanchthon in collaboration with him wrote the “\emph{Instructions for Visitors}” which appeared in print the following year. The institutions for the training of the common people and the young which so far had been introduced locally, were now extended to the whole of Saxon territory, with limitations even to the villages. It is enjoined that servants and young people attend church on Sunday afternoons in order to study Decalogue, Creed, and Lord’s Prayer, subsequently also Baptism and Communion, and that every catechetical sermon end with a recitation of the trilogy. In cities and market towns where day schools existed, Saturday or Wednesday was to be devoted exclusively to religious instruction. Here also some of the easier Psalms were to be committed to memory, and the catechism was to be explained with the aid of Bible stories. In the villages, according to the visitation records, the verger (Kuester) was directed to instruct the children in the Decalogue, the Creed, and the Lord’s Prayer, and also in the more important German hymns. For the time being this was the limit of the attainable; erection of regular village schools was not to be thought of. Gradually knowledge of the trilogy had come to be regarded as a condition for the admission to the Holy Supper. Thereby the way was prepared for a new educational goal, admission to Holy Communion.

In conformity with the injunction of the \emph{Instructions for Visitors} Luther preached three series of catechetical sermons in 1528. Following a precedent set by Bugenhagen, he delivered them during the Ember weeks of May, September, and December; the last series he preached after personally taking part in the visitations. The results of the visitation were such that Luther felt the imperative need of publishing a \emph{catechism pro pueris et rudibus}, and he decided to do it on basis of the three series of catechetical
\quellseite{113}
sermons of 1528. In a letter bearing the date of January 15, 1529, Luther writes that he was engaged in the preparation of a catechism, and on the twentieth of the same month Roerer, the Wittenberg proof-reader, expressed the hope that the work would soon be on the market. Undoubtedly, therefore, about the close of 1528 or the beginning of 1529 Luther began the composition of this \emph{catechismus praedicatus pro rudibus et simplicibus} (= for the crude heathen, the uneducated masses), that is, the work that was later called the Large Catechism. But during its preparation, Luther evidently concluded that this work would be too elaborate for the “crude heathen” and quickly resolved to write also a brief catchism \emph{pro pueris et familia}. This work, or rather the first three parts were printed in the first week of 1529\footnote{The time of publication is fixed definitely by the Schoenewalde visitation records. In the constitution which the visitors, among them Luther, gave this city, the stipulation is found that on Sunday afternoons “the Ten Commandments, the Creed, and the Lord’s Prayer are to be recited for the people and, after that, to be explained in a very simple manner as shown by the table published.” Now this constitution cannot have been promulgated earlier than January 7 nor later than January 9. On January 12, Roerer sent these tabulae catechismi to Spalatin at Altenburg; according to a letter of January 20, Roerer had charts hanging on the walls complectentes brevissime simul et crasse catechismum Lutheri \emph{pro pueris et familia}. February 12, he announced that the supply of tables was exhausted and that one could not buy one even for a gold guilder (originally they sold at two or three pennies).} and published in the form of tables. In choosing this form of publication Luther followed the practise of the later Middle Ages (cf. also p. 75); printed on tables or placards and fixed upon the wall in the home or the school, they might be seen and read by all at all times. It is probable that now already the superscription of the charts was addressed “to the fathers that they may instruct their children and servants.” The tables issued at this time contained no more than the three parts which he had repeatedly designated as genuine \emph{Kinderlehre}, viz., Decalogue, Creed, Lord’s Prayer, together
\quellseite{114}
with an explanation. Not until March 16, probably in consequence of Luther’s illness and the pressure of other urgent work, did the \emph{tabulae confessionis} together with the \emph{Litania Germanica}, and the \emph{tabulae de sacramentis baptismatis et sanguinis Christi} appear. They were ready for transmission just in time for the Easter confession and communion although this factor did not exclusively prompt their publication at this time. Luther had observed since compulsory confession had been abrogated how frequently evangelical freedom was used as a cloak to cover one’s contempt of the Sacrament; to promote a correct understanding of the Sacrament and to arouse a heart-felt desire for it, Luther added these parts to the first three parts (cf. the Preface to the Small Catechism). Probably at the same time two additional tables were published, the one containing the morning and evening prayers, the other the \emph{Benedictine and the Gratias}; a Low German table, printed at Wittenberg in 1529, containing the morning and evening prayers has come down to our times. When the last part, the table of domestic duties (\emph{Haustafel}) was published, is unknown. The first to unite all these tables as far as he knew them, in book form, apparently was Bugenhagen; occupied at the time with the regulation of church and school affairs at Hamburg he issued the five Chief Parts and the Table Prayers in Low-German; a copy of this book, with a Hamburg imprint, is extant.

So far as Luther is concerned, he issued in book form first what was later called “The Large Catechism.” The completion of this book had been somewhat delayed, but it was ready for transmission on April 23, 1529. A companion volume for the German Mass of 1526, its title was simply \emph{Deudsch Catechismus}. Mart. Luther. There was as yet no occasion to designate this book as the Large Catechism for the reason that at this time no book with the title “Small Catechism” existed. Its component parts still circulated in
\quellseite{115}
the form of tables; and the above mentioned Hamburg print was designated \emph{Eyn Catechismus effte underricht}, not “Small Catechism.” After Luther had published the substance of the tables in bookform and called it Small Catechism, the new title “Large Catechism” gradually replaced the original title “Deudsch Catechismus.”

This “\emph{German Catechism}” contained (1) the texts of the Ten Commandments, the Creed, and the Lord’s Prayer, “the most necessary parts which every Christian should first learn to repeat word for word”; (2) the words of institution of Baptism and Communion, introduced by the remark: “When these three parts are apprehended, we ought to know how to speak concerning our Sacraments which Christ himself instituted”; (3) the detailed exposition of these texts. A brief preface served to remind especially the fathers of their duty. In the same year 1529 a second edition appeared, augmented by a confessional exhortation (cf. Mueller, \emph{Symbol. Buecher}, pp. 773 ff.) and an addition to the exposition of the Lord’s Prayer (cf. Mueller, \emph{op. cit.}, p. 463). In 1530 a third edition was published, containing a second, rather lengthy preface (cf. H. F. Jacobs, \emph{Book of Concord} i, pp. 383 ff.) in which Luther also justifies catechism study in the well known words, “I am also a doctor and preacher etc.” (cf. Jacobs, \emph{op. cit.}, p. 384). From that time its compass has remained the same.

Having grown forth from Luther’s catechetical sermons, the “\emph{German Catechism}” offered excellent models for such sermons, and often was read in the place of original sermons. Even today it is the best commentary on Luther’s Small Catechism.

After issuing the Large or \emph{German Catechism}, Luther arranged for the issue in book form of the educational material hitherto published in the form of tables. The issue was completed on May 16, 1529. Already on June 13, an “enlarged and revised edition” was in circulation (the second [?] or third Wittenberg edition). The original Wittenberg edition is no longer extant, but we possess one Marburg and two Erfurt reprints (possibly these texts are reprints of the second Wittenberg edition which is likewise extinct). According
\quellseite{116}
to these reprints, the booklet had the title \emph{Der kleine Catechismus fur die gemeine Pfarherr vnd Prediger. Mart. Luther. Wittemberg.} As is seen, Luther dedicated the book to the preachers although the superscriptions at the head of the several parts\footnote{“In the plain form in which they are (resp. it is) to be taught by the head of a family.” The precursors of the Small Catechism, the tables, had been addressed primarily to the heads of families.} were not eliminated. No doubt he did this, in part at least, because of the ignorance of many preachers and because they were in need of a definite form upon which to base the instruction of the common people and especially the young.

The Original Edition (May 16, 1529) evidently had the following contents: (1) Preface to pastors and preachers; (2) text and exposition of the five Chief Parts (the three reprints omit the third question pertaining to the Lord’s Supper, possibly on account of a misprint in the original Wittenberg edition); (3) Morning and Evening prayer, and the Blessing and Thanksgiving at Table; (4) a selection of appropriate “passages for various holy orders and conditions of men”; (5) the Marriage Booklet. Confession received no attention in this edition; and the Lord’s Prayer lacked the introduction and its explanation.

The Third (or second?) Wittenberg Edition (June 13, 1529) bears the word “\emph{Enchiridion}” on its title page and is designated as “enlarged and revised.” The following parts were added: (1) The third question pertaining to the Lord’s Supper; (2) the Baptismal Booklet in the form of 1526; (3) “A brief form of Confession before the Pastor, for the Simple” (not identical with the statement concerning confession which was later incorporated); (4) the German Litany with music, and three concluding collects. This is the earliest Wittenberg edition which is extant though poorly preserved. By the addition of the Baptismal Booklet, the Form of Confession, and the Litany, the little book had grown to the size of a pastoral manual. For this reason perhaps the word \emph{Enchiridion} was added on the title page. But Luther had in mind not only the pastors, but also the laity; he was in hopes that these additions would bring about an intelligent appreciation of, and active participation in the various acts outlined. In fact, during the sixteenth century, the Small Catechism largely served the purpose of a “church book”; other liturgical
\quellseite{117}
parts were added from time to time, as Psalm 111 which was sung at the celebration of the Holy Supper, the \emph{Tedeum}, the \emph{Magnificat}, the Prayer against the Turk (“Lord, keep us steadfast in thy word”). We observe here the faint beginning of liturgical training as a part of Christian instruction (cf § 11, p. 66). Finally, this edition was also embellished with biblical pictures corresponding to the parts of the Catechism and offering welcome illustrative material; whether this applies also to the first and second Wittenberg editions, we do not know.

The Edition of 1531 brought the following changes: (1) The introduction to the Lord’s Prayer with its explanation was added; (2) the paragraph on Confession which was found in the third edition immediately after the Baptismal Booklet, was replaced by the part entitled: “How the Simple should be Taught to Confess” (i.e. the three questions (1) What is Confession? (2) What Sins should we Confess? (3) Which are these?, and the section “Please show me a short way to confess”). Luther inserted this part between the Fourth and Fifth Chief Parts.

In all essentials the Enchiridion retained the form of 1531. The few changes which were made and in part were dropped again, are rather insignificant. They are chiefly attempts to harmonize the Bible passages in the Catechism with the phraseology of the German Bible; Luther himself had simply adopted the text form current at Wittenberg. Little consistency, however, may be observed in these efforts. Thus we find missbrauchen in 1536, but in 1540 the original unnuetzlich fuehren; in 1540 the promise is added to the Fourth Commandment; the Table of Domestic Duties is enlarged in 1540 and again in 1542; no Wittenberg print of Luther’s time has geluesten in the Ninth and Tenth Commandments or the threat in the Second Commandment. Most of these changes (also the addition of the Introduction to the Ten Commandments: “I am the Lord thy God”) can be traced back to the catechism text of the very important \emph{Nuremberg Sermons for Children} of 1533. To what extent Luther himself authorized these changes, or whether the typesetter and the proof-reader are responsible for them, is impossible to determine. At any rate, the Eisenach Conference was ill-advised when it adopted the text of 1542 as the basis for the uniform text of the Small Catechism.

From the foregoing statements it is readily inferred that the material appearing now as Fifth, now as Sixth Chief Part, under the title “\emph{Of the Office of the Keys and Confession}” belongs only in
\quellseite{118}
part to the catechism issued by Luther. The paragraph on Confession is indeed Luther’s product, but that on the Office of the Keys is traced back to two other sources. The Nuremberg Sermons for Children (composed by Osiander and Schlenpner, not by Brenz or even Veit Dietrich; appended to the Nuremberg Church Order of 1533) regularly conclude with the words of Luther’s Catechism. Now this series contains between Baptism and Holy Communion a sermon on the Office of the Keys which is summarized in the words John 20:22–23 and the explanation: “I believe when the called ministers etc.” Thus the appearance of these words as part of Luther’s Catechism is easily accounted for. The question “What is the Office of the Keys” with its answer is found in the exact form in which we know it today, for the first time in an edition of Luther’s Catechism for Saxony-Altenburg of 1584 (1582?); but the substance is found much earlier, even before 1550, in original catechetical compositions; owing to Knipstro’s endeavors it forms a regular part of the edition of Luther’s Catechism for Pomerania (1554).

The so-called “\emph{Questions and Answers for Prospective Communicants}” so frequently appended to Luther’s Catechism and even ascribed to him (since 1549), were not composed by Luther, but possibly by his friend, Dr. John Lange of Erfurt.

As soon as Luther’s Catechism appeared, they were translated into Latin. John Lonic­er of Marburg, and a little later Obsopoeus of Ansbach translated the Large Catechism. An unknown author translated the Small Catechism for a Latin edition of the above mentioned Booklet of Prayers, taking the liberty of changing the question and answer form into the thetic form. Probably dissatisfied with this translation, Luther commissioned a certain John Sauermann with preparing a Latin translation and assisted him in this work. Its title \emph{Parvus catechismus pro pueris in schola} indicates that it was intended for school use; wherever in the German original the father of the family was mentioned, this edition substituted the \emph{paedagogus}. Since 1529 this Latin version has gone through many editions, and even today deserves to be consulted.

The appearance of Luther’s Catechisms, especially the Small Catechism to whose significance we shall devote a special chapter, is sufficient to stamp the year 1529 as the most momentous in the history of religious instruction. But so far we have not even touched upon another very important
\quellseite{119}
work which was issued also by Luther, in the same year, a work which may well be called the first Biblical History for the Christian Home. We have in mind Luther’s \emph{Passionale}. The importance of history, especially sacred history, for the education of the young had been stressed repeatedly by Luther and Melanchthon before this time; Otto Braunfels of Strassburg had made an effort to introduce instruction in Bible History in his Latin school and had collected in his \emph{Heldenbuechlein} a number of exemplary and godless characters from the Bible; and Luther himself—perhaps unconsciously following the cue given by Melanchthon—had inserted biblical pictures in his Small Catechism and advised the teacher to adduce many biblical examples showing how God punishes the wicked and rewards the godly (cf. the Preface to the Small Catechism). But when he issued a new edition of his Booklet of Prayers in 1529, he added to it the \emph{Passionale}. This is a collection of 49 pictures with explanatory texts; 11 of them are from the Old Testament, the remainder from the New Testament; no less than 15 illustrate phases of the Passion of our Lord, hence the name of the book. The work was written for the Christian home with special reference to the instruction of children. In the preface the Reformator says explicitly that he collected the pictures and supplied them with biblical texts “chiefly for the sake of the children and the simple folks who will remember the sacred stories more readily when you use pictures and illustrations in teaching.” It is true that Luther followed medieval precedent in the selection and publication of this material; but the finished product was something new and reformatory, something totally different from any medieval picture book, for (1) Luther offered only biblical pictures expunging all legendary matter; (2) he augmented the number of pictures and arranged them according to the order of sacred history; (3) he added carefully selected portions of
\quellseite{120}
Biblical History. Luther blazed the way—though perhaps not quite as successfully as in catechetical instruction—and unquestionably deserves to be honored as “Father of instruction in Bible History in the Evangelical Church.” Except for its name this significant booklet was completely forgotten until the present writer had it reprinted in his \emph{Quellen zur Geschichte des kirchlichen Unterrichts im evangelischen Deutschland zwischen 1530 und 1600}. In the same work about thirty different editions dating from 1529 to 1600 are noted, giving eloquent evidence of the wide circulation of this book.

\quellseite{121}
\section[§ 14. LUTHER’S SMALL CATECHISM]{§ 14. THE SIGNIFICANCE OF LUTHER’S SMALL CATECHISM}
G. v. Zezschwitz ii, 1, ch. 37–40. *G. v. Zezschwitz, \emph{Luthers Kleiner Katechismus. Seine Bedeutung u. s. Urgestalt}, 1881. J. Gottschick, \emph{Luther als Katechet}, 1883. K. Knoke, \emph{Ueber Katechismusunterricht}, 1886. G. v. Rohden, \emph{Ein Wort z. Katechismusfrage}, ²1890. G. v. Rohden, \emph{Ueber christocentrische Behandlung des luth. Katechismus}, 1891. Th. Kaftan, \emph{Auslegung des luth. Katechismus}, 1892, ⁶1913. Th. Hardenland, \emph{Die katechetische Behandlung des Kleinen Katechismus Luthers im Geist s. Verfassers}, 1899. F. Cohrs, \emph{Katechismusversuche vor Luthers Enchiridion}, 1900 ff. Koestlin-Kawerau, \emph{Martin Luther. Sein Leben u. s. Schriften}, 1903. H. Huebner, \emph{Was d. Kleine Katechismus fuer ein grosser Schatz ist}, 1904. M. Reu, \emph{Die Eigenart des Katechismus Luthers u. s. katechetische Behandlung} in Kirchliche Zeitschrift, 1906. J. Gillhoff, \emph{Zur Sprache u. Geschichte des Kleinen Katechismus}, 1909. *F. Rendorff, \emph{Das Problem der Konfirmation u. der Religionsunterricht in der Volksschule}, 1910. O. Albrecht, \emph{Luthers Katechismen}, Weimar Edition of Luther’s Works, vol. xxx, pp. 635 ff., 1910. *M. Reu, \emph{Die grosse Bedeutung des Kleinen Katechismus Luthers}, 1913.

In issuing the Small Catechism Luther \emph{gave the Church} a book of singular, superlative importance. Its great significance is found first in the fact that it represents \emph{the crowning consummation of the educational endeavors put forth by the Church} during the course of fifteen centuries. Although Luther did not include in his Catechism all those elements of instruction which had been employed during this period of development, he gathered all the necessary and essential material and issued it in handy and convenient form. Quite correctly he eliminated the mass of auxiliary material accumulated during the Middle Ages, retaining, however, the Decalogue which had not until lately been in universal usage. Luther had a conservative and historical turn of mind; and he was well aware of the intimate relation of his catechetical work to that of the past (cf. his preface to the Short Form). This connection shows itself not only in the
\quellseite{122}
adoption of the doctrinal parts (Decalogue, Creed, Lord’s Prayer, words of institution of the two Sacraments, in part also the prayers), but also in the wording of the catechetical texts, and, in the Third Chief Part, even in the explanation.

Thus Luther retained the traditional enumeration of the Commandments, omitting the prohibition of images and keeping separate the Ninth and Tenth Commandments; this arrangement had won favor since Augustine’s times and at the close of the Middle Ages was the only one in use. He retained the abbreviated form of the first Two Commandments, “Thou shalt have no other gods, Thou shalt not take the name of the Lord thy God in vain.” He did not, as is frequently stated, cast the Third Commandment in a new form when he wrote “Thou shalt sanctify the holy day”; in this instance the reformatory element is found in the explanation. In the First Article he clung to the form \emph{Vater allmechtiger} (correct Middle High German where the attributive adjective often follows the noun), although other catechisms of his day adopted the form \emph{Allmaechtiger Vater} which grammatically is more correct. He adhered to the words “resurrection of the flesh” though stating in the Large Catechism that “resurrection of the body” might be preferable. With all countrymen of his he continued to pray \emph{Vater Unser} (Old High German grammar places the pronoun after the noun in the vocative case) while in his German Bible he transposed the two words in harmony with the grammatical rule which gradually became recognized. He retained the broader form of the Seventh Petition “And deliver us from evil” and explains it in this sense although in the Large Catechism he states that the original speaks of the evil one, the devil; when Bugenhagen introduced this form at Wittenberg in 1533, he raised no objection. In the Fifth Petition he prays with the Christian people of his time “\emph{Verlasse uns unsere Schuld}” (\emph{Schulde} is plural) while in his German Bible he wrote “\emph{Vergib uns unsere Schuld},” the form which later gained greater currency. Finally, he omits the Doxology at the end of the Lord’s Prayer (it was missing in the old manuscripts and in the Vulgate, hence also in church forms), retaining and explaining only the Amen. Luther knew himself called to something greater than a philological reformation of linguistic forms; with pedagogical insight and historical piety he preserved whatever might be safely retained.

The traditional form, however, was for Luther no inviolate
\quellseite{123}
shackle. He changed the “un-German” \emph{Gemeinschaft der Heiligen} in the Third Article into \emph{Gemeinde der Heiligen}, and he added the Conclusion of the Ten Commandments which was not found at all in the traditional forms of the Decalogue, removing it from its place in the Old Testament and placing it at the end. In both instances he was actuated by pedagogic considerations which to him appeared of greater importance than devotion to the customary form.

The second and most significant feature of the Small Catechism is the \emph{deeply evangelical interpretation of its constituent parts which issues from the article of justification}. This is made evident in the \emph{First Chief Part} by three important points: (a) by the freedom with which he treats the Old Testament text. For in omitting the prohibition of images Luther did not merely follow medieval precedent, but actually exemplified his evangelical attitude toward the Old Testament law. According to Luther, the Christian needs no longer any prohibition of image worship, and whatever permanent value this inhibition possesses, is found in the First Commandment correctly understood. Attempts have been made, are being made even in the present, to correct Luther in this matter; but they all betray more or less distinctly a legalistic trend, in harmony with the Old Testament rather than the New. Their authors and advocates ought to study Luther’s writing \emph{Against the Heavenly Prophets, of Images and the Sacrament}. Here he demonstrates with all possible clearness that the Old Testament Decalogue contains also ceremonial and civil regulations the observance of which was restricted to Israel. Among such temporary laws he classifies the prohibition of images and the Sabbath law; and he states that the Decalogue is obligatory only in so far as it corresponds to the natural law, and that therefore the New Testament determines what features are of permanent force.

Luther’s thesis is: “All the laws which Moses added to the natural
\quellseite{124}
law, as those concerning images and the Sabbath, are now abrogated and no longer in force because they are not part of natural law and were intended for Israel alone. An emperor or king may give his people special laws and ordinances, like the Saxon mirror of Saxony; but natural law is in force everywhere, as honoring one’s parents, abstaining from murder and adultery, serving God, etc. France pays no attention to the Saxon mirror; yet it agrees with Saxony in the observance of natural law. Therefore let Moses remain the Saxon mirror of the Jews and let no one trouble us gentiles with it. Why then do we keep and teach the Ten Commandments? Answer: Because natural law is nowhere so well arranged and composed as in Moses.”

In the Third Commandment he ignores the concept “day” entirely. He does not consider Sunday as a substitute for the Sabbath: the Sabbath is abrogated (Gal. 4:10 f.; Col. 2:16 f.). The establishment of the holy days of the Church would merely serve to sanctify all the days of the week.

(b) In the second place Luther’s evangelical interpretation of the Decalogue shows itself in the attitude toward the letter of the law. Luther follows the precedent set by Jesus in the Sermon on the Mount: he does not stop at the individual, overt act named in the commandment; he rather takes the sin mentioned in the commandment as the coarsest outgrowth of the wrong forbidden: he traces all sinful thoughts, words, and deeds back to sin dwelling in man’s heart; he demands that the Christian battle against sin and strive after perfection. He requires of him a new attitude: external deeds may be done even by natural man, but true Christian morality grows forth from the proper disposition of a godly heart, from a holy attitude which influences every phase of life.

(c) Those two factors would suffice to lift Luther’s explanation above most of the medieval explanations of the Decalogue; but its evangelical character is most plainly evidenced, thirdly, by the fact that he points out the “fear and love of God” as the one and only root of Christian morality. By “fear” he means reverence, veneration (Ehrfurcht), or filial
\quellseite{125}
awe. That Luther traces the words “we should fear and love God” through all the Commandments and thus welds them into a unity, this is an achievement of truly reformatory significance, an absolutely unique accomplishment. Thereby an end is made of the “bargain features” of Romanism; no longer work is added to work in the hope that the sum of them all will suffice to earn God’s favor. The Christian life is no more regarded as a series of separate works, but as the organic development of fear and love which God himself kindled in the heart. Instead of bewildering the child, or common man, or mature Christian with an almost countless number of virtues, Luther shows him the one and indispensable requirement; fear and love embraces, and produces everything else. And by representing as perfectly worthless all other works if detached from this source, he emphasizes most effectively the important evangelical principle that works are not acceptable to God unless the person has first become acceptable.

The fairest pearl in the whole Catechism is the explanation of the Second Chief Part, and in the whole realm of catechetical literature we know of nothing that equals or even approximates its chaste beauty. In two ways Luther’s evangelical understanding of Holy Writ has found clear expression. (a) His first merit is that he groups all the material of the Creed which up to that time had been treated as a mass of unrelated detail, around the three great saving deeds of the Triune God — creation, redemption, and sanctification. As early as 1520 Luther indeed divided the Creed into three parts, and most of the catechisms brought out since that time had followed this method, but neither Luther nor his followers had advanced beyond a mere tripartite division; during the process of explanation they would dissolve the three articles again and again into a number of independent particles and ascribe to each one of them a specific redemptive
\quellseite{126}
significance. In Brenz’s Catechism, e.g., we read: “What benefit do you derive from this article: ‘Conceived by the Holy Ghost, born of the Virgin Mary’? Answer: I receive this benefit that my conception in sin is not accounted as sin in the judgment of God, and that my sinful birth is sanctified through Jesus Christ. — What benefit do you derive from this article: ‘Suffered under Pontius Pilate’? Answer: That his life, in God’s judgment, is a satisfaction for my sin, and that all my sufferings are blessed and sanctified through his.” In the Small Catechism, however, Luther connects all the individual statements, subordinates them to the main thought, conceives them as the organic members of a unit, and, in truly classical fashion, sets forth the meaning of this unit for the Christian life. In the First Article he selects the idea of creation, places it in the center, and makes it the controlling factor. It is true, he also confesses God as “Father Almighty” (“and all this purely out of fatherly, divine goodness and mercy”; as a matter of fact, the idea of God’s fatherhood is present also in the foregoing statements), but he carefully refrains from making this the governing aspect. How exactly this mode of treatment harmonizes with the New Testament conception of faith, may be seen especially plainly in the Second Article. Viewed separately, the events of Christ’s life here mentioned, are without saving significance, but as joint factors of the life of Jesus they become supremely important. Of what benefit were the conception of Jesus by the Holy Spirit and his birth of the Virgin Mary if they had not been followed by his passion, death, and resurrection? None of these facts isolated from the rest, has brought about our salvation; all of them together have done so. Hence Luther assigns the central position to the word “I believe in our Lord,” subordinates to this central thought all the rich material of the Article, and uses it to explain who our Lord is, and by what
\quellseite{127}
means he became and will ever more be our Lord. Luther’s method of explaining the Third Article is similar. All previous catechists, only a few church fathers excepted, divided this article into five coordinate particles. Luther, however, discovered and demonstrated the underlying unity; without straining a single particle he obtained the outlines of the genuinely evangelical way of salvation, i.e., the fundamentals of sanctification or of the way in which souls, in time and eternity, are brought to Christ, their Lord and Redeemer. For it is through the church, through forgiveness of sins, through resurrection, and through reception into life eternal that the one work of the Holy Spirit is wrought who in the past has led us to Christ, in the present daily brings us to him again, and in the future will bring us to him forever. Thus the explanation is at the same time a classical expression of the basic truth that salvation from beginning to end depends not on human agencies but on Christ and his Spirit.

(b) The second feature which reveals Luther’s evangelical treatment of the Creed, is the vital relation which he establishes between the facts confessed and the individual confessor. He is not satisfied with a mere \emph{fides historica} concerning facts of salvation belonging to the past. In confessing the First Article, the Christian, according to Luther, does not merely signify his assent to the fact that God once upon a time created the world, but he confesses that he himself owes life and every blessing to no one but God, his creator and protector, that his whole life depends upon God, and that he is in duty bound to praise and to obey him. He who confesses the Second Article according to Luther, does not simply declare that Christ performed a number of saving acts and so became the Lord of Christendom, but he connects all this with his own life and confesses that Christ has redeemed him, a lost and condemned creature, purchased and won him and thereby has become his Lord and that he
\quellseite{128}
would serve him in everlasting righteousness, innocence, and blessedness. The same personal note rings through the Third Article. But this is the very characteristic of saving faith truly evangelical; the evangelical Christian correlates the redemptive facts of the past with the present and with his own life, he relies exclusively upon grace divine, and he dedicates himself to God as a living sacrifice, holy and acceptable.

Luther’s exposition of the \emph{Lord’s Prayer} is less original. Owing to his dependence upon early Christian models (esp. Tertullian and Cyprian) he does not always penetrate to the deepest meaning of the words (e.g., in the Second Petition). Also here, however, everything is genuinely evangelical. This may be seen from the arrangement as a whole which takes up the First and Second Chief Parts and makes that the subject for prayer which has there been disclosed as the saving and commanding will of God; furthermore from the explanation of the Fifth Petition and its emphasis upon divine grace; and especially from the most precious interpretation of the name “Father” and of our sonship in the Introduction. The explanation of the Introduction was not added until 1531, but it is the most beautiful part of the exposition of the Lord’s prayer, in the sunburst of the Catechism one of the most resplendent gems.

The very addition of the exposition of the Fourth and Fifth Chief Parts is a loud protest against the Roman \emph{opus operatum} and every kind of mechanical training of the masses for outward observance. In particular, however, Luther shows his evangelical understanding in the \emph{Fourth Chief Part} by placing the proper emphasis upon the subjective element, faith, which conditions the saving appropriation of the baptismal blessing (“to all who believe what the words and promises of God declare,” “and faith which trusts this word of God in the water”); and he does this without in
\quellseite{129}
the least infringing upon the objective efficacy of baptism as a saving act of God. Moreover, he sets forth the initial sacrament as a factor influencing the whole Christian life and as the source of numerous ethical motives, so that the Christian life virtually passes under the aspect of a daily \emph{regressus ad baptismum} (“it signifies that the old Adam in us should, by daily sorrow and repentance, be drowned and die”).

Also the \emph{institute of \emph{Confession}} which was completely obscured by medieval additions of every description, was restored to its pristine purity and simplicity. To Luther, its center of gravity is absolution following upon sincere confession of sins before the Lord. Abolishing compulsory enumeration of all sins, he gives the advice to confess to the pastor such sins as press heavily upon the heart because by such private treatment the individual will be more effectively counseled and assured of forgiveness. The Decalogue evangelically interpreted is to be used as the mirror in which the soul may see its sin. Resolutely he stripped confession of its sacramental character; from the position of chief sacrament he reduced it to a mere episode between baptism and communion. It bears upon the one as well as upon the other, and has a deep pedagogical significance; but after all, in contrast to the preaching of the Word and to the administration of the Sacraments, it is no more than an act instituted by the Church.

In the \emph{Fifth Chief Part} the evangelical understanding of the Reformator is seen in the correct definition of the nature of the Sacrament of the Altar as grounded in the words of institution; in the stress laid upon repentant faith as the prerequisite of a blessed reception; further, in the reiteration of the phrase “for you”; most conspicuously, however, in the statement, “Where there is forgiveness of sins, there is also life and salvation.” This sentence is the most
\quellseite{130}
forcible evidence of evangelical understanding; it places the reformatory conception of justification by faith in glaring contrast to Rome’s view of justification. Forgiveness of sins is not the initial stage on the way to God’s favor from which one must work his own way to justification and God’s favor and full salvation; but he who has received forgiveness of sins has everything else and is a saint in the sight of God.

Truly evangelical ideas, finally, prompted Luther to append, resp. to expurgate, the \emph{remaining parts of the Catechism} which today are neglected altogether too much. The sanctification of daily labor and daily food through God’s word and prayer is necessary to a genuine Christian life, and the more completely the life of prayer is regulated by system, the better and healthier it is. Regularity, habituation in matters of Christianity is of incalculable value; although it is true that nothing is more detrimental than the shriveling of the inner life to a routine of mere lifeless habit. The prayers received by Luther into the Catechism or composed by him for the purpose are free from unbiblical medieval additions and genuinely evangelical; when he advocates kneeling, folding of hands, and the sign of the cross, he exhibits a noteworthy appreciation of symbols in family or private devotions. — A thought diametrically opposed to medievalism is expressed in the incorporation of the Table of Domestic Duties, namely the biblical and decidedly reformatory idea that a man must express and prove his Christianity in the divinely ordered calling and profession, not in monastic world-flight. There are no two ethical standards, a higher and a lower one, the former consisting in the renunciation of earthly blessings, the latter in their use; there is \emph{but one standard}, and that is the same for all. Likewise there is but one sphere in which this ethical standard is to be applied, namely the several forms of the earthly calling or vocation which are divinely ordered and therefore well-pleasing to
\quellseite{131}
God. For this reason Luther supplied the Table of Domestic Duties with the superscription: “Table of Duties, Certain Passages for Various Holy Orders and Estates.”

While the great significance of the Small Catechism is largely found in the second meritorious feature just considered—the evangelical interpretation of the individual parts, this book would never have gained or maintained such significance but for the additional merit of possessing \emph{great pedagogic excellences}. It would never have attained such wide-spread use as a school book merely because it was written by Luther or acknowledged as one of Symbolical Books of the Church or regarded as a truly evangelical treatise, if it had not possessed a number of pedagogic excellences.

Of the pedagogic merits we mention first its \emph{consummate linguistic form}, the all but architectural beauty of its several parts. From this point of view the explanation of the Second Article appears as a hitherto unparalleled model, characterized by rhythmic euphony and a noble plastic form. Read it aloud, by way of illustration, and note the word painting so truly in harmony with the deep meaning especially in the passage “who has redeemed me, a lost and condemned creature, purchased and won me from all sins, from death and from the power of the devil.” Observe the miniature word painting (\emph{Kleinmalerei}), born out of the principle of intuition and demanded by it, when he unfolds the term “daily bread” in the explanation of the Fourth Petition. Here and there indeed a lengthy period is found which is memorized with difficulty; soon after publication some pedagogues have therefore attempted to alleviate this task by resolving such periods into shorter sentences,\footnote{See Reu, \emph{Quellen} i 1, p. 516, cf. p. 198, 483; i 2, pp. 165 ff.} but these attempts are anything but improvements; besides, as a rule Luther’s periods are lucid and architectonic, and when they have once
\quellseite{132}
been committed to memory, they constitute a treasure for the mind which will prove far less elusive than a series of asynthetical sentences. The glow of personal testimony which breathes from these lines, the picturesque concreteness of expression which at the proper time relieves pithy succinctness, the hymn-like sweep of phrase minimize these difficulties and cause them to be all but forgotten.

Other pedagogic excellences are even greater. We name, secondly, the \emph{absence of any polemics whatever}. However fiercely Luther fulminated against false teachings, not recoiling from an occasional vehement onslaught even in his sermons, just as assiduously he avoided every trace of polemics in the Small Catechism where in spirit he stood in the nursery or in the school room. He does not inveigh in the First Chief Part against the spurious good works conceived by men; but merely aims to show at the hand of suitable examples how genuine good works must grow forth from the fear and love of God. He does not battle in the Second Article against those who have usurped authority over Christendom, as pope and bishops; but is intent on one thing alone, namely to impress firmly the positive and fundamental truth that Jesus alone is our Lord and Master, and that we simply cannot own anyone else as master if we remember under what circumstances he became our Lord. In the Third Article he does not contend against Romanist or enthusiast perversions of the way of salvation; but he emphatically impresses the fact upon the heart that only the Holy Spirit leads us to Christ and that he uses no other means than the gospel. No doubt, Luther was most strongly tempted to polemics when discussing Baptism and Communion. At the time of writing the Small Catechism, the Anabaptists and enthusiasts had been checked in Saxony, but the havoc made by them was still plainly visible everywhere, and outside of Saxony they were at work undismayed
\quellseite{133}
and not without success. Nevertheless, Luther did not discuss their pernicious doctrines; only the third question on Baptism and its answer contain, as it were, faint echoes of the strife, but also here the thought is given a purely positive turn. The sacramentarian conflict was raging when Luther wrote the Fifth Chief Part, but he refrains from polemicising and merely presents with special care the biblical doctrine of the Lord’s Supper; he knows that the more positively boys and girls are grounded in biblical truth, the more constant they will be when the time comes to maintain and to defend their biblical faith. Luther banishes all cold scholasticism and polemics from the school room, he would rather let youth feel the breath of true noble godliness and the throb of personal faith and life; for he would find the way to the heart of the young and exert permanent plastic influence upon their souls. A breath of the life from God and in God—such is Luther’s catechetical motive.

A third pedagogic excellence of the Small Catechism is the fact that, in distinction from most of his predecessors and successors, Luther \emph{refrains from combining the five chief parts into a systematic organism} through questions connecting the parts with one another. In the Short Form of 1520 Luther too had established such a connection at least so far as the Trilogy is concerned, in the following renowned passage: “Three things a man must needs know in order to be saved. First he must know what to do and what not to do. Secondly, when he recognizes that by his own power he can not do what he should nor leave undone what he should not do, he must learn where to receive, seek, and find such power and strength. Thirdly, he must know how to seek and obtain it. Just as a sick man must know first of all what his sickness is and what he may do and what he may not do. Next he must find where the remedy is which may cure him so that he can do what any healthy person does.
\quellseite{134}
Thirdly, he must desire such remedy and seek and fetch it or have it brought to him. Accordingly the Commandments teach man to know his sickness so that he may see what he can do and what he cannot do, and to acknowledge himself as a sinner and a wicked man. After that faith shows and teaches him where to find the remedy that helps him to become godly and to keep the commandments; it shows him God and his mercy, revealed and offered in Christ. Thirdly, the Lord’s Prayer teaches him how to seek, fetch, and obtain such mercy, namely through proper humble confident prayer: thus it shall be granted him, and he shall be saved through the fulfillment of God’s commandments.” Similar statements are found in the Large Catechism.\footnote{See Jacobs, \emph{Book of Concord}, p. 439 and 448: \emph{Tr.} p. 679 and 697.} But in the Small Catechism written \emph{pro pueris et familia}, he does not only refrain from any such allusion, but excludes every attempt at reducing its material to a system by treating the Decalogue both as the norm of the new life (First Chief Part) and as a mirror of sin (Confession). For, classification into a system becomes possible only when the Ten Commandments are considered either as the norm of the new life or as a mirror of sin; it is impossible when they are treated as both at one and the same time. To Luther, each chief part represents the whole of Christianity, but every time from a different point of view. With sound pedagogical insight Luther recognized that young people and ordinary folks need clear and compact statements, but no system.

Of still more importance is the fourth pedagogical merit of the Small Catechism: its \emph{wise restriction to the essentials of Christian faith and life} and the greatest possible elimination of the technical terms of dogmatics. More than all the catechists that preceded him, more than almost all the catechists that followed, Luther remained conscious of writing not for prospective theologians, but for the youth
\quellseite{135}
and the common people. He restricted himself to the actual food needed by the Christian as a child of God, to the actual requirements for a Christian life and a blessed death. He might have dilated upon the divine essence and attributes, upon the wonderful interpenetration of the divine and human natures in Christ down to an elaborate statement concerning the \emph{communicatio idiomatum}; he might have brought in all the finest shadings of the several stages of the way of salvation, or all the manifestations of Christian morals in the diverse forms of social life. On other occasions he would indeed blaze the path in these respects, but in his Catechism he was silent about all this in order to bring out what is central, to make it clear in all its aspects, and to impress it upon the young and ordinary people. He knows that otherwise they “would not see the forest for its trees.” It was this exquisite pedagogic tact no less than his deep evangelic understanding that prompted Luther to reduce Christian morals to one root, the fear and love of God; and to group all the creedal facts around the three outstanding deeds of creation, salvation, and sanctification. Luther, always active and vigorous, knew how to restrict himself to essentials; it was Melanchthon, the theorist, who introduced dogmatics into catechetical instruction.

It was likewise pedagogical wisdom and tact and by no means accidental when Luther refrained from stressing the negative side in the exposition of the First and Sixth Commandments. In the First Commandment, more than in any other, the negative feature as far as it has a place at all in the instruction of children, is contained in the positive one, especially in the phrase “above all things”; when teaching children, “gross idolatry” must not be discussed, and by stressing the negative element one is likely to displace the positive element from its central position and to obscure it. And the omission of the negative side of the Sixth Commandment is to be commended because thus the children are not made acquainted with sins which hitherto may have been strange to them.

\quellseite{136}
As far as \emph{method} is concerned, the merits of Luther’s Small Catechism are much smaller. Because Luther employed questions and answers in his Catechism, it has been claimed that Luther is the author of the erotematic form of instruction; some even would make him the father of the didactic conversation. Such assertions, however, are wholly unfounded. The erotematic form of instruction had been used long before Luther, far in the remote Middle Ages; and the questions employed in the Small Catechism are in reality examination questions or confessional questions, not developing questions in the sense of modern pedagogics. The didactic conversation as an exact educational method can be traced to the Reformation only in its faintest beginnings. The pronunciation and repetition of sentences was at that time deemed quite sufficient. The father or the teacher would pronounce a sentence, and the child would repeat it until the subject matter was appropriated. Nevertheless, the question may have served to arouse the attention of the learners.

From the standpoint of method it is noteworthy that Luther added pictures, chiefly biblical, to his Catechism. Thereby he exemplified that he adhered to the important pedagogical principle of intuition. Also in his preface he advises teachers to adduce as many illustrations from the Scriptures as possible.

Another point is worthy of imitation: he insists upon a gradual introduction into the catechetical material. He advises the pastor first to take up only the text proper of the Catechism for the purpose of impressing it upon the memory. In order to avoid confusion, the teacher should refrain from introducing changes into the text of the Catechism but leave it unaltered once for all. When the children have become familiar with the words of the text, the pastor must proceed to the explanation so that the children
\quellseite{137}
may understand the meaning of the Catechism. For this purpose he should use “these tables,” i.e., Luther’s explanation, or “some other brief explanation”; and again Luther adds the admonition to abide by the form once selected without altering even a single syllable. When this explanation has been mastered, the “Large Catechism” shall follow, i.e., an ampler explanation, as found in “a great number of books” among which, of course, also Luther’s Large Catechism though left unmentioned is to be reckoned. By means of this large Catechism the pastor shall furnish the children and the plain man “with fuller and more comprehensive explanation,” and Luther advises him in a special manner to insist on such commandments or other parts as are neglected or misunderstood by his people. But if the pastor because of the lack of response on the part of the people and because of their ingratitude is in danger of losing his courage and joy, Luther reminds him that “Christ himself will be our reward if we labor with fidelity”; those, however, who refuse to receive instruction are to be told plainly that “they deny Christ and are not Christians; such persons shall not be admitted to the Lord’s Table, nor present a child for baptism, nor enjoy any of our Christian privileges, and are to be sent back to the Pope and his agents, and, indeed, to Satan himself.”
\quellseite{138}
\section[§ 15. Religious Instruction, 1530–1600]{§ 15. Religious Instruction in the Church Between 1530 and 1600}

E. Sachsse pp. 201 ff. E. Chr. Achelis, pp. 303 ff., 368 ff., 401 ff. G. Schumann u. E. Sperber, \emph{Gesch. d. Religionsunterrichts in der Volksschule}, pp. 35 ff., 1890. G. Langemack, \emph{Historia catechetica} ii, 1733. K. J. Loeschke, \emph{Die religioese Bildung der Jugend u. d. sittl. Zustand d. Schulen im 16. Jahrh.}, 1846. C. Moenckeberg, \emph{Die erste Ausg. v. Luth. Kleinem Katechismus}, 1851, ²1868. J. Brustlein, \emph{Luthers Einfluss auf d. Volksschulwesen u. d. Religionsunterricht}, 1852. F. Ehrenfeuchter, \emph{Zur Geschichte d. Katechismus}, 1857. H. Ernesti, \emph{Zur Orientierung ueber die Katechismusliteratur}, 1859. E. Sachsse, \emph{Zur Gesch. d. ev. Katechismus} in \emph{Halte, was du hast}, 1889. M. A. Gooszen, \emph{De Heidelbergsche Catechismus}, 1890. A. Ernst u. J. Adam, \emph{Katechetische Geschichte des Elsasses}, 1897. *Friedrike Fricke, \emph{Luthers Kleiner Katech. in s. Einwirkung auf die katechetische Literatur d. Reformationsjahrhunderts}, 1898. R. Neumann, \emph{Der ev. Religionsunterricht im Zeitalter d. Reformation}, 1899. Th. Wotschke, \emph{Brenz als Katechet}, 1900. F. Cohrs, \emph{Katechismen u. Katechismusunterricht in PRE³}, 1901. K. Knoke, \emph{Die deutsch-luth. Katechismen i. d. braunschweig-hannoverschen Landen waehrend d. 16. Jahrh.} in \emph{Zeitschr. fuer niedersaechs. Kirchengeschichte}, 1901, 1905. *M. Reu, \emph{Quellen z. Geschichte d. kirchl. Unterrichts im ev. Deutschland zwischen 1530 u. 1600}, 8 vols., 1904-1926. O. Frenzel, \emph{Zur katechetischen Unterweisung i. Zeitalter d. Reformation u. Orthodoxie}, 1915. M. Lauterburg, \emph{Der Heidelberger Katechismus in PRE³}, 1901. A. Lang, \emph{Der Heidelberger Katechismus u. vier verwandte Katechismen, mit einer historisch-theologischen Einleitung} (here more literature concerning the catechisms of the Reformed Church), 1907. Chr. Moufang, \emph{Katholische Katechismen d. 16. Jahrh. i. deutscher Sprache}, 1881. P. Bahlmann, \emph{Deutschlands kath. Katechismen bis z. Ende d. 16. Jahrh.}, 1894. O. Braunsberger, \emph{Entstehung und erste Entwickelung der Katechismen des Petr. Canisius}, 1893. F. X. Thalhofer, \emph{Entwicklung d. katholischen Katechismus in Deutschland}, 1899. A. Richter, \emph{Ev. Kirchenordnungen d. 16. Jahrh.}, 2 vols., 1846. H. Heppe, \emph{Gesch. d. deutschen Volksschulwesens}, 5 vols., 1858-1860. *A. Vormbaum, \emph{Die Schulordnungen d. 16. Jahrh.}, 1860. F. Paulsen, \emph{Gesch. d. gelehrten Unterrichts vom Ausgang d. Mittelalters}
\quellseite{139}
 i, ²1896. *G. Mertz, \emph{Das Schulwesen d. deutschen Reformation i. 16. Jahrh.}, 1902. E. Sehling, \emph{Die evang. Kirchenordnungen d. 16. Jahrh.}, 5 vols., 1903 ff. *Fr. Seefeldt, \emph{Zur Entstehung d. biblischen Geschichts-unterrichts im deutschen Protestantismus}, 1924.

When Luther laid before the Church his Catechism, he was convinced that it met the needs of \emph{oral instruction} of the common people. The very name “Catechism” which he had chosen for this booklet indicates that it was not intended to be anything but a set of oral lessons in literary form. But it was far from his mind to have produced a standard book which must be used in all evangelical churches and schools of Germany. A considerable number of independent catechetical publications made their appearance even after Luther had issued his catechism; and this applies not only to Southern Germany where, e.g., the influential preachers of Strassburg published their catechisms (Butzer in 1534, 1537, and 1543; and Matthaeus Zell in 1535, 1536, and 1537), or to sections largely Reformed, but even to purely Lutheran territory. It is true that most Lutheran catechisms, as those by Moibanus Rhegius, Melanchthon, Lossius, Trotzendorf, Chytraeus, Wigand, were intended for the higher classes of the Latin schools and therefore composed in Latin; but several independent text-books, as those by Chr. Vischer, Hermann Bonnus, Judex, and others, were issued for the primary classes of the German schools and for the instruction of the common people in general. Gradually, however, Luther’s Small Catechism gained predominance not only because Luther was its author but chiefly because it surpassed all others in excellence and because it formed a suitable counterpart to the Heidelberg Catechism (1563) which had welded the various branches of the Reformed Church into a unit. It was only in the upper classes of Latin schools that systematic text-books of Christian doctrine could maintain themselves. Eventually all codes of regulations for churches and schools prescribed the study of Luther’s Catechism
\quellseite{140}
; the Lutheran \emph{Corpora Doctrinae}—these precursors of the Book of Concord of 1580— incorporated it, and after the appearance of the Crypto-Calvinistic \emph{Wittenberg Catechesis} (1571) it found a place among the Symbolic Books of the Church in the Book of Concord. In the schools it was used as First Reader: it contained the alphabet and some reading exercises and offered, in accordance with Luther’s advice, first the catechetical texts and then the texts with Luther’s explanation. In the Latin schools the same was done, but the medium of instruction was, of course, Latin. The advanced pupils, especially in the Latin schools, were given a rather copious analysis and exposition of Luther’s explanation. Among the best of such explanations of Luther’s Catechism in the German language are those by Moerlein (1547), Friedrich (1572), Mathesius (1574), and Bischoff (1599). The most excellent explanation is John Tetelbach’s \emph{Gueldene Kleinod} (Golden Jewel) of the year 1568; here an expository principle of great importance is correctly applied, viz., that an explanation of the Catechism must introduce into, and unfold, the treasure of religious truth contained in the \emph{verbal contents} of Luther’s Catechism. Tetelbach’s booklet was widely used far into the seventeenth century. In the higher classes of the Latin schools the Catechism was taught in Greek; even a polyglot edition (German, Latin, Greek, Hebrew) was issued, and not only once but repeatedly. A German-Latin edition had been prepared soon after the publication of the German Catechism and enjoyed great popularity. Yes, it happened that even in universities it formed the basis for the discussions. It was not long till Luther’s Catechism had spread beyond the boundaries of Germany. Within about 60 years after its appearance it was found in Bohemian, Danish, English (in the form of the Nuremberg Sermons for Children), Finnish (?), Esthonian, French, Greek (three different translations), Hebrew, Icelandic (in
\quellseite{141}
the form of the Nuremberg Sermons), Italian, Lettish, Lithuanian, Latin (four translations), Dutch, Old-Prussian, Slovenian, Spanish (?), Swedish, and Wendish. In but few regions the Catechism of Brenz, in the briefer form of 1535, maintained itself for some time, in Wuerttemberg down to the present day in the form of 1681 which combines with it Luther’s Catechism. Recent specific investigations in this field have disclosed the faithful energy with which the Church in this period everywhere disseminated the Catechism among the young and common people.

The Church, however, did not only provide the needed text-books, but also the institutions required for teaching. Again and again, the home was reminded of its God-given duty of instruction with the result that the inculcation of the Catechism became an important element in family-worship; children and servants would often know the Catechism by memory without, or before, being able to read. In the Latin schools the inculcation and explanation of the Catechism was a regular part of the curriculum, and the school regulations almost without exception contained a specific provision concerning such instruction. In the villages it became an established rule, even beyond Saxony, that the vergers (cf. p. 90) impart religious instruction. Here we have the faint beginnings—not more than that—of the public school. Regular inspections contributed toward preserving and improving these institutions; one code of regulations even required the vergers to give four hours’ instruction daily which was, of course, an excessive demand and for the time being simply could not be enforced. Catechetical instruction was given also in Church: every Sunday the Catechism or one of its Chief Parts was read from the pulpit, special catechetical lessons were conducted on Sunday afternoons, now and then also on week-days, and catechetical sermons were preached during the ember seasons or throughout the
\quellseite{142}
year on Sundays or week-days. All the extant catechetical sermons intended for children are excelled by the above mentioned Nuremberg Sermons for Children; they were translated into several languages (into English in 1548), and in some respects may serve as models even today. When in these institutions the children had learned to understand the Catechism, an examination, called \emph{Katechismus examen} or \emph{Glaubensexamen}, was conducted (as a rule during Lent and usually in the parsonage) which determined whether or not the children might be admitted to Holy Communion. This examination was regarded as the genuine evangelical confirmation. At some places it was preceded by a special course of instruction ranging over several weeks (Schleswig-Holstein 1544, Brandenburg-Ansbach 1564). Gradually public examination before the assembled congregation supplanted the more private examination in the parsonage; it was connected with a simple confession on the part of the children and prayer on the part of the congregation (\emph{exploratio, confessio, oratio}). In this form it was prescribed, for the first time, in 1535 for Liegnitz, and proposed in the “Wittenberg Reformation” of 1545; later it spread also to distant places, but it was by no means universally accepted (Electoral Saxony, e.g., did not introduce it). Chemnitz, in \emph{Examen Concilii Trident.} ii 3, strongly recommended it as the true evangelical counterpart of the Catholic confirmation. In Hesse (1539, 1566, 1574) and in Strassburg (at least since 1543), owing to Butzer’s influence, the imposition of hands was connected with this act and also the vow of obedience to the Church of Christ. Even sacramental character was imputed to confirmation since the reception of the Holy Spirit was identified with the imposition of hands; fortunately, however, this form of confirmation did not spread very far. Further details concerning this point will be found in § 39.

The first communion thus clearly developed as goal of
\quellseite{143}
instruction. It must be admitted, however, that in many cases rote learning of the five Chief Parts, or even of the catechetical texts, was deemed sufficient. At other places the five Chief Parts were supplemented by special questions pertaining to confession and the Lord’s Supper.

Finally, private confession which was practiced almost universally (cf. Augustana xi) enabled the pastor to revive his parishioners’ knowledge of what they had learned in early youth in as much as he examined them in the five Chief Parts of the Catechism; the same was done, here and there, in the examination preliminary to marriage. Now and then circumstances rendered necessary an examination, in private confession, of all adults, but as a general rule to be applied everywhere it could not be justified because it would militate against the Evangelicals’ aim of all educational efforts, viz., the maturity of the Christian congregations.

In the era under discussion also the other tendencies developed that had manifested themselves before and since 1529. In many Latin schools \emph{Biblical History} was taught although it was not yet placed on the same level with the Catechism. Castellio’s excellent \emph{Dialogi Sacri} were edited more than sixty times between 1543 and 1600—although not always for the sake of instruction in \emph{Biblical History}, and eleven editions of Fabricius’ \emph{Historia Sacra} were issued between 1564 and 1600. The same purpose was served by the study of Pericopes which had been introduced in many schools as part of the regular curriculum; in fact, instruction in the pericopes which had been introduced in nearly all Lutheran regions in the 16th century and which was imparted generally, though not exclusively, on Saturdays, was the form of instruction in Bible History in the 16th and 17th centuries. Formerly it was customary among historians to deny that instruction in Bible History was given in this period; but recent research has revealed the error of this view.
\quellseite{144}
More yet, it has been shown that also those children who did not attend Latin schools, frequently studied the pericopes, i.e., Bible History; occasionally the pastors would even explain the pericopes in the churches to children and young people. Furthermore, \emph{collections of Bible Passages} were used in lower as well as higher schools; some of these collections followed the plan of the Catechism, some the order of pericopes. Aside from the feeble beginnings in the twenties, the earliest selections in Latin were prepared by Lossius of Lueneburg and by Trotzendorf in Silesia; Veit Dietrich of Nuremberg composed the first German selection for children (1546). We also find school editions of the Proverbs of Solomon, the Psalms, and Ecclesiasticus. Only by recent research has this wealth of catechetical literature been rescued from oblivion. The young were even taught to sing \emph{hymns} and had to commit some of them to memory; we find hymnals specially designed for children and also catechisms with an appendix containing the chief \emph{hymns}. The diligent and efficient work done between 1530 and 1600 is simply astounding and, in part at least, a model for today.

Faithful work was done also within the \emph{Reformed Church}. As early as 1522, Zwingli instituted a sort of \emph{Kinderlehre} in Zuerich. Here Leo Jud, adopting Luther's arrangement of the Chief Parts, wrote his Catechism in 1534 which was issued in abridged form in 1535. For the church of Geneva Calvin composed one French Catechism in 1537 and another one in 1541 which in the Latin translation of 1545 and in the German translation of 1563 was employed in various districts of Germany. Consisting of 373 questions and answers it was divided into four parts: (1) Faith; (2) the Law; (3) the Lord's Prayer; (4) the Word of God and the Sacraments. It was further divided into 52 sections for the purpose of studying it completely in one year. John a Lasco's Catechism of 1546 (printed in 1551) is exceptionally lucidly arranged
\quellseite{145}
 and in exemplary manner attempts to arouse faith; it contained the following parts: (1) Decalogue; (2) Creed; (3) Baptism, Communion, and Church Discipline; (4) Prayer. All these catechisms were superseded by the Heidelberg Catechism which by electoral rescript was introduced in 1563 in all churches and schools of electoral Palatinate; since the Synod of Dort (1618) it has enjoyed almost symbolical dignity throughout the Reformed Church. Composed by the Heidelberg theologian Ursinus in collaboration with others, notably Olevian, this catechism represents a skillful summary of the previous catechetical efforts in the Reformed Church. It consists of three parts: (1) Man’s Misery (question 1-11); (2) Man’s Redemption (questions 12-85; Creed and Sacraments, including the notorious eightieth question); (3) Man’s Gratitude (questions 85-129; Decalogue and Lord’s Prayer). This famous threefold division was patterned after a catechism by a Lutheran author.\footnote{See M. Reu, \emph{Quellen etc.} i 1, pp. 201–203.} The Heidelberg Catechism is a fair statement of the milder type of Reformed teaching, and here and there rises to heights of linguistic power and beauty, particularly in the famous first question; but its lengthiness and clumsiness and especially the hairsplitting and spinous nature of many of its questions are decided disadvantages. In later editions the proof texts which in the original had been merely indicated, were printed in full and the whole was divided into 52 lessons so as to be covered in the course of 52 Sundays. An abridged edition for children was issued in 1585. This catechism marks the conclusion of catechetical efforts within the Reformed Church in Germany. The needed institutions for the teaching of the catechism were also provided, among them the catechetical examination of old and young on Sundays.

By these educational efforts of the reformatory Church even the Roman Church was induced to make some earnest
\quellseite{146}
endeavors toward providing more general and thorough instruction for the young. A considerable number of catechisms appeared in the sixteenth century, especially since the Jesuits (1540) began to devote themselves to this work. Canisius, the father confessor of Ferdinand I, deserves special mention. In 1556 he issued a small catechism (i.e. a synopsis of his \emph{Summa doctrinae Christianae}) under the title \emph{Summa doctrinae Christianae per quaestiones tradita et ad captum rudiorum accommodata}; soon it was issued also in German (\emph{Der Klain Katechismus} etc.). For two centuries this book was the chief catechism of the Roman Church. In 1559 Canisius issued a catechism for higher schools of learning which despite its bulky size bore the title \emph{Parvus catechismus catholicorum}; it became the standard exposition of the smaller Summa. Even the \emph{Catechismus Romanus} (also called \emph{Catechismus Tridentinus}) though composed by order of the Council of Trent and intended as a manual for priests, or the text books subsequently based on this catechism and designated for exclusive use (1598; pope Clement VIII) did not succeed in dethroning the text books by Canisius.

\quellseite{147}
\section[§ 16. Lutheran Orthodoxism]{§ 16. Catechetical Instruction During the Period of Lutheran Orthodoxism (1600–1676)}

G. Langemack, \emph{Historia catechetica} ii. iii, 1733. 1740. Ph. Schuler, \emph{Gesch. d. katechetischen Religionsunterrichts}, pp. 93 ff., 1802. *F. Ehrenfeuchter, \emph{Zur Geschichte d. Katechismus}, 1857. F. W. Bodemann, \emph{Katechetische Denkmale der ev.-luth. Kirche}, (containing the catechism of Tetelbach 1568, for Nuremberg 1628, of Justus Gesenius 1639, of M. Walther 1653, of Glassius 1640, for Mecklenburg 1711), 1861. *G. v. Zeschwitz ii 1. Schuman-Sperber, pp. 39-49. *E. Sachsse, pp. 186 ff. M. v. Nathusius i, pp. 57 ff. E. Chr. Achelis, pp. 319 ff. J. Steinbeck, pp. 37 ff. A. Ernst u. J. Adam, Katechetische Geschichte des Elsasses, 1897. Buenger, \emph{Entwicklungsgeschichte des lutherischen Katechismusgebrauchs in Hannover}, 1912. *O. Frenzel, \emph{Zur katech. Unterweisung im Zeitalter der Reformation und Orthodoxie}, 1915. *Vormbaum, \emph{Ev. Schulordnungen des 17. Jahrhunderts}, 1863. F. Paulsen, \emph{Geschichte des gelehrten Unterrichts}, 21896. F. Sander, \emph{Geschichte der Volksschule}, besonders in Deutschland in K. A. Schmidt, \emph{Gesch. d. Erz. v 3}, 1902. A. Heubaum, \emph{Gesch. d. deutschen Bildungswesens seit Mitte des 17. Jahrh.}, 1915. A. Matthias, \emph{Gesch. d. deutschen Unterrichts}, 1907. Joh. Arndt; F. Hashagen, \emph{Ausgew. Predigten} in \emph{Predigt der Kirche} xxvi, 1894. G. Hoelscher, \emph{J. Arndt in PRE³}, 1897. W. Koepp, \emph{J. Arndt, e. Untersuchung ueber d. Mystik im Luthertum}, 1912. Joh. Val. Andreae: J. Ph. Gloeckler, \emph{Joh. Val. Andreae, e. Lebensbild}, 1886. J. Bruegel, \emph{Joh. V. Andreae in K. A. Schmidt, Gesch. d. Erz. iii 2}, 1892. G. Hoelscher, \emph{J. V. Andreae in PRE³}, 1896. C. Huellemann, \emph{V. Andreae als Paedagog}, 1884. His \emph{Theophilus}, written 1622 (prior to Comenius!) and printed in 1649, may be found in German in Oehler, J. V. Andreae's \emph{Theophilus}, 1878. His \emph{Evangelische Kinderlehre} upon which the Booklet for Confirmands of Wuerttemberg is based, was published 1621; the third edition (1648) is reprinted in Ph. Schuler, op. cit., pp. 329-352. Wolfgang Ratichius: G. Voigt, \emph{Wolfg. Ratichius, der Vorgaenger d. A. Comenius}, 1894. Joh. Mueller, \emph{Handschriftliche Ratichiana in Paedag. Blaetter}, 1878 ff., esp. 1880, pp. 497 ff. J. Lattmann, \emph{Ratichius und die Ratichianer}, 1898. A. Prall, \emph{Paedag. Schriften d. W. Ratichius}, esp. pp. 61 ff., 1902. H. Barnard, \emph{German Teachers and Educators}, pp. 343-370. R. H. Quick, \emph{Educational Reformers}, ch. ix, 1896. Joh. Amos
\quellseite{148}
Comenius: S. S. Laurie, \emph{A. Comenius, His Life and Educational Works}, 1885. W. S. Monroe, \emph{Comenius}, 1900. *J. Kvasala, \emph{J. A. Comenius, s. Leben u. s. Schriften}, 1892. J. Kvasala, \emph{Die paedag. Reform d. Comenius in Deutschland bis z. Ausg. d. 17. Jahrh.}, 1904. F. Pappenheim, \emph{J. A. Comenius}, 2 vols., 1902. 1905. \emph{Monatshefte d. Comeniusgesellschaft}, 1892 ff. His \emph{Didactica Magna} has been translated into German by Lion, ⁵1904, into English by Keatings, 1896; \emph{Schola ludus} ed. by Boetticher, ²1907; \emph{Mutter Schul} ed. by Lion, ²1907, in English: W. S. Monroe, \emph{The School of Infancy}, 1896; \emph{Orbis pictus} ed. by C. W. Bardeen, 1891. Joh. Kromayer: L. Weniger, \emph{Ratichius, Kromayer u. der Neue Methodus in Weimar} in Zeitschr. f. thuering. Geschichte, 1896. L. Weniger, \emph{J. Kromayers Schulordnungen von 1614 u. 1617}, ibid., 1901. Chr. Waas, \emph{Die Generalvisitation . . . von 1641 bis 1647}, ibid., esp. pp. 245 ff., 1909. Sigm. Evenius: P. Stoetzner, \emph{Sign. Evenius, e. Beitrag zur Gesch. d. Ratichianismus}, 1895. G. Schmid, \emph{S. Evenius in Monatshefte d. Comeniusgesellsch.}, 1895. K. Knoke, \emph{E. christl. gottselige Bilderschule} in Katech. Zeitschrift, 1907. Cf. also Lattmann, \emph{Ratichius u. d. Ratichianer}, 1898. Ernest the Pious, Duke of Gotha: G. Kreyenberg, \emph{Ernst d. Fromme, ein Lebens- und Kulturbild des 17. Jahrh.}, 1890. H. Petrich, \emph{Herzog E. d. Fr., s. Leben und Wirken}, 1901. W. Boehne, \emph{Die paedagogischen Bestrebungen Ernsts d. Fr. v. G.}, 1888. H. Schroedel u. H. Moeller, \emph{E. d. Fr., ein Paedagog unter d. Fuersten}, 1901. M. Ehr, \emph{Beitr. z. Kirchen- und Schulverfassung d. Herzogtums Gotha bis z. Tod Ernst des Fr.}, 1891. *Braem, \emph{Der gothaische Schulmethodus}. \emph{E. kritische Untersuchung ueber die ersten Spuren d. Pietismus i. d. Paedagogik d. 17. Jahrh.}, 1897. The Methodus of 1642 has been edited with an excellent critical and historical introduction by Johannes Mueller in \emph{Sammlung selten gewordener paedagogischer Schriften frueherer Zeiten}, No. 10, 1883; it bears the title: “Herzog Ernst des Frommen Spezial und sonderbarer Bericht, wie die Knaben und Maegdelein auf den Dorfschaften und in den Staedten im Fuerstentum Gotha kurz und nuetzlich unterrichtet werden koennen und sollen”. The Methodus of 1672 has been edited by A. Prall, \emph{Der Schulmethodus des Herzogs Ernst v. Gotha}, ²1912. S. Glassius’ \emph{Kurzer Begriff der christlichen Lehre} is reprinted by Ehrenfeuchter, op. cit., Supplem. pp. 41-62. Michael Walther: Ehrenfeuchter, op. cit., pp. 82 ff. R. Buenger, op. cit. His Catechism is reprinted by Ehrenfeuchter, op. cit., supplem. pp. 89-126, and by Bodemann. Justus Gesenius: Ehrenfeuchter, op. cit., pp. 79 ff. G. v. Zezschwitz ii 2,1, pp. 93 ff., 1874. Buenger, op. cit. E. Bratke, \emph{Just. Gesenius, sein Leben u. s. Einfluss auf d. Hannoversche Landeskirche}, esp. pp. 44 ff.; 82 ff.; 146 ff., 1883.
\quellseite{149}
C. Bertheau, Justus Gesenius in PRE³, 1899. K. Kayser, Die General-visitation des Just. Gesenius in Zeitschr. fuer niedersaechs. Kirchen-gesch., 1906. His Catechism is reprinted by Bodemann, op. cit., and in parts by Ehrenfeuchter, op. cit., supplem. pp. 62 ff.

When in opposition to its obscurations and perversions, Lutheran doctrine had been confessionally fixed in the Book of Concord of 1580 or more specifically in the Formula of Concord, the temptation was very enticing to work all this doctrinal material into the catechetical text books though they were not designed for theologians but for children. Generally speaking, the Church succumbed to this temptation. As a result, the doctrinal element was so preponderant in the catechisms of the time that the religious and practical elements were often completely overshadowed. An abhorrent example of this kind is found in Conr. Dietrich’s Institutiones Catecheticae of 1613. In itself it is an excellent book, helpful alike to student and pastor; but the author’s method of dragging into it every technical term of dogmatics and history of dogma has had baneful consequences; even today we are not quite free from this injurious practice. In catechetical instruction the catechist found a welcome opportunity to review his dogmatics. This proceeding may have been salutary for himself, but the children were given stones rather than bread. Even the epitome of Dietrich’s Institutiones designed especially for school use and employed in some regions to our day, is far too doctrinaire. The same criticism applies to the catechetical sermons of this period. The explanation of the Fifth Chief Part in Aegidius Hunnius’ catechetical sermons would be a highly creditable university lecture; others even surpass Hunnius in this respect. Catechetical sermons like those of John Arndt (1620) are rarely found in this era. Noble men were grieved by these conditions; it was especially Valentin Andreae, his soul on fire for the welfare of the Church, who made bitter complaint of the parrot-like recitation of a catechism and an explanation
\quellseite{150}
which were not understood, and severely scourged such practice in his Theophilus (1622). To improve these deplorable conditions he published his Kinderlehre (1621). Compare also § 35.

Then came the Thirty Years’ War with its woeful, its unspeakable devastation of church and school. The educational institutions that had developed in the course of time were dangerously undermined and some of them disappeared entirely, for instance the preparation of the children for their first communion. But God’s dire visitation through war once more filled the hearts with a desire to care for the country’s youth so shockingly neglected during the incessant tempest of battle. Thus a revival of education came to pass in the midst of the turmoil of war. As early as 1628 Nuremberg issued the equally instructive and edifying \emph{Kinderlehrbuechlein} (Booklet for the instruction of children); it must be admitted, however, that this booklet promoted the one-sided evaluation of the Decalogue as a mere mirror of sin by concluding the explanation of each commandment with the question “Have you kept this commandment?” Strassburg, likewise, distinguished itself by taking good care of its youth; especially J. Schmidt and J. K. Dannhauer, professors at the university, took a warm interest in this matter. Stirred by their example, faithful Justus Gesenius of Hanover devoted himself particularly to the improvement of education in his country. In his \emph{Catechetical Questions} which in the issue of 1639 exerted far reaching influence, he induces his pupils to think by changing the form of every question two or three times, and he lays stress upon the application of the truth to life. Also Michael Walther of Celle helped to raise catechetical instruction to a higher level. His explanation of the Catechism which appeared first in 1653 and was reissued in 1861 and since then frequently used in America (the so-called \emph{Stohlmann Catechism}), excels that of Gesenius in clearness
\quellseite{151}
 of definitions and exactness of method, but introduces too much dogmatical material.

In the Thuringian region it was especially general-super\-intendent Kromayer who made earnest efforts at improving religious instruction. Still better, more systematic, and more permanent work was done by the excellent duke Ernest the Pious of Gotha (1601-1675) and the splendid educators Evenius and Reyher. They succeeded in putting into practice the creative thoughts of the eminent pedagogue Ratichius († 1635) and the still greater Amos Comenius (1592-1671) and built up model schools throughout the region. The school regulations of 1642, and subsequent extensions, stipulated that in the villages all boys and girls from five to twelve years of age must attend school; the munificent duke furnished the needed equipment, in part even the school books; he insisted on the employment of efficient teachers; and he was the first to incorporate (1642, resp. 1662) scientific branches, as arithmetic and natural history, in the curriculum of the village school which formerly had included only religion, reading, writing, and singing. Thus in every respect he laid the foundation of the Christian common school of today. As regards religious instruction, Comenius had demanded that all children should learn by heart those Psalms and hymns which were most frequently sung in their respective localities, and that they should be able accurately to recite not only the Catechism but also the Bible stories and the chief passages of Scripture. In Gotha these requirements were met by teaching the children the following subjects: (1) the text of the Catechism; (2) Luther’s explanation; (3) the \emph{Kurze Begriff} by Sal. Glasius (an explanation of Luther’s Catechism); (4) 160 Bible passages; (5) Psalms, hymns, and prayers. In 1636 Evenius (?) or possible Brunchorst had published a sort of introduction into Biblical History under the title “Christian, Godly Instruction through
\quellseite{152}
Pictures, i.e., Instruction of the Younger Children in Godliness through Biblical Pictures,” but it was never employed in the schools; a little later, however, books of Bible stories were issued which at least the older children were (6) expected to read (1662) whereby Bible reading was introduced into the regular work of the school. Subsequently the children were given (7) a “Gospel Book” containing the pericopes taken from the Gospels. These lessons were read in school and explained to the children in order that they might intelligently appreciate the sermons preached on these texts. Finally (8), they were examined at school in regard to the chief contents of the sermon. During the long reign of Ernest, these regulations took root in his domain, and the Christian public schools found today in the land of the Reformation owe their origin, directly or indirectly, to his institutions.

Even in Gotha, this model country, instruction in the Catechism occupied the foreground rather exclusively; Bible instruction, however, was at least part of the curriculum. Elsewhere, too, selections of Bible passages and collections of pericopes were used in the instruction of children; and that instruction in Biblical History was not totally ignored becomes evident from the repeated publication in this era of Castellio’s \emph{Dialogi Sacri} and Fabricius’ \emph{Historia Sacra} (cf. p. 121). The above mentioned Gesenius issued in 1656 a Biblical reader for home and school under the title: “Bible Stories of the Old and New Testaments, Chronologically Arranged for the Benefit of the Young and Unlearned, Divided into Two Parts each with 54 Lessons.” In the preface Gesenius emphasizes the fact that the child’s mind craves stories, that without knowledge of Bible History neither catechetical instruction nor sermons can be properly understood, and that chronological knowledge of the stories of the Bible is indispensable if one would comprehend the progress in the economy of salvation.

\quellseite{153}
Of the better catechisms of this period we mention those of Altenburg, Quedlinburg, Frankfort-on-Oder, and Danzig. Some of these offer also selections of Bible passages and a number of Psalms and hymns. Rector Maukisch of Danzig suggested ways of reforming the method of catechetical instruction (1653), and Theo. Grossgebauer wrote on the betterment of education in general (1661).

\quellseite{154}
\section[§ 17. The Period of Pietism]{§ 17. Religious Instruction in the Church During the Period of Pietism (About 1677–1750)}

G. Langemack, \emph{Historia catechetica}, iii, 1740. Ph. Schuler, pp. 107–207. F. Ehrenfeuchter, pp. 48–61. *G. v. Zezschwitz i, ch. 26; ii 2, 1, ch. 32; ii 2, 2, ch. 16. *Schumann-Sperber, pp. 49 ff. F. Cohrs, \emph{Katechismusunterricht} in PRE³, 1901. M. v. Nathusius, \emph{Das Ziel des kirchlichen Unterrichts}, pp. 56 ff., 1903. E. Chr. Achelis, pp. 320 ff. J. Steinbeck, pp. 43 ff. *R. Vormbaum, \emph{Evang. Schulordnungen des 18 Jahrh.}, 1864. K. v. Raumer, \emph{Geschichte der Paedagogik} ii, ⁶1897. PHIL. SPENER: *P. Gruenberg, \emph{Phil. Jac. Spener}, 3 vols., esp. ii, pp. 58–90, 1893. 1905. 1906. W. Thilo, \emph{Spener als Katechet}, 1840. W. Thilo, \emph{Der Bibelspruch im Dienst des evangelischen Religionsunterrichts}, 1846. *W. Caspari, \emph{Die evangelische Konfirmation}, esp. pp. 84 ff., 1890. W. Diehl, \emph{Zur Geschichte der Konfirmation}, pp. 90 ff., 1897. A. Ernst u. E. Adam, \emph{Katech. Gesch. des Elsasses}, p. 210, 1897. Spener’s \emph{Kurze Katechismuspredigten} (i.e. the exordia of his sermons delivered during the seventies but published later), 1689; new edition, St. Louis and Leipzig, ²1869. His *\emph{Einfältige Erklärung der christlichen Lehre nach der Ordnung des kleinen Katechismus Luthers}, 1677; new edition, Berlin, ⁵1864; translated into Tamil by Barth. Ziegenbalg, 1719, revised edition Tranquebar, 1872. \emph{Tabulae Catecheticae quibus 5 capita catechismi minoris et subnexa tabula oeconomica in certa pensa distributa . . . tractantur}, 1683; in German by J. G. Pritius, \emph{Katechismustabellen}, 1713. \emph{Gedanken von der Katechismus-Information}, aus Speners Schriften von einem Freund katechetischer Uebungen zusammengetragen, 1715. \emph{Theologische Bedenken}, 4 parts, 1700–1702. \emph{Letzte theologische Bedenken}, 1711. \emph{Consilia et judicia theologica latina}, 1709; a selection of his German and Latin \emph{Consilia} was edited by F. A. E. Hennicke, 1838. His \emph{Pia Desideria} in condensed and modernized form edited by Gruenberg in \emph{Bibliothek theolog. Klass.}, 1889. AUG. HERM. FRANCKE: G. Kramer, \emph{A. H. Francke, ein Lebensbild}, 2 vols., 1880–1882. M. E. Richard, \emph{Aug. Herm. Francke and His Work}, 1897. G. F. Herzberg, \emph{A. H. Francke und sein Hallisches Waisenhaus}, 1898. O. Schulze, \emph{Franckes Paedagogik}, 1898. *A. Otto, \emph{A. H. Francke, vol. ii: Beurteilung und Bedeutung der Franckeschen Paedagogik}, 1904. *G. Kramer, \emph{Paedagogische Schriften Franckes}, ²1885. \emph{Der}
\quellseite{155}
\emph{kurze und einfältige Bericht, wie die Kinder zur wahren Gottseligkeit und christlicher Klugheit anzuführen sind} (1702), according to the edition of 1748 re-edited by O. Frick, 1889. The two important \emph{School Regulations} are reprinted by Vormbaum, op. cit., pp. 1–116. Johann Huebner: F. Brachmann, Johann Huebner, Johannei Rector 1711–1731, 1899. G. v. Zeschwitz ii 2, 1, ch. 32. Johann Jacob Rambach: C. Bertheau, \emph{J. J. Rambach} in PRE,3 1905. *\emph{Der Wohlunterrichtete Katechet}, 1722, ⁶1730. \emph{Darmstadt School Regulations} of 1733, reprinted by Vormbaum, op. cit., pp. 343 ff. \emph{Das Erbauliche Handbuechlein fuer Kinder} (1. Order of Salvation; 2. Treasures of Salvation; 3. A little Hymn-book; 4. A new Prayer-book; 5. Stories of Pious Children; 6. Rules for a Christian Life; 7. Necessary Rules for a Good Life), 1734, ⁵1735 (!). *\emph{Der Wohlunterwiesene Informator}, containing not a few modern ideas, published after Rambach’s death by Neubauer, 1737, with a good introduction in which he shows that this work as well as \emph{Der Wohlunterrichtete Katechet} have grown forth from the catechetical and pedagogic lectures which Rambach delivered at the university of Jena. Frederick the Great and the Prussian School Regulations: J. B. Meyer, \emph{Friedrich der Grosse, paedagogische Schriften und Aeusserungen}, ²1890. Huebler, \emph{Friedrich d. Gr. als Paedagog}, 21900. E. Clausnitzer, \emph{Zur Geschichte der preussischen Volksschule unter Friedr. d. Grossen}, 1901. A. Heubaum, \emph{Geschichte des deutschen Bildungswesens} etc., esp. part 4, 1905. F. Fischer, \emph{Preussische Volksschulverordnungen}, 1907. H. Kiehl, \emph{Joh. Julius Hecker}, Jahresbericht des Wilhelm-Realgymnasiums, 1908. Val. Ernst Loescher: *M. v. Engelhart, \emph{Val. Ernst Loescher nach seinem Leben und Wirken}, ²1856. P. Schulze, \emph{Zur Neubegruendung der Dresdener Volksschule im Beginn des 18. Jahrh.} Franz Wilhelm Kockel, \emph{Aus dem Leben e. saechs. Schulmanns}, pp. 151–176, 1900. W. Mueller, \emph{V. E. Loescher} in PRE,3 1902.—Compare H. Stephan, \emph{Neuzeit} (part 4 of Krueger’s \emph{Handbuch der Kirchengeschichte}), 1909. F. Uhlhorn, \emph{Geschichte der deutsch-lutherischen Kirche} i, 1911. P. Monroe, \emph{Text-book in the History of Education}, 1905. *S. C. Parker, \emph{History of Modern Elementary Education}, 1912. F. P. Graves, \emph{History of Education in Modern Times}, 1913. F. P. Graves, \emph{A Student’s History of Education}, 1915, 21917.

One frequently hears the assertion that \emph{Spener} (1635–1705), the “father of Pietism,” inaugurated a new era in the history of religious education, and that he and the other Pietists were the first to direct the attention of the Lutheran Church to the necessity of religious instruction. Such claims,
\quellseite{156}
however, can be made only by men who are wholly ignorant of the catechetical labors of the Lutheran Church between 1530 and 1600 and between 1600 and 1676. It is true, indeed, that Spener emphatically stressed the obligation of educating the youth and energetically devoted himself to that task. But others before him had done the same; and he himself repeatedly admitted that the exemplary educational institutions of Strassburg aroused his zeal. Again it is true that he conscientiously endeavored to prompt the young to apply to heart and life what they had recognized as saving truth and thus “to bring the head down to the heart” (cf. the question “What comfort and inducement to godliness do I find in this?” with which the discussion of any topic generally concludes). But he explicitly states that the catechisms of Gotha and Quedlinburg excel by virtue of such elements of edification. It is furthermore true that he insisted on using the Bible in school and on proving every doctrine from the Bible. But the same end had been in view when, long before his time, selections of Bible passages and collections of pericopes had been prepared; the only advance made by Spener in this respect is his effort to place the New Testament into the hands of the pupils so that by looking up the passages they might become more thoroughly acquainted with the Bible. Finally, it is true that it was Spener who secured general diffusion for \emph{confirmation} as the concluding ceremony of \emph{religious instruction} and as preparation for the first communion. But (a) according to Lutheran principles the chief feature of confirmation is not what Spener stresses, the solemnity proper and the vow, but the examination of the catechumens’ faith and knowledge; this examination, however, had been reestablished in most localities during the restoration period after the Thirty Years’ War. Moreover (b), it was not Spener at all who first evaluated confirmation as the solemn conclusion of religious instruction; he found
\quellseite{157}
it in that form in a Hessian village near Frankfort-on-Main where it had maintained itself, as generally throughout Hesse, since 1539, and simply adopted it. It is clear therefore that his catechetical efforts are not to be regarded as “a voice in the wilderness” but—and this applies equally well to all other complaints and counsels of Spener—rather as the dominant note in a concert of godly suggestions and desideria. The reforms which others before him and beside him had advocated and which some had long ago enacted, were energetically sponsored by Spener and increasingly carried into effect owing to the weight of his personality; but he is not entitled to be honored as “reformer of religious instruction.” His Catechism which was based on catechetical sermons and was published in 1677 under the title “Simple Explanation of Christian Doctrine” does not even compare with some of the catechisms mentioned in § 16; it is beautiful in some respects, but its merits are neutralized by an over-abundance of dogmatical and edifying matter; as a result it exceeds in bulkiness all limits set by precedent and convenience even if we keep in mind that the book was not primarily intended for children, but for parents, schoolmasters, and pastors. His Tabulæ Catecheticæ have the same fault. Though he could, at times, warmly recommend Bible reading, he did nearly nothing for the dissemination of Biblical History.

In point of method, Spener marks a slight advance. He was, indeed, no model catechist as his own catechism abundantly shows, nor was he conversant with the leading question whereby the unknown is reached through the known. But more resolutely than all predecessors he commenced to discard the learning by rote of ready-made answers, and he introduced the “analyzing” or “dividing question” by which the subject matter was at least logically divided (cf. his Tabulæ).

\quellseite{158}
At the same time the disciples of Spener, more than he himself, started out upon dangerous paths. Confirmation was exalted at the expense of baptism. According to the Pietists, the latter constitutes only the beginning of regeneration, and confirmation is the consummation of this work of the Spirit. Thus they regarded confirmation as a renewal, on the part of God and man, of the baptismal covenant. As a result, the instruction of confirmands completely overshadowed all other forms of religious instruction; moreover, the catechist was expected to lead the catechumens through the several stages of the way of salvation in such a manner that they actually experienced conversion and regeneration and on basis of their experience vowed henceforth to be God's own. In consequence of such treatment, on the one hand, the gradual introduction into the objective acts of God was dangerously curtailed, and on the other hand there resulted an undue stressing of the subjective element of inner experience and much emotional jugglery and self-deception. Certainty of salvation was grounded, not in the objective act of God, but rather in the — inconstant, vacillating — subjective feeling of grace. Another grievous consequence was that many pietists ignored Luther's Catechism altogether in their instruction and substituted for it the socalled \emph{Ordnungen des Heils} (successive stages of the order of salvation) which were rather frequently used till 1750.

Among the catechisms based upon Spener's principles and still fairly sound in doctrine, we mention the Dresden Cross Catechism which has been in use in many places in Germany and America down to the present day, and the Flensburg Catechism. The Catechism of Herford (1690) and of Mecklenburg (1717) have kept rather free from pietistic influence. The latter, although in revised form, and here and there also the former, are used to this day. On basis of Osiander's Booklet for Communicants (1590) and J. V. Andreae's \emph{Instruction}
\quellseite{159}
\emph{for Children} (1621), Hiemer issued the Booklet for Catechumens in 1723 which, in revised form, even today forms the basis for pre-confirmation instruction in Wuerttemberg while the Brenz-Luther Catechism is employed for other religious instruction in both church and school.

In part Spener, in part also Gotha where he had been reared, and unquestionably also some opponents of dead orthodoxy influenced \emph{August Hermann Francke}, the eminent educator and theologian whose activity was a great blessing for the cause of school and religion throughout Germany. His orphans’ school soon developed into a veritable model. In accordance with the Gotha regulations, he made room for scientific branches; children, he thought, should not only write a good hand, but also be able to compose letters independently, they must know something about their natural environment, about geography and about history, boys ought to learn the elements of surveying, and boys as well as girls must be able to knit and patch. One half of the school time, however, he devoted to religious instruction; hence we meet here all the features which Ernest of Gotha had introduced in his country, also a number of improvements: more stress was laid upon training in prayer; the catechism was analyzed by questions and consistently applied to life; sections from the Bible were read and explained daily; surveys of the more important books of the Bible and of the whole Bible were given; the children were instructed in the private devotional use of the Bible; and in the Paedagogium, the higher school, Castellio’s and Fabricius’ books were used as the basis for regular instruction in Biblical History. From all parts of Germany pupils flocked to the institutions at Halle to receive their normal training in the Paedagogium or Normal School. Teachers who had been educated in Francke’s schools were found in every section of the country, and additional Normal Schools patterned after Halle were
\quellseite{160}
established in other states of Germany. Francke’s views were given widest publicity and in most cases cordially received. In fact, when they were largely incorporated in the Prussian Regulations for Schools (1763; in Saxony 1773), they had, even if with some important modifications, conquered nearly all of Germany.

In the first half of the eighteenth century instruction in Biblical History received more attention also in the common schools. School regulations prescribed it for Dresden in 1711, for all electoral Saxony in 1713, for Prussia in 1716. In 1714 Huebner, the Hamburg rector, issued “\emph{Twice 52 Stories from the Old and New Testaments}”; while this book is not, as has been claimed, the “first Biblical History”, it has been a favorite text book till far into the nineteenth century on account of the skillful selection of stories and the questions appended to them. The Prussian Regulations mentioned above fixed Saturday as the day for instruction in Bible History. Also Valentine Loescher, an eminent champion of orthodoxy, advocated the introduction of Biblical History as a branch of religious instruction, and pleaded for a general regeneration of the education of the young. In the first quarter of the eighteenth century, probably owing to Francke’s influence, the first courses in catechetical theology were given in universities with Buddeus and Rambach as lecturers.

\quellseite{161}
\section[§ 18. The Period of Rationalism]{§ 18. Religious Instruction in the Church During the Period of Rationalism (About 1750–1830)}

Ph. Schuler, pp. 226 ff. G. v. Zeschwitz ii 2, 1 ch. 17.18. *Schumann-Sperber, pp. 61–87. J. Steinbeck, pp. 48–53. K. F. A. Kahnis, \emph{Der innere Gang des deutschen Protestantismus}, ³1874. *K. v. Raumer, \emph{Geschichte der Paedagogik} ii, ⁶1897. Chr. E. Luthardt, \emph{Der Rationalismus im Gebiet der Schule}, ³1891. *M. Schian, \emph{Die Sokratik im Zeitalter der Aufklaerung}, 1900. CHRIST. WOLFF AND HIS THEOLOGY: H. Stephan, \emph{Christ. Wolff} in PRE³, 1908. JOHANN LORENZ MOSHEIM: K. Heussi, \emph{J. L. Mosheim}, 1906. JOH. BERNH. BASEDOW: G. Hahn, \emph{Basedow und sein Verhaeltnis zu Rousseau}, 1885. O. H. Lang, \emph{Basedow, His Life and Educational Work}, 1891. H. Lorentz, \emph{Entwicklung und Bedeutung der Paedagogik Basedows} in Jahn-Fleckeisens Jahrbuecher, 1893. A. Pinloche, \emph{Geschichte des Philanthropinismus}, 1895. Balike, \emph{Die Stellung der Philanthropisten zum Religionsunterricht}, 1901. J. B. Basedows \emph{Paedagogische Schriften} ed. by H. Goering, 1880. J. J. ROUSSEAU: W. Boyd, \emph{The Educational Theory of Rousseau}, 1911. His \emph{Emile}, or treatise on education is translated into English by W. H. Payne, 1893. KARL FRIEDR. BAHR T: G. Frank, \emph{K. Fr. Bahrt} in Raumers Historisches Taschenbuch, 1866. L. Leyser, \emph{K. Fr. Bahrt, der Zeitgenosse Pestalozzis, sein Verhaeltnis zum Philanthropismus und zur neueren Paedagogik}, ²1870. Moses Mendelsohn, \emph{Phaedon oder ueber die Unsterblichkeit der Seele}, 1767; also reprinted in Reclams Universalbibliothek. CHR. G. SALZMANN: *E. Ackermann, \emph{Salzmanns ausgewaehlte Schriften}, 2 vols., 21901. Besides his \emph{Krebsbuechlein} (1780, reprint in Reclam), \emph{Ameisenbuechlein} (1806, repr. in Reclam), and \emph{Konrad Kiefer} (1796), compare esp. \emph{Ueber die wirksamsten Mittel, Kindern Religion beizubringen} (1780; new edition by P. Schuetze, 1905), \emph{Die Familie Ehrenfried, oder erster Unterricht in der Sittenlehre fuer Kinder von 8–10 Jahren} (1808), \emph{Heinrich Gottschalk in seiner Familie, oder erster Unterricht fuer Kinder von 10–12 Jahren} (1804), and \emph{Christliche Hauspostille} (1792 f.). JOHANN HEINRICH PESTALOZZI: L. W. Seyffarth, \emph{J. H. Pestalozzi nach seinem Leben und seinen Schriften dargestellt}, 1872, 61903. H. Kruesi, \emph{Pestalozzi, His Life, Work and Influence}, 1875. A. Pinloche, \emph{Pestalozzi and the Foundation of the Modern Elementary}
\quellseite{162}
\emph{School}, 1901. P. Natorp, J. H. \emph{Pestalozzi}, 3 vols., 1905. W. S. Monroe, \emph{History of the Pestalozzian Movement in the United States}, 1907. J. Natorp, \emph{Pestalozzi, sein Leben und seine Ideen}, 1908, ²1912. H. Holman, Pestalozzi, 1908. J. A. Green, \emph{Life and Work of Pestalozzi}, 1913. H. Debes, \emph{Das Christentum Pestalozzis}, 1880. L. W. Seyffarth, \emph{Saemtliche Werke Pestalozzis}, 12 vols., 1899–1902. *F. Mauer, \emph{Ausgewaehlte Werke Pestalozzis}, 4 vols., ⁵1897–1906. J. A. Green, \emph{Pestalozzi’s Educational Writings}, 1912. \emph{Abendstunde eines Einsiedlers} (1780), reprinted in Ostermann, \emph{Paedagogisches Lesebuch}, 1905. \emph{Lienhard und Gertrud} (1781–1787), reprinted in Reclam; translated into English and abridged by Eva Channing, 1896. \emph{Wie Gertrud ihre Kinder lehrt} (1801) reprinted in Reclam; English translation by L. E. Holland and Fr. C. Turner, 1898. A. Israel, \emph{Pestalozzi-Bibliographie}, 3 vols., 1903. GUST. FRIEDR. DINTER: Chr. Palmer, \emph{G. Fr. Dinter} in K. A. Schmid’s Enzyklopaedie, ²1878. W. Amelung, \emph{Dinters Grundsaetze der Erziehung und des Unterrichts}, 1881. Fr. Seidel, \emph{Dinters Schriften in Auswahl}, ²1887–1889. *A. Schultz, \emph{G. Fr. Dinter, sein Leben und seine Schriften}, 1908. FR. H. CHR. SCHWARZ: K. B. Hundeshagen, F. H. Chr. Schwarz in PRE,³ 1906. F. D. E. Schleiermacher: A. Heubaum, F. D. E. Schleiermacher in Rein’s Enzyklopaedisches Handbuch der Paedagogik, ²1902. G. v. Rohden, \emph{Schleiermachers Paedagogik}, 1884. *E. Platz, \emph{Schleiermacher’s paedagogische Schriften}, ³1902.—Compare H. Stephan, \emph{Neuzeit} (part 4 of Krueger’s \emph{Handbuch der Kirchengeschichte}), 1909. F. Uhlhorn, \emph{Geschichte der deutsch-lutherischen Kirche} ii, 1911. P. Monroe, \emph{Text-book in the History of Education}, 1905. *F. P. Graves, \emph{History of Education in Modern Times}, 1913. S. C. Parker, History of Modern Elementary Education, 1912. F. P. Graves, A. Student’s History of Education, 1915, ²1917.

By yielding ever increasingly to subjectivism, the disciples of Spener prepared the soil for the rise and spread of the Wolffian School (Christian Wolff, 1679–1754) and eventually for Rationalism. Wolff and his followers attempted to apply the methods of mathematical demonstration to theology, believing in all sincerity that they would thereby fortify the doctrine of the Church against attacks. Their method, however, led them inevitably to the point where they made reason the arbiter of revelation; thus they became the forerunners of Rationalism which recognized no other source of knowledge than reason and regarded as morally good whatever
\quellseite{163}
is useful. As far as religious instruction is concerned, the earlier exponents of Rationalism left the body of Christian doctrine intact, but they deemed an analyzing process sufficient for disclosing and setting forth its vital values (Loe-secke, Baumgarten); as a matter of fact, the children were by no means always taught to think, and very little permanent influence was exerted upon their will. Later Rationalists openly broke with the doctrine of the Church, demanded the abolition of Luther’s Catechism and substituted other text books for it. This turn of affairs was in no small degree due to the prestige which the so-called philanthropist Basedow (1723–1790) had gained; he aroused much interest in education, but thought that the sole aim of education was to train useful citizens of this world, and therefore advocated the divorce of school and church, abolition of confessional instruction, and the introduction of courses in morality. His principles were cordially received far and wide. That morality is well-pleasing to God only when based upon true religion, that real power for good can grow only from a heart which is devoted to God, had been forgotten. But even in those circles where orthodox doctrine was preserved, serious defects were apparent: Luther’s Catechism was still used, indeed, but it was not organically unfolded for the children; instead they were given independent dogmatical and ethical expositions (sometimes anything but congruous) which only too frequently expunged or at least concealed the very core of Christian truth, viz., faith in the crucified and risen Saviour as the one and only root of true morality. Questions which even today haunt some catechetical expositions, such as “How many religions are there?” “What is natural religion?” “What additional sources of religion are there besides Scripture?”, and the division of “duties” into “love of God, love of self, and love of the neighbor”, are remainders of the heritage from this period. The Hanoverian Catechism of
\quellseite{164}
1790 and Seiler’s Catechism of 1779, long used in Bavaria, are among the better products of this age. Salzmann in Schnepfental (Thuringia) and Dinter in Koenigsberg (1760–1831) are strong exponents of the rationalistic type of religious instruction. Dinter at least tried to come to terms with revealed truth, but the catechetical works of Parisius, Mayer, and others contain deistic philosophy and physiology rather than Christian saving truth.

Yet, in one respect this period marks an advance beyond all the catechetical labors of the past and is to be considered as the foundation of the better conditions obtaining in the present. From what has been said, this element of progress can not be found in the sphere of the material of instruction although we meet now and then with wise and commendable restriction to the essentials of Christianity, but rather in the \emph{didactic method}. The method was that of Socrates who like the obstetrician endeavored to draw forth into the light of day the germs of truth slumbering in the minds of his disciples, and has therefore been named “Socratic Method”. Already in his Ethics of Holy Scripture Mosheim, as early as 1735, had recommended this method for religious instruction of children; but he presupposed that the matter of instruction was first imbedded in the memory; it was then to be explained through the question method with the aim of bringing about firm convictions of faith and virtuous maxims of life. Apparently, however, Mosheim’s principles did not affect the catechetical practice of the decades immediately following. The Socratic method did not win favor until the appearance of Basedow’s \emph{Method of Instruction} (1764), of the \emph{Zuerich Questions} (1772), and the \emph{Philanthropinischer Erziehungsplan} (Philanthrophy as the Basis of Education) by the frivolous and licentious K. Fr. Bahrdt (1776). It is quite likely that Moses Mendelsohn’s \emph{Phaedon} (1767) is responsible for the name given to the method. The members of the Socratic
\quellseite{165}
school, at least its earlier exponents, actually attempted by skillful questioning to draw forth from the child’s mind all articles of faith and rules of life; in their opinion it was unnecessary to put anything into the child’s mind since everything, even if somewhat hazy and inchoate, is found in it by nature. However erroneous this conception, and however inadequate and defective this method, one important result was achieved: war was declared upon all senseless, mechanical memorizing. The proposition “Nothing is to be memorized unless it is understood” was and is sometimes exaggeratively applied, it is true; but if correctly employed, it helps develop an educational method which appeals to the intellect of the child, engages his co-operation, and by proceeding from the near to the far, from the known to the unknown, succeeds in arousing his interest, i.e., his mental presence and sympathy, and in moulding his whole future development. Also in this respect Dinter’s achievements are noteworthy, especially because he, more than all other rationalists who followed the Socratic method, proceeded from Biblical foundations instead of a supposititious innate body of knowledge possessed by the child. Even today one may learn much about the developing method of teaching from his \emph{Discourses upon the Chief Parts of Luther’s Catechism} (1806 ff.).

While the Socratic method ruled undisputedly till the end of the eighteenth century, a counter movement was ushered in at the beginning of the new. Its protagonist was \emph{Pestalozzi} (\emph{How Gertrude Teaches her Children}, 1801). He had much in common with the Socratics, but he endeavored to reach not only the head but the whole man, including the heart; and dissatisfied with an occasional, more or less accidental employment of intuition (\emph{Anschauung}) he made the principle of intuition fundamental for the whole process of education. The movement was furthered by Schwarz of Heidelberg through his important \emph{Catechetics} (1818): he
\quellseite{166}
made some concessions to the Socratic method, but demonstrated clearly that after all many religious truths simply cannot be elicited from the child’s mind and that they postulate another source. The Socratic method was still further discredited when Schleiermacher (1764–1834) notwithstanding his demand for the retention of the developing method, lifted up his mighty voice against it and maintained that the acroamatic method had to go hand in hand with the developing; the teaching of general truths he declared insufficient: the children are to become Evangelical Christians who will recognize, at the end of instruction, as their own the faith of the Church summarized in the Apostolic Creed. And when finally throughout the land men returned to the faith and theology of the Bible, the Socratic period so closely linked with Rationalism was at an end.

\quellseite{167}
\section[§ 19. Catechetical Labors Since the Renewal of Faith]{§ 19. Catechetical Labors Since the Renewal of Faith. Since About 1830}

*G. v. Zezschwitz ii 2, 1, ch. 34–50; ii 2, 2, ch. 19. Schumann-Sperber, pp. 83–122. J. Steinbeck, pp. 53–63. *G. Thomasius, \emph{Das Wiedererwachen des evangelischen Lebens in Bayern}, 1867. R. Rocholl, \emph{Geschichte der evangelischen Kirche}, 1897. G. Ecke, \emph{Die evangelischen Landeskirchen Deutschlands im 19. Jahrh.}, 2 vols., 1897. 1904. H. Stephan, \emph{Die Neuzeit} (part 4 of Krueger’s \emph{Handbuch der Kirchengeschichte}), 1909. F. Uhlhorn, \emph{Geschichte der deutsch-lutherischen Kirche}, ii, 1911. F. Schindler, \emph{Kritischer Wegweiser durch die Literatur des Konfirmandenunterrichts}, 1899. F. Cohrs, \emph{Katechismusunterricht} in PRE³, 1901. H. Scherer, \emph{Fuehrer durch die Stroemungen auf dem Gebiet der Paedagogik und ihrer Hilfswissenschaften}, part ii: Religions- und Moralunterricht, 1907. J. Berndt, \emph{Methodik des Unterrichts in der evangelischen Religion}, pp. 131–140, 1909. H. Lewin, \emph{Geschichte der Entwicklung der preussischen Volksschule}, 1910. *W. Ostermann, \emph{Geschichte der Paedagogik}, pp. 120–291, ⁴1910. E. Thraendorf, \emph{Allgemeine Methodik des Religionsunterrichts}, pp. 14–36, ⁵1912. WILHELM HARNISCH: R. Rissmann, \emph{W. Harnisch in seiner Bedeutung f. d. Entwicklung der deutschen Volksschulpaedagogik}, 1887. Metzmacher, \emph{Weiter- und Umbildung der Pestalozzischen Grundsaetze durch Harnisch}, 1901. Harnisch’s \emph{Volksschulwesen} has been republished by Bartels, 1893. F. A. W. DISTERWEG: G. Voigt, \emph{Diesterweg und die evangelische Volksschule}, 1891. *E. v. Sallwuerk, \emph{Leben, Lehre und ausgewaehlte Schriften Diesterwegs}, 3 vols., 1899–1900. JOHANN FR. HERBART: W. Kinkel, \emph{Joh. Fr. Herbart, sein Leben und Wirken}, 1903. *O. Fluegel, \emph{Herbarts Leben und Lehren}, 1907. G. Voigt, \emph{Die Bedeutung der Herbartschen Paedagogik fuer die Volksschule}, ⁴1908. *Chr. de Garmo, \emph{Herbart and Herbartianism}, 1895. Chr. Ufer, \emph{Introduction to the Paedagogy of Herbart}, ed. by Chr. de Garmo, 1894. W. J. Eckhoff, \emph{Herbart’s ABC of sense-perception}, 1896. *Sallwuerk u. Bartholomaei, \emph{Herbart’s paedagogische Schriften}, 2 vols., ⁷1903, ⁷1906. H. M. and E. Felkin, \emph{Herbart’s Science of Education}, translated, with a biographical introduction, 1895. TUISKON ZILLER, T. Ziller, \emph{Grundlegung zur Lehre vom erziehenden Unterricht}, 1865. *T. Ziller, \emph{Allgemeine Paedagogik}, ²1892. *T. Ziller, \emph{Einleitung in die allgemeine Paedagogik}, ²1901. G.
\quellseite{168}
Froehlich, \emph{Die wissenschaftliche Paedagogik Herbart-Ziller-Stoys}, 1901. E. v. Sallwuerk, \emph{Das Ende der Zillerschen Schule}, 1904. Compare the short but good chapter on Herbartianism in F. P. Graves, \emph{A Student’s History of Education}, pp. 333–369, 1915, and the excellent chapter on The Herbartians in S. Ch. Parker, \emph{History of Modern Elementary Education}, pp. 375–430, 1912. F. W. Doerpfeld: A. Carnap, \emph{F. W. Doerpfeld. Aus seinem Leben und Wirken}, ²1905. E. Hindrichs, F. W. Doerpfeld, ²1906. *F. Wienstein, F. W. Doerpfeld, \emph{Sein Leben und seine Schriften}, ²1907. Doerpfelds \emph{Gesammelte Schriften}, 12 vols., ¹894 ff.; the most important writings of Doerpfeld’s are \emph{Denken und Gedaechtnis}, ¹1906, \emph{Der didaktische Materialismus}, ⁵1915, \emph{Religioeses und Religionsunterrichtliches}, ¹895, \emph{Zur Methodik des Religionsunterrichts}, ⁵1906. Compare the pertinent articles in Rein, \emph{Enzyklopaedisches Handbuch der Paedagogik}, ²1902 ff.; also W. Rein, A. Pickel, u. E. Scheller, \emph{Theorie und Praxis des Volksschulunterrichts nach Herbarts Grundsaetzen}, 8 vols., ⁵1901–1908. History and Practice of Sunday Schools in Germany: H. Dalton, \emph{Geschichte, Wesen und Weise der evangelischen Sonntagsschule}, 1887. H. von der Goltz, \emph{Das Beduerfnis besonderer Jugendgottesdienste}, 1888. Reinhard, \emph{Zur Geschichte der Kindergottesdienste und Sonntagsschulen in Deutschland}, 1888. L. Tiesmeyer, \emph{Die Praxis der Sonntagsschule}, 1877. *F. Dibelius, \emph{Der Kindergottesdienst}, 1881. *Al. Schumann, \emph{Der Kindergottesdienst in seiner gesteigerten Bedeutung fuer Gegenwart und Zukunft der evangelischen Kirche}, 1909. \emph{Der Sonntagsschulfreund}, ed. by Fleischmann, 1869 ff. Der Kindergottesdienst, ed. by Tiesmeyer, Zaule ck, and Volkmann, 1891 ff.

Amid the wreck and ruin in church and school caused by Rationalism, the old catechisms preserved their hold upon the home and supplied the pious with food. A number of \emph{catechetical} works with Biblical content had appeared during the first quarter of the nineteenth century, and with the revival of confessional Lutheranism (about 1830) genuinely evangelical principles were once more quite generally applied to the work of training the young. Men like Nitzsch, Klaus Harms, and Hengstenberg made their influence felt in Northern Germany, and Brandt, Bomhard, Thomasius, Harless, and Loehe in Southern Germany. In Prussia Harnisch combatted Dinter’s influence (Hilfsbuch fuer Lehrer, Aid for Teachers, 1834); in Saxony Rudelbach advocated the restoration
\quellseite{169}
of \emph{Christenlehre} (instruction in evangelical doctrine on Sunday afternoons; 1840); in Bavaria Ackermann defended the use of the Chief Parts of the Reformers (1832) and Boeckh issued a good exposition of the Catechism. In 1845 Loehe published his excellent Catechism; one of the chief merits of this book is that it is free from dogmatical and historical encumbrances; however pressing the temptation may have been to include such material, Loehe refused to yield to it, but rather enunciated the following important principle: An exposition of the Catechism is no more than a presentation, according to the rules of exegesis, of the treasures of religious truth contained in Luther’s Enchiridion. Jaspis in Pomerania (1850) reminded again of the fact that no better illustrative material for catechetical instruction may be found than Biblical History. Mecklenburg-Strelitz introduced a splendid catechism in 1849. Between 1850 and 1860 the number of confessional Lutheran catechisms increased steadily. Special mention must be made of Caspari in Bavaria, Ernesti in Brunswick, Seebold in Hanover, Nielsen in Schleswig-Holstein. Also within the “Union” there was renewed activity. A good catechism was issued, e.g., for Baden in 1856 which served as the basis for the Rhenish Catechism. In Wuerttemberg the Brenz-Luther Catechism and the Booklet for Catechumens were revised. The School Regulations included catechetical instruction in the curriculum, and most of the old educational institutions in church and school were resuscitated and adapted to the new conditions. Expositions of the catechisms employed in the different regions were issued for the children; and bulky volumes for teachers and pastors were written by Wangemann, Arndt, Kaehler, Nissen, and others. Even catechetical sermons were preached once more; those of Loehe, Heubner, Ahlfeld, and Caspari are particularly excellent. The origin of Luther’s Catechism was investigated, the earliest attainable editions were republished
\quellseite{170}
(Harnack, Moenckeberg, Schneider); and collections of stories were issued as illustrative material (e.g., Caspari, \emph{Geistliches und Weltliches}). Of the expositions issued between 1860 and 1870 Buchrucker’s Explanation and his handbook \emph{Der Katechismusunterricht} in which he strongly pleads for the developing method, are the best. In the subsequent decades an almost endless variety of catechetical works appeared. The best books for children are those by Steinmetz and Beck, for the use of pastors those by Schuetze, Muenchmeyer, Schumann, Zezschwitz, and Steinmetz; Th. Kaftan’s exposition, however, deserves the palm for its discriminating and exquisite presentation of material. The most instructive product of the Ritschlian school is B. Doerries, Der Glaube. At present, certain theologians and pedagogues again demand with loud voices the elimination of Luther’s Small Catechism from the instruction of children. Compare § 35.

Another result of the renewal of faith was the increasing appreciation of \emph{Biblical History} as educational material. In their conflicts with the Rationalists the champions of the old faith stressed the fact that Christianity is based upon divine revelation, and quite naturally they came to evaluate the introduction into the history of revelation, into the genesis of Christianity, as an integral part of the education of the young. Especially Harnisch and F. L. Zahn advocated this branch. Harnisch emphatically declared that instruction in the Catechism can be profitably given only where Bible History has first been taught. In the common schools a number of stories were to be devotionally explained and recited; the teachers were given many valuable hints in Nissen’s Unterredungen ueber die biblische Geschichte (Talks upon Biblical History). Huebner’s stories of 1714 were widely used as text-book, also those issued by the Calw Publishing House. In the higher schools Biblical History was expanded into a history of the kingdom of God; here the books by Harnisch,
\quellseite{171}
Zahn, and later by Kurtz and Thomasius were often used. Franz Wiedemann’s How I Tell My Children the Stories of the Bible assisted mothers and other teachers of the little ones. In proportion as the views of Hofmann, the great Bible theologian of Erlangen, spread, it was deemed desirable even in the common schools, at least in the higher classes, to connect the individual stories of the Bible into a connected history of salvation. It was particularly Buchrucker who advocated this plan and composed excellent manuals for this purpose, a Biblical History for the pupils, and \emph{Biblischer Geschichtsunterricht} for the teacher. Where this seemed unattainable, the upper classes would at least study some outstanding Bible characters, as Wangemann had advised. At present Biblical History is so strongly emphasized in some quarters that all instruction in the Catechism is to be excluded until the pupils reach the higher classes. Of late “modern” theology and pedagogy begin to affect also instruction in Biblical History: the value of the Old Testament and in part the New Testament as records of the history of redemption is denied, yes even their historicity is frequently rejected; large parts, such as the account of creation or of the patriarchs, even the narratives of Jesus’ infancy and of his resurrection, are dismissed as legendary; there are many who expressly demand fairy tales as material of religious instruction in the lower grades and others desire the child to experience inwardly all the stages of religious “evolution” from naïve paganism to evangelical Christianity (Kulturstufentheorie). Compare § 34.

\emph{Instruction in the Bible} was also retained in this period, both in the form of Collections of Bible Passages, used rather exclusively as proof-texts for the Catechism, and in the form of Bible reading and of Bible Literature. A question that was much discussed was whether or not children should be given a complete Bible, or a School Bible or Biblical Reader
\quellseite{172}
which might be replaced by a complete Bible at the close of the school period. The \emph{Biblical Reader} by Voelker-Strack, and those for Bremen and for Wuerttemberg have found wide circulation.

Also the \emph{Hymnal} continues to be used in instruction. Well known hymn-tunes are practiced, and owing to the efforts of Lauritz and J. Zahn, in their original rhythmic setting; a treasure of choice hymns (the reduction of which to a smaller number is frequently demanded) is learned by heart and explained; the liturgy of the chief service is discussed, and often the pericopes are read and briefly interpreted. Also the most important events of \emph{Church History} are now usually treated in school, thus realizing one of Augustine’s ideals. Increasingly the need is felt of instruction concerning the polity and constitution of the Church, her tasks and duties in the present time, the work of missions, etc.; but specific forms for such instruction have not as yet been developed (cf. Warneck, \emph{Missions in the School}). The conviction is gaining ground that one cannot do justice to all these topics in the time previous to confirmation. They are therefore often studied in \emph{Christenlehre} where this ancient institution still exists, also in the Young People’s or Young Men’s Societies which owe their origin to the Church’s interest in the young after their confirmation. Others would remedy the situation by changing the confirmation practise or at least the time of confirmation so that confirmation will not any more coincide with graduation from grammar school. The supplementary care of the confirmed is still largely an un-solved problem which, like many other problems in the sphere of religious instruction, is highly complicated by the recent recrudescence of unbelief.

Since 1863 \emph{Sunday Schools}, introduced from America, have been built up in many cities of Germany. In most cases
\quellseite{173}
they have developed into religious \emph{services for children}, with or without division into groups.

Much was accomplished during the last century in the field of \emph{method}. The fundamentals of the Socratic Method were retained, but the attempt to elicit all truths from the child was made no longer; instead the material was first presented to the child, then made clear by the developing method, and thus utilized as an incentive for his will and food for his heart. During the whole century predominant influence was exerted by the great pedagogue Herbart (1776–1841; \emph{Umriss paedagogischer Vorlesungen}, 1835) who had been inspired by Pestalozzi. Owing to his and his followers’ efforts the peculiarities of the child’s soul life were given due — often one-sided — attention, as is seen from the \emph{Formalstufentheorie} (compare § 32). The progress made during the last century in the education of the German public school teachers rendered possible a fairly universal application of this method to religious instruction. But the theologically trained catechists rather avoided it; they tried to reach their aim of gripping and permanently influencing the whole soul by a combination of the acroamatic and erotematic methods and by adapting the educational material, as the catechism, and hymns, and esp. Bible History, to their pupils’ mental ability. Among the outstanding educators it was chiefly Doerpfeld who reduced the Herbartian theory to sane limits.

Of works embracing the whole sphere of Catechetics we mention the text-books by Kraussold (1843), Palmer (1844), Zezschwitz (1863–1872), Schuetze (1876), Kuebel (1876), Th. Harnack (1882), Buchrucker (1889), Sachsse (1897), v. Nathusius (1904), Gottschick (1908), and Steinbeck (1914), to which must be added the sections on Catechetics in the textbooks of Practical Theology, especially those of Krauss, Knoke, and Achelis. O. Baumgarten’s Neue Bahnen (1903) is suggestive and instructive; like Kabisch’s \emph{Wie lehrt man}
\quellseite{174}
\emph{Religion?} It discusses the subject from the standpoint of modern theology. The works on Didactics or Public School Methods by Doerpfeld, Kehr, Schumann, Staude, Thraendorf, Rein, Reukauf, and others contain much valuable material on methods of religious instruction. For the history of catechetical methods compare also § 34–36.

\quellseite{175}
\section{§ 20. Catechetical Labors in America}

G. Weber, \emph{Geschichte der akatholischen Kirchen und Sekten in Grossbritannien}, 1845f. \emph{Original Letters relative to the English Reformation}, edited by the Parker Society, vol. 2, 1848. \emph{Three Primers put forth in the Reign of Henry VIII.}, 1848. \emph{First Prayer Book of Edward VI compared with successive revisions}, 1883. A. F. Mitchell, \emph{Catechisms of the Second Reformation}, 1886. H. Bonar, \emph{Catechisms of the Scotish Reformation}, 1886. *B. M. Schmucker, \emph{Luther’s Small Catechism, The Editions, Translations, and Explanations Published in America} in Lutheran Church Review 1880. *H. E. Jacobs, \emph{The Lutheran Movement in England}, 1890. P. L. Ford, \emph{The New England Primer}, reprint of the earliest known edition, 1899. E. F. K. Mueller, \emph{Bekenntnisschriften der reformierten Kirche}, 1903. M. Reu, \emph{Quellen zur Geschichte des kirchlichen Unterrichts}, 1 and 2, 1904, 1911. A. Lang, \emph{Der Heidelberger Katechismus und verwandte Katechismen}, 1907. F. Watson, \emph{English Grammar Schools before 1600}, 1908. W. H. Free and W. H. Kennedy, \emph{Visitation Articles and Injunctions of the Period of the Reformation}, 1910. John W. Prince, \emph{Wesley on Religious Education}, 1926. L. G. Pray, \emph{History of Sunday Schools and of Religious Education}, 1847. I. C. Power, \emph{The Rise and Progress of Sunday Schools, a Biography of Rob. Raikes and William Fox}, 1863. S. Gilbert, \emph{The Lesson System: the Story of its Origin and Inauguration}, 1879. H. C. Trumbull, \emph{Yale Lectures on the S. S.}, 1888. J. H. Harris, \emph{Rob. Raikes, The Man and his Work}, 1899. M. C. Brown, \emph{Sunday School Movement in America}, 1901. S. B. Haslett, \emph{Pedagogical Bible School}, 1903. O. G. Michael, \emph{Sunday School in the Development of the American Church}, 1904. F. W. Rice, \emph{Origin and Expansion of the S. S.}, 1906. *H. F. Cope, \emph{The Modern S. S. in Principle and Practice}, 1909. *H. H. Meyer, \emph{The Graded S. S. in Principle and Practice}, 1910. H. F. Cope, \emph{Evolution of the S. S.}, 1911. J. R. Sampey, \emph{International Lesson System}, 1911. J. R. Hurlbut, \emph{Organizing and Building up the S. S.}, 1911. R. P. Shepherd, \emph{Religious Pedagogy in the Modern S. S.}, 1911. Proceedings of the Religious Education Association 1903 ff. \emph{The Encyclopedia of Sunday Schools and Religious Education}, 3 vols., 1915. De Garmo, \emph{Principles of Religious Education}, 1901. *G. U. Wenner, \emph{Religious Education and the Public School}, 1907. A. N. Riley et al., \emph{The Religious Question in}
\quellseite{176}
\emph{Public Education}, 1911. H. Chapell, \emph{The Church Vacation School}, 1915. *W. S. Athearn, \emph{Character Building in a Democracy}, 1924. *F. O. Gift, \emph{Week Day Religious Education}, 1926. State Teachers College of Colorado, \emph{Special Bulletin of Religious and Moral Education}, 1913. *S. C. Parker, \emph{History of Modern Elementary Education}, 1913. W. P. Burris, \emph{The Public School System of Gary, Ind.}, U. S. Bureau of Education Bulletin 591, 1914. W. J. Mutch, \emph{Religious Day School} in Encyclopedia of S. S. etc., 1915. R. S. Bovine, \emph{Daily Vacation Bible School Association}, ibid., 1915. *H. H. Meyer, \emph{Cooperation in Christian Education}, 1917. S. A. John, \emph{Die Sonntagsschule, ihre Bestimmung u. Arbeit}, 1897. C. W. Hertzler, \emph{Die religioes-sittliche Erziehung d. kirchl. Jugend}, 1913. F. W. Schneider, \emph{Barclays Handbuch zur Heranbildung von Sonntagsschullehrern}, 1916. H. E. Jacobs, \emph{History of the Ev.-Luth. Church in the U. S.}, 1893. A. Graebner, \emph{Geschichte d. luth. Kirche Amerikas}, 1892. G. I. Fritschel, \emph{Geschichte der luth. Kirche Amerikas}, 1896f. J. L. Neve, \emph{Kurzgefasste \emph{Geschichte der luth. Kirche Amerikas}}, 1915 (here further literature). J. C. W. Lindemann, \emph{Schulpraxis}. J. Schaller, \emph{Notwendigkeit d. christl. Gemeindeschule fuer die christl. Familie, Kirche und Staat}. G. M. Grossmann, \emph{Die christl. Gemeindeschule}, 1895. J. R. E. Hunt, \emph{The Luth. S. S. Handbook}, 1911. R. Bunge, \emph{Weide meine Laemmer}, 1912. E. H. Engelbrecht, \emph{Why Lutheran Parish Schools}, 1915. G. H. Gerberding, \emph{The Lutheran Catechist}, 1910. J. H. Hertzer, \emph{Evang.-luth. \emph{Katechetik}}, 1911. M. Reu, \emph{Katechetik}, 1915. *\emph{Symposium on the Sunday School} in Lutheran Church Review, 1896. Theo. Schmauk, \emph{The General Council Sunday School System of Grading and Instruction}, 1897 ff. M. Reu, \emph{Wesen u. Aufgabe der Sonntagsschule} in Kirchl. Zeitschrift, 1900. 1901. M. Reu, \emph{Grundsaetze zur Herstellung von Sonntagsschulliteratur}, ibid., 1911. M. Reu, \emph{Wartburg Lehrmittel fuer ev.-luth. Sonntagsschulen}, 1911 ff. M. Reu, \emph{Wartburg Lesson Helps for Luth. Sunday Schools}, 1914 ff. Chr. P. Wiles, \emph{The Challenge of the S. S.}, 1916. *M. Reu, \emph{A New English Translation of Luther’s Small Catechism}, 1926.

The same inconclusiveness which characterizes the religious history of England between the years 1530 and 1550 confronts us also when we investigate the status of religious instruction in that period. Marshall’s \emph{Goodly Primer} of 1534 was uncompromisingly evangelical. Also the second edition (1535) which contains unabbreviated Luther’s \emph{Short Form} of 1520, is permeated by the \emph{sola gratia}; it is true that in
\quellseite{177}
the Litany which bears the superscription “\emph{\emph{The Invocation of the Saints}}”, a concession is made to the weak, but in the preface (pp. 123 ff.) it is explicitly stated that it is a concession, no more. Romanizing tendencies are not in evidence; the enumeration of the prohibition of image worship as a specific commandment (Dialogue, pp. 216-221) rather indicates Swiss influence. Henry VIII’s injunction of 1536, however, in no way transcends those of medieval times since the necessity of explaining the parts is ignored; the injunction reads as follows: “The curates shall, in their sermons, deliberately and plainly recite of the said \emph{Pater Noster}, the Articles of our Faith, and the Ten Commandments, one clause or article one day, and another another day, till the whole be taught and learned by little; and shall deliver the same in writing, or show where printed books containing the same be to be sold to them that can read or will desire the same.” In 1537 Archbishop Cranmer and a commission of bishops edited an authorized explanation of the Lord’s Prayer, the Ave Maria, the Creed, the Decalogue, and the seven Sacraments under the title \emph{The Godly and Pious Instructions of a Christian Man} (commonly called Bishop’s Book), but the explanation was Catholic, and not Evangelical. Of the same type is \emph{The Necessary Doctrine and Condition of a Christian Man} (also called King’s Book), an explanation issued in 1543. When Henry VIII died, in January 1547, and Edward VI acceded to the throne, more resolute measures of reformation were taken. In 1548 Cranmer issued an English translation of Justus Jonas’ Latin version of the Nuremberg Sermons for Children (Osiander whose niece had married Cranmer in 1532, was their chief author) under the title \emph{A Short Instruction into the Christian Religion; for the Singular Commoditie and Profite of Children and Young People}. Thus Luther’s Catechism and its best extant exposition was made accessible to the English people no less than to most
\quellseite{178}
of the other European peoples (see Reu, Quellen, etc.). It was not Luther’s Catechism, however, that gained vogue in subsequent years, but the \emph{Instruction to be Learned of Every Child}, a catechism which had been issued as part of the Order of Confirmation appearing in the first edition of \emph{the Book of Common Prayer} (1549), and which had to be recited from memory at the time of confirmation. It begins with a statement concerning sponsorship and then succinctly explains Creed, Decalogue, and Lord’s Prayer. After 1552, owing to Reformed influence (Butzer had been brought into the country, Calvin had become the chief adviser of Edward and Cranmer; but cf. also Marshall’s Primer), the commandments which hitherto had been compressed into short sentences, were given in the exact wording of Exod. 20. After Bishop Overall (1604) had added the section pertaining to the Sacraments, this catechism remained unaltered, and it is in force in the Anglican (Episcopal) Church to this day; it is, however, no longer printed as part of the Confirmation service, but as an independent section of \emph{the Book of Common Prayer} (since 1662). Authorized by the king, Poinet issued a larger catechism for adults and teachers in 1553, but it never gained wide circulation. Greater influence was exerted by the catechism issued by Nowell with the approval of archbishops and bishops, but it likewise failed to become an official text book for instruction; it is noteworthy, however, that Nowell’s catechism furnished the basis for the above-mentioned section on the Sacraments added to the Catechism in 1604. For many years, if we may judge by the curricula extant, Nowell’s catechism in its original Latin form as well as in English and Greek translation was used for instruction in religion and in the ancient languages.

In Scotland it was Knox who emphatically advocated religious instruction of the young. In his Book of Discipline (1560) we find the following rule for all places not large
\quellseite{179}
enough to support day schools: “Either the reader or the minister must take care of the children and youth of the parish, instructing them in their first rudiments, and especially in the catechism, as we have it now translated in the Book of our Common Order, called the \emph{\emph{Order of Geneva}}.” He also enjoined that on Sundays “after noon the young children must be publicly examined in their catechism in audience of the people, and in doing this the minister must take great diligence, to cause the people to understand the questions propounded, as well as the answers, and the doctrine that may be collected thereof.” Nor did he forget home instruction; he wrote: “Every master of household must be commanded either to instruct, or else cause to be instructed his children, servants and family, in the principles of the Christian religion.” The catechism to which he referred is Calvin’s Catechism of 1545 (cf. p. 122). The Presbyterian Church of England, subsequently also the Church of Scotland introduced the Shorter Westminster Catechism sanctioned by parliament in 1648 which is still in use. While the Large Westminster Catechism (also approved by parliament in 1648) resembles a textbook of dogmatics rather than a catechism (it is based on the \emph{\emph{Compendium theologiae}} of the Basle theologian Wolleb, issued in 1628 and contains 196 questions and answers in part very lengthy), the Shorter Westminster Catechism has justly been praised for its clearness and careful wording and its commendable restriction to essentials, but with its 107 questions it is, like the Heidelberg Catechism, still too voluminous and doctrinal. It contains a large number of proof-texts. The title of the authorized edition is: “The Shorter Catechism, agreed upon by the Assembly of Divines at Westminster, with the assistance of commissioners from the Church of Scotland, as a part of the covenanted uniformity in religion betwixt the churches of Christ in the kingdom of Scotland, England and Ireland, and approved
\quellseite{180}
anno 1648 by the General Assembly of the Church of Scotland, to be a directory for catechizing such as are of weaker capacity, with the proofs from the Scripture.” It is divided into two parts, (1) What man is to believe concerning God, and (2) What duty God requires of man. The first part is joined rather loosely to the Creed, the second part discusses the Ten Commandments, faith and repentance, the Sacraments, and the Lord’s Prayer; at the end of the book the words of the Creed, the Decalogue, and the Lord’s Prayer are given. — The Westminster Catechism was used for a long time also by the Congregationalists. — In Holland and in the Palatinate the Reformed Church used the Heidelberg Catechism which at Dort (1618) had been made the confession of all Reformed churches.

From these premises conclusions are readily drawn regarding the catechetical helps and the manner of religious instruction obtaining in America among the immigrants coming from England, Holland, and the Palatinate. The families that had been members of the Episcopal Church brought with them the Book of Common Prayer and instructed their children prior to confirmation in the Catechism of 1549 which was contained in it. Presbyterians coming from Scotland, Ireland, or England, also the Congregationalists brought the Shorter Westminster Catechism. The Dutch Reformed and the Palatines brought the Heidelberg Catechism. It is likely that the largest part of instruction was given in the home, but in accordance with the usage in the mother countries, also Sunday was devoted to religious instruction. We hear, for instance, that at Roxbury, Mass., the young were instructed on Sunday as early as 1674, at Norwich, Conn., in 1676, and in the Pilgrim’s Church, Plymouth, Mass., in 1680. At Plymouth it was resolved “that the deacons of the church be requested to assist the minister in teaching the children during the intermission on the Sabbath” (i.e., between the
\quellseite{181}
forenoon and afternoon services). The Rev. M. Jones opened a Sunday School at Newton, L. I., N. Y., in 1683. The Schwenkfelders in Berks and Montgomery counties, Pa., had a similar institution in 1734. It is claimed that John Wesley during his stay in America (1735-1738) established a Sunday School at Savannah, Ga. The Church of the Brethren at Germantown, Pa., gave its youth instruction on Sundays as early as 1738; it is said that Christopher Saur printed tickets or cards for use in this school. In 1740 Rev. Joseph Bellamy opened a Sunday School at Bethlehem, Conn., which exists to this day; his Memoir states that “he is believed to have been the first pastor in the land, if not in the world, who began and through all his ministry kept up, a Sabbath school in his congregation, regularly spending an hour in the interval of public worship, on the Sabbath, in catechizing and instructing one class of children and another of adults, in the word of God”; this is indeed an overstatement, but it furnishes proof that at that time such instruction on Sundays was not yet widely prevalent. Ludwig Hoecker, a Seventh Day Baptist or Dunker, established a Sunday School at Ephrata, Lancaster county, Pa., in 1740 in order “to give instruction to indigent children as well as to give religious instruction to those of better circumstances”. In 1744 Mrs. Greening opened a school in Philadelphia. Among the records of the early Presbyterian church at Westerly, R. I., there is a note dated 1752 according to which it was required “statedly to hear the children read a portion of ye Holy scriptures, and repeat ye Assembly’s Catechism”.

Besides the Catechisms enumerated above, of which the Shorter Westminster Catechism was specifically mentioned in the last quoted note, and the Bible, the catechetical helps most frequently used in home and school were the writings of Cotton, Harris, and Watts. John Cotton arrived in America in 1632 and published a catechism in 1646 which bore the
\quellseite{182}
title “\emph{Milk for Babes, Drawn out of the Breast of both Testaments, chiefly for the spiritual nourishment of Boston babes in either England, but may be of like use for any children, a Catechism}.” This little writing was reprinted even in England (Cambridge 1656), and it went through very many editions in New England. The author’s grand-son called it “peculiarly the catechism of New England” and said: “It will be valued and studied and improved, until New England ceases to be New England.” It does not surprise, therefore, that we find it incorporated in the \emph{\emph{New England Primer}}, the book which for over a century was the most important school book of New England. The \emph{\emph{New England Primer}} was composed probably by Benjamin Harris, and very likely it was printed to take the place of his \emph{The Protestant Tutor}. The first edition must have appeared between 1687 and 1690 as the “second Impression of The \emph{\emph{New England Primer}}, Enlarged” was advertised in 1691. The contents of the book vary in different editions, but generally it contained the simple alphabet, also rhymed and illustrated alphabets, the vowels, consonants, double letters, italic and capital letters, “easy syllables for children,” forming a syllabary of words from one to six syllables in length, the Lord’s Prayer, the Apostles’ Creed, pictures of birds, animals, and fish, each with a rhyme, “lessons for youth,” which are admonitory Scripture verses, Dr. Watts’ cradle hymn, “verses for little children” and prayers and advice for them, the Westminster Assembly’s Shorter Catechism, John Cotton’s “Spiritual Milk for Babes,” and a “Dialogue between Christ, Youth, and the Devil.” Some editions included proper names of men and women, the Ten Commandments, and the names of the books of the Bible given in their order. This primer reminds one strongly of school editions of Luther’s Catechism from the period when it was the sole school book in Germany.

Among the contributors to the \emph{\emph{New England Primer}}
\quellseite{183}
we meet also the third author who powerfully influenced school and home through his writings, the well known hymn writer Isaac Watts of whom Samuel Johnson said: “For children, he condescended to lay aside the scholar, the philosopher, and the wit, and to write little poems of devotion, and systems of instruction adapted to their minds and capacities, from the dawn of years through its gradations of advance in its morning of life.” His juvenile hymns published 1715 with the title \emph{Divine Songs} and republished 1720 with additions under the title \emph{Divine and Moral Songs for the Use of Children} had profound religious effects. Owing to their child-like simplicity and their spiritual warmth they enjoyed phenomenal circulation. And not only in England but also in America they constituted a most important factor in religious education of the young. More than 100 editions had been issued by 1750; Wilmer (\emph{Life of Watts}, p. 372) states that at the end of the century about 80,000 to 100,000 copies were printed annually; and Julian (\emph{Dictionary of Hymnology}, first edition, 1892) calls this hymn-book and two other writings of Watts which will be mentioned directly “the most popular text-books for religious education 50 years ago.” The two other important works of Watts are his \emph{Plain and Easy Catechisms} published in 1730 and his \emph{Scripture History} of 1732. The “First Catechism” presents in brief form the Christian way of salvation and a “Catechism of Scriptural names” (Who was Adam? Who was Eve? etc.); the “Second Catechism” contains an exposition of the Ten Commandments and of the Sacraments and some prayers, among them the Lord’s Prayer. His \emph{Scripture History} undoubtedly is the first Bible History in English; his collection of hymns is indeed not the first juvenile hymnal — Spangenberg had issued one as early as 1544, and several other hymnals for children had appeared after that, — but the first book of the kind to be published in the English language.

\quellseite{184}
New impulses for the religious education of the young proceeded from the \emph{Sunday School Movement} called into being by Robert Raikes of Gloucester (1735 or 1736-1811) and soon afterward transplanted to America.

\begin{quote}
E. W. Rice, in \emph{Schaff-Herzog Encyclopedia} vol. xi, summarizes the results of Harris’ investigations relative to Raikes’ life and work as follows: “Robert Raikes” (Sept. 14, 1735 is the date generally given for his birth; in the baptismal records of the Church of St. Mary de Crypt at Gloucester, however, the following entry is made: “Sept. 24, 1836, Robert, son of Robert and Mary Raikes, of this parish”; is it possible that he was a year old before he was baptized?) “was editor and proprietor of the \emph{Gloucester Journal}, Gloucester, England, and a strange mixture of the “dandy” and of the Reformer. Before he was of age, he began to visit the two prisons of Gloucester, to relieve the horrors of prison life, and to reform the prisoners. His sympathies were widened, his charities deepened; his failures in prison reform set him thinking, until he reached the conclusion that ‘vice is preventable.’ Twenty-five years later, when he was forty-four years of age, he began a ‘new experiment’ as he called it, of ‘botanizing in human nature.’ Going into the suburbs of the city, where many youths were employed in the factories, his heart was touched by the groups of ragged, wretched, cursing children. He knew their parents, homes, and habits; none ever entered the House of God. It was useless to appeal to such parents. He had tried to reform adults and had failed. George Whitefield had tried to reach the masses in Gloucester, but with meager results. Raikes was moved, therefore, to apply his maxim that ‘vice is preventable.’ ‘Begin with the child, for idleness is the parent of vice,’ and ‘ignorance is the cause of idleness’; therefore, ‘begin by instructing the child.’ These seem like trite statements now, but they were the result of long, deep thought by Raikes. Here was his mission. But the Rev. Thomas Stock, head-master of the Cathedral school, whom he met in his walk, was the man to help him. He started his first Sunday School in Sooty Alley in 1780, paying Mrs. Meredith for teaching the wretched little street children whom he persuaded to come to her kitchen for instruction. Mrs. Meredith found the boys ‘terribly bad,’ and soon the pupils were transferred to Mrs. Mary Critchley’s home in Southgate Street, whose house extended to Grey Friars. Bad as the boys were, the ‘girls were worse.’ The children were required to come with clean hands and faces, hair combed, and
\quellseite{185}
with such clothing as they had, though shoes and clothes were sometimes provided. The boys were ‘strapped’ or ‘caned’ by Raikes himself, for misbehavior; the girls were subdued by other means. The children were to remain in school from ten to twelve, then go home; to return at one, and, after a lesson, to be conducted to church; after church service to repeat portions of the catechism; then to go quietly home about five, without playing in the streets. Attentive scholars received rewards of Bibles, Testaments, books, combs, shoes, and clothing. The head teachers were paid a shilling a day. Raikes engaged four women in his schools, and procured other employment for them as rewards of diligence which ‘may make it worth sixpence more.’ The Rev. Thomas Stock ‘went around to the schools Sunday afternoon,’ says Raikes, ‘to examine the progress made, and to enforce order and decorum among such a set of little heathen.’ The boys were in classes of five, the advanced pupils acting as ‘monitors’ or teachers, teaching the younger pupils their letters. The girls in a separate room, with white tippets on their shoulders and white caps on their heads, were in classes also, with ‘monitors’ or sub-teachers over them. The monitors and sub-teachers were unpaid and voluntary, selected and directed by the paid master or mistress. For about three years Raikes looked upon his schools as an experiment. When William Fox, William Wilberforce and the Wesleys—John and Charles—and Jonas Hanway came as the guests of Raikes and his neighbors, he explained his plan, asked their counsel, and took them to the school to hear the children repeat prayers, the catechism, answer Bible questions, and sing Watts’ hymns. It is recorded that they were astonished, ‘caught the fire,’ and extended the movement. When satisfied that his scheme had passed the experimental stage, Raikes published a brief notice of it in the \emph{Gloucester Journal} of Nov. 3, 1783, which was copied into the London papers. \emph{The Gentleman’s Magazine} also published a letter of Raikes of November 25, 1783, in full, and a little later, another description by Raikes was given in the \emph{Arminian Magazine} edited by John Wesley. These and many other published accounts extended knowledge concerning the new movement, while many pulpits repeated the story and praised the institution.”
\end{quote}

Raikes had no intention to build up schools for all children, but he wanted to lift the neglected, wretched, unrestrained children of the city out of their misery; and although he was not chiefly concerned with religious instruction
\quellseite{186}
, but their mental and moral training, he gladly employed such aids as Bible reading, Catechism study, singing of hymns, and attendance of divine services. Moreover, it is very probable that in the course of time other children beside those whom he had in view originally, began to attend his schools because at that time England had no state-supported free public schools (until 1833 elementary education in England was left to the home; the Church, or private schools charged tuition). Certain elements in the Church at first violently resisted the movement: the archbishop of Canterbury called a convocation to determine methods of stopping the movement; William Penn seriously considered introducing a bill in parliament “for the suppression of Sunday Schools”; and in Scotland laymen who taught in Sunday Schools were held guilty of “Sabbath breaking.” But the movement was not to be retarded: in a surprisingly brief span of time the thoughts of Raikes had swept over the country, opposition was brushed aside, and schools began to spring up in large numbers. Bolton had a Sunday School with 80 volunteer teachers as early as 1787, the Sunday School at Stockport had 30 teachers in 1794, and in 1800 there were 156,400 Sunday School pupils.

William Fox, a Baptist, vigorously promoted the Sunday School cause and called a representative meeting in London on Sept. 7, 1785; its chairman was the philanthropist Jonas Hanway, and it was here that “\emph{The Sunday School Society}” was founded “for the support and encouragement of Sunday Schools in the different counties of England.” In 1803 “\emph{The British Sunday School Union}” was organized in London. Though founded originally by a group of Sunday School teachers for mutual assistance and encouragement, it soon began to consider as its chief purpose the establishment of new schools. Almost from its beginning this society has been a publishing agency, issuing lesson plans, lesson helps, etc.;
\quellseite{187}
it also arranged lecture courses. In 1818 the official figures for England were 5463 Sunday Schools with 477,225 pupils. Various factors contributed to this extraordinary development; H. F. Cope mentions the following four: (1) The awakening interest in the general education of the young which arose in various places; (2) a wide-spread development of humanitarian sentiment which led to the organization of many important relief and betterment societies; (3) the remarkable religious revival which is today best remembered in the work of Whitefield and Wesley; (4) the industrial revolution, together with upheavals of the French revolution and American independence.” The Sunday Schools of England have continued on the general pattern of the schools of Raikes; Cope enumerates the following characteristics: (1) Organized and conducted independently of church control; (2) without denominational oversight and promotion; (3) designed to combine elementary general education with religious instruction; (4) lacking the urgency of a secular system of general education.

The waves of the new movement soon touched America. Yes, here the Sunday School developed much more vigorously and rapidly than in England, and from here powerful impulses have been communicated to other countries. The stupendous growth was caused, (a) on the one hand by the fact that in the United States the government at first paid very little attention to the training of the young, leaving elementary education to the family and the Church. It was estimated that in Indiana less than one-sixth of the children attended any school in the first quarter of the nineteenth century. In the larger cities the Lancasterian system of monitors was applied; in Philadelphia, e.g., in 1834 there averaged one teacher to 218 pupils, the work being directed by the teacher with the aid of “monitors” who were responsible for certain groups of pupils. Such herding of pupils in large numbers was, of course, extremely ineffectual.
\quellseite{188}
The government of Massachusetts required towns to maintain schools by the eminently important law of 1647. This law reads, in part, as follows: “It being one chief point of that old deluder, Satan, to keep men from the knowledge of the Scriptures, as in former times, by keeping them in an unknown tongue, so in these latter times, by persuading from the use of tongues, that so at last the true sense and meaning of the original might be clouded by false glosses of saint-seeming deceivers, that learning might not be buried in the grave of our fathers in church and commonwealth, the Lord assisting our endeavors—It is therefore ordered that every township in this jurisdiction, after the Lord has increased them to the number of fifty householders, shall then forthwith appoint one within their town to teach all such children as shall resort to him to write and read, whose wages shall be paid either by the parents or masters of such children, or by the inhabitants in general, by way of supply, as the major part of those that order the prudentials of the town shall appoint; providing those that send their children be not oppressed by paying much more than they can have them taught for in other towns.” Parker shows that many of the Massachusetts towns failed to obey the spirit of the law of 1647. The reading texts used were religious books like the \emph{New England Primer} and similar primers, until spelling books, above others \emph{Webster’s Speller} of 1783, took their place.—“In colonial Pennsylvania, in contrast with Massachusetts,” Parker says (pp. 62 f.), “no general system of public schools was developed, but elementary education remained entirely in the hands of the churches and neighborhood organizations which were actuated by religious motives. Wm. Penn had contemplated the organization of a system of public schools, but his Utopian ideals were not realized. The second general assembly of the colony (1683) passed a law requiring that all children be taught so that they could read the Scriptures and write by twelve years of age. The law soon became a dead letter, however, owing to changes in the government and the conflicting interests of such a cosmopolitan colony. The assurance of liberty of religious worship attracted in Pennsylvania a great many Protestant religious emigrants and exiles. These included Quakers, Baptists, Presbyterians, Methodists, German Lutherans, members of the German Reformed Church, Moravians, and others. They were generally believers in the Protestant principle of training to read the Bible as the road to salvation, and each group of religious enthusiasts set up a school as an essential part of its religious organization. These church schools predominated in the
\quellseite{189}
eastern portions of the state. In the more thinly settled and frontier parts of the state, where these compact religious communities were not found, the more mixed communities tended to establish subscription (voluntary), or ‘neighborhood’ schools. These were generally the result of the voluntary cooperation of a few families, often stimulated by some energetic and wealthy father who desired that his children should have at least an elementary education. The neighborhood schools were most common in the western part of the state. Together with the Church schools they provided nearly all the elementary education available down to 1834.”

On the other hand, (b) the success of the American Sunday School was due to the fact that the Church fostered the new institution, regarding it from the very beginning as the real church school and subsequently as the one and only church school. The earliest American Sunday Schools were quite similar to those founded by Raikes; the “First Day or Sunday School Society” founded in Philadelphia in 1791 by Bishop White (Episcopalian), Mathew Carey (Roman Catholic), and Benjamin Rush (Universalist) was organized in order to maintain schools for those “who on Sundays were employed in the worst of purposes, the depravity of morals and manners,” but soon the sphere of activity was widened so as to embrace all children. Again, wherever no elementary public school existed or where the Church had not established any day school, the Sunday School also gave some elementary instruction, but this feature disappeared when the number of public state-supported schools began to increase, and the Sunday School became a school exclusively religious. And when the ever improving secular school had crowded out of existence the religious day school the Sunday School became the one and only school of the Church because religious instruction was excluded from the public school. Small wonder that many hands were busily engaged in improving, extending, and systematizing the work of the Sunday School.

It appears that the Methodist Church was the first to cultivate the Sunday School in its newer form in the United
\quellseite{190}
States. More than once John Wesley had written in praise of Raikes’ schools (for instance, under date of July 18, 1784: “I find these schools springing up wherever I go. Perhaps God may have a deeper end therein than men are aware of. Who knows but that some of these schools may become nurseries for Christians”). Every Methodist pastor, in the first \emph{Book of Discipline of Methodist Episcopal Churches} (1784), was directed “where there are ten children whose parents are in the society, to meet them at least one hour every week.” In 1786 Francis Asbury, the first bishop of the Methodist Church to be ordained in the United States, organized a Sunday School in the home of Thomas Crenshaw in Hanover county, Va. In 1790 the Methodist Conference at Charleston, S. C., ordered that these schools should be established “in or near the place of worship,” provided for the appointment of teachers, and set the sessions “from six in the morning till ten and from two o’clock in the afternoon till six, when it does not interfere with public worship.” The next year (1791) a Universalist Sunday School was established in Philadelphia, a Friends’ school in the same city in 1792, and a Baptist school at Pawtucket, R. I., in 1797. At an early time organizations were formed which had as their purpose the establishment of new Sunday Schools. On basis of Brown’s investigations, Cope furnishes the following summary: “The organization of the “First-Day or Sunday School Society” of Philadelphia, in 1791, was followed in 1808 by the organization of the “Evangelical Society” to promote Sunday evening schools in Philadelphia; the Pennsylvania Union in Pittsburgh, 1809; the Female Union Society in New York, 1816; the New York Sunday School Union, 1816; the “Boston Society for the Moral and Religious Instruction of the Poor,” 1816. All these were of local influence only, but a wider scope was designed for the “Philadelphia Sunday and Adult School Union,” 1817, for it planned to promote
\quellseite{191}
the organization of new schools in villages and the country. In 1821 it employed what was probably the first American Sunday School missionary; he organized sixty schools in six states. Doubtless this work prepared the way for the American Sunday School Union.” It is quite evident that the American Sunday School Union has adopted the chief principles of the Philadelphia Union and continues to apply them to this day. One of the most important of these was the idea that the Sunday School should be interdenominational. Thus we read in the Statutes of the Philadelphia Sunday and Adult School Union that one of its objects was “to cultivate unity and Christian charity among those of different names,” and in presenting their appeal to the public for generous support its leaders said: “The comparative fewness of Christians calls for all practicable and profitable union among themselves. Divide and conquer is the maxim of their great foe. Unite and triumph be then the motto of Christians.” Tracts in large numbers were issued to promote the formation of Sunday Schools and S. S. Associations, for instance a “\emph{Model Constitution},” a “\emph{System for the Internal Regulation of Sunday Schools}” including the “grading of the school in four grades, providing for two sessions of the school each Sunday, and thorough rules with a complete scheme of rewards and penalties for pupils, teachers, and officers, based upon record of attendance, lessons, and behavior.” In 1827 the Union issued 25,000 Sunday School Hymn books, 8000 class books, 2000 \emph{Teacher’s Guides} by John Angell James, 81,000 premium books, and for the third time, half a million of red and blue tickets.

The ever expanding work of this Philadelphia Union led to the formation, in 1824, of the American Sunday School Union, an organization that was national in name and scope. The objects of this Society were stated by the founders thus: “To concentrate the efforts of Sabbath-school associations
\quellseite{192}
in different sections of the country; to strengthen the hands of the friends of pious instruction on the Lord’s Day; to disseminate useful information; to circulate moral and religious publications in every part of the land; and to endeavor to plant a Sunday School wherever there is a population.” The basis of this Union was declared to be “No sacrifice of principle essential to salvation; no compromise of duty; no interference with the internal management of smaller associations; all discordant elements must be banished, and (there must be) union with Christ, and union with each other.” Upon this basis, the essential truths of Christianity held in common by all evangelical denominations, were to be fully inculcated. The individual members of the A.S.S.U. could maintain the integrity of their denominational relations, but in this work they were to unite with members of other churches for the purpose of teaching “the truths that Christ taught and as plainly as he taught them.” The significance of the A.S.S.U. for the development of the Sunday School can hardly be overestimated. It published a vast amount of literature for pupils and teachers, prepared lesson plans and recommended them to all schools, issued Bibles, text books, and other lesson helps down to the most modest text cards, and promoted the organization of new schools with brilliant zeal and enthusiasm. At the meeting in 1830 it resolved “within two years to establish a Sunday School in every destitute place where it is practicable throughout the valley of the Mississippi”; at once 78 missionaries started out from Cincinnati and founded 2876 Sunday Schools in the Mississippi valley during the next two years. In 1832 the A.S.S.U. invited representatives of the S. S. cause from all parts of the country to attend a meeting in New York; thus the first National Convention was brought about (October 3). An exhaustive questionnaire which had been sent out furnished valuable data concerning the actual conditions and
\quellseite{193}
needs of the S.S. On the basis of this material, Cope summarizes the situation at that time in the following manner: “(1) There were several distinct types of schools; infant schools, mission schools, adult schools, evening schools, as well as the general schools. (2) Lesson plans were many. Some spent over an hour in class memorization of long passages of Scripture; others simply told Biblical and other stories; some assigned one verse for each day of the week, the whole being the basis of recitation and comment on Sunday; the tendency was to adopt the limited portions designed by the A.S.S.U., all schools aided by this society being required to use the same lessons; everywhere the emphasis was on Biblical material. (3) Nearly all schools were directed by superintendents. (4) Libraries of general literature were established in many schools.\footnote{No public libraries existed in the early part of the nineteenth century, nor was there any juvenile literature; President Humphrey of Amherst could not recall over six books in this class, even counting Webster’s \emph{Spelling Book}, \emph{Robinson Crusoe}, and \emph{Pilgrim’s Progress}! The first S.S. library was established at Boston in 1812; in 1830 the A.S.S.U. issued 200 books for S.S. children.} The type of book afterwards to be known as a “Sunday School” book was almost a special creation for them. (5) Special “Bible classes” for church members and “adult classes” were organized. (6) Many schools met on Sunday afternoon, the sessions often being two hours in length. (7) In the frontier regions schools were organized before churches and became the parents of the latter. (8) In the cities there were many mission schools not immediately connected with specific churches.” This first national convention also resolved, “that the S.S. should embrace all classes of the community” and not only destitute and neglected children. The second national convention met in Philadelphia in 1833. It urged the establishment of Sunday Schools in jails and reformatories and the promotion of groups for the study of the Bible in the homes.
\quellseite{194}
For several reasons no national convention was held during the ensuing twenty six years. One important reason was the reawakening of denominational consciousness in this period: the more vigorously this consciousness was aroused, the more inevitable was the conflict with an association which was avowedly un-denominational and virtually super-denominational, which professed to preach a gospel of Christ free from all denominational “stains and encumbrances.” and which arrogated to itself some of the Church’s inalienable functions. So it came about that the various denominations organized their own S.S. Societies. In 1827 the Methodist Episcopal Church created its S.S. Union, the Unitarians followed in the same year, the Lutherans (v.i.) in 1830, the Congregationalists in 1832, the Baptists in 1840 (resp. 1826), the Southern Baptists in 1857, the Presbyterians organized, in 1833, “The Board of Managers of the Presbyterian Tract and Sabbath School Book Society” which subsequently promoted the S.S. work of this denomination. The reawakening of confessional consciousness thus made these 26 years a period of battle against the super-denominational claims and tendencies of the A.S.S.U. The historians of the A.S.S.U. contemptuously designate this movement as “sectarian”; as a matter of fact, however, this opposition was not only justified, but clearly necessitated by the very nature of the Church. While the A.S.S.U. thus was seriously limited in many of its activities, it devoted itself with redoubled energy to the publication of periodicals, teachers’ manuals, and S.S. literature in general and also to the formation of Sunday Schools in frontier regions where no separate denominational schools had been organized. Of special importance was the publication of the \emph{Sunday School Journal} (since 1831) in which the pedagogic theories of such great educators as Sturm, Comenius, Milton, Francke, Rousseau, Pestalozzi, Froebel were set forth and the current systems of Lancaster, Bell, Gall,
\quellseite{195}
Stow, and Mimpriss were examined as to their applicability to the S.S. Naturally the method and the standards of teaching were greatly improved by these efforts; also the denominational schools profited thereby. Another tremendously important service rendered by the A.S.S.U. to the Sunday School and the cultivation of a Christian spirit in general was the creation of a juvenile literature of a religious type. According to E. W. Rice, the following principles governed this phase of the Union’s work: “(1) The literature must be suited to the development of the child mind; (2) pure in tone; (3) serious rather than sensational; (4) ingenious, but not absurd; (5) popular rather than polished; (6) thoroughly Biblical and evangelical; (7) truly American, written by American writers, statesmen, and philanthropists, redolent of American mountains, forests, prairies, rivers, history, and yet filled with the spirit of the Word, not of the world.” In 1845 the Union began to issue a one-hundred bound-volume (72 to 272 pages per volume) library for Sunday Schools at the low price of ten dollars. Three similar sets followed. Add to these publications the periodicals for children (\emph{The Youth’s Penny Gazette}, with full page illustrations, was issued from 1843), the hymnals for which men like Mason and Hastings composed melodies, the large number of new Sunday Schools founded by its representatives, and the arousing of missionary zeal (as early as 1835 the Union declared one of its objects was “to assist in carrying the Gospel to every family in the world, and to insure the religious instruction of every child that is born”)—and one readily sees what a rich stream of blessing proceeded from the A.S.S.U. even after its activity had been circumscribed by the newly founded denominational schools and societies. Nor should it be forgotten that many denominational societies adopted the lists of lessons issued by A.S.S.U. and explained them from their own point of view.
\quellseite{196}
In 1859 the third national convention was in session at Philadelphia. On this occasion also the western section of our country was well represented; the most prominent men besides Henry Clay Trumbull (Congregationalist, since 1871 secretary of the A.S.S.U. and chairman of the executive committee of the National S.S. Convention, soon afterward editor and part-owner of the \emph{Sunday School Times}, died 1903) were Benjamin Franklin Jacobs of Chicago (a Baptist business man, prominent S.S. leader in Illinois, conceived the idea of Uniform Lessons, since 1872 active member of the Lesson Committee, later president of the International S.S. Association, died 1902) and John Heyl Vincent of Joliet, Illinois (Methodist, one of the leaders in the movement for the International Uniform Lessons, helped found the Chautauqua Assembly (1874) where training courses for S.S. teachers were offered, Bishop since 1888). After the Civil War the fourth national convention met at Newark, N. J., in 1869; and at the fifth convention at Indianapolis, in 1872, the plan of Uniform Lessons was adopted, a Lesson Committee (composed of twelve members, later increased to fourteen by the addition of two Canadians) was appointed, and the national convention was made International with biennial meetings. In 1889 the first World Sunday School Convention was held in London, chiefly owing to Jacob’s efforts. At the fifth convention in Rome (1907) it changed its name to World’s Sunday School Association. The purpose, policy, and field of the Association were defined as follows: “That this Association shall hold conventions and gather information concerning the conditions of Sunday Schools throughout the world by correspondence, visitation, and other methods; that it shall seek to extend the work and increase the efficiency of Sunday Schools by cooperation with Sunday Schools and missionary organizations and otherwise, especially in those regions of the world most in need of help; that it shall seek
\quellseite{197}
to improve, so far as possible, the methods of organization and instruction in Sunday Schools and promote the formation of S.\,S. Unions and Associations.” By agreement the world field was divided for the purpose of financial administrative responsibility: Europe, Australia, South Africa, India—the British Section of the Executive Committee; North and South America, Japan, Corea, Philippines, Turkey, North Africa—the American Section of the Committee; China to be jointly administered by the British and American Sections. The chief feature of the seventh convention at Zuerich, Switzerland, in 1913 was the report of the Six Commissions of the Association on the Sunday School conditions, needs, and opportunities of the world, covering continental Europe, South Africa, India, the Orient, Latin America, and the Mohammedan lands. According to the latest available statistics there were in the United States 143,035 Protestant Sunday Schools with 17,729,216 pupils. To the World Association belonged in 1924 32,677,611 members in 347,001 schools with 3,520,192 teachers and officers.

Efforts are made to introduce everywhere within the International S.\,S. Association the identical lessons which are interpreted independently by denominational boards according to their convictions. In this way the autonomy of the cooperating denominations is preserved and at the same time expression is given to the unity of the Churches composing the Association—a compromise terminating the conflict between the superdenominational aspirations of the A.\,S.\,S.\,U. and the denominational demands. Most of the Churches have adopted these lessons, but there are many Sunday Schools numbered in the statistics of the I.\,S.\,S.\,A. (e.\,g. most of the Lutheran schools) that use their own systems. In order to do justice to the common tasks of all Sunday Schools (especially the improvement of methods) more fully than is possible at the International Conventions, “\emph{The Sunday}
\quellseite{198}
\emph{School Council of Evangelical Denominations}” was organized in 1910 which included in its membership about 150 S.\,S. editors, publishers, and secretaries representing all the more important evangelical denominations in the United States and Canada. The chief object of this Council was to apply the results of modern psychology, pedagogics, and sociology to the S.\,S. lesson courses without interfering with denominational autonomy in matters of administration. Its resolutions though helpful and salutary now and then, on the whole evinced unionistic tendencies shifting the emphasis from the center of the New Testament. This Council and the I.\,S.\,S.\,A. were merged into a new body in February 1922, “\emph{The International Sunday School Council of Religious Education}.”

In reference to the study material used by American Sunday Schools it may be said that in general the Bible formed the basis of instruction ever since the S.\,S. did not need to devote its time to the study of elementary branches. Only a few recent lesson plans depart from this practice, offering general materials (natural history, tales, daily life, etc.) which are to be interpreted and applied religiously. In regard to the method employed in introducing the pupils to the Bible, several periods stand out more or less distinctly. At first the memory method was deemed sufficient. The pupils memorized and recited shorter or longer sections of Scripture. Spurred on by rewards, some children would learn a hundred or more verses, or even a whole Biblical book in the course of one week. Now, to store in the mind a rich treasure of Bible passages is undoubtedly beneficial; but when the entire S.\,S. period, even when it lasted two hours, was consumed by the recitation of memory material, no time was left for explanation and application, and very little religious insight was brought about and still less influence was exerted upon the children’s religious life. A step of progress was taken
\quellseite{199}
when in 1825 Gall’s method was introduced in the Sunday Schools of America. Dr. James Gall of Edinburgh, Scotland, had prepared a series of Bible lessons on a method called “\emph{Nature’s Normal School}”; the lesson consisted of a brief Bible story to be told by the teacher and a series of questions on that story. The system, however, was so complicated and mechanical that little actual knowledge was imparted and the soul of the pupil remained cold. In 1829, Rev. Albert Judson projected for the A.\,S.\,S.\,U. a scheme of lessons to be completed in five years; all classes studied the same lesson and reviews were conducted at the end of the quarter. This plan was adopted by the A.\,S.\,S.\,U. and widely used during the following years. In distinction from the Uniform Lessons introduced in 1872, the period inaugurated by the acceptance of Judson’s Lessons may be designated as the \emph{First Uniform Lesson Period}. This method actually provided for the study of the Bible in the course of five years, called for the actual reading of the Bible, and stimulated the children to think about the stories for themselves. But, of course, everything of a denominational nature was omitted. Among the numerous helps for these lessons, the most popular were \emph{Judson’s Questions} in three grades, and especially the \emph{Union Questions} edited by John Hall (in five little volumes, later after an extension of the course in seven, and finally the review year having been added, in twelve). For beginners \emph{The Child’s Scripture Question Book} was issued, and for adult Bible classes McDowell’s and Tyng’s lessons. In point of method Cope’s judgment is: “The examination of a text book of 1830, with the assigned lesson, analysis of the narrative, series of questions, explanations of unfamiliar words, exposition on the doctrine and applications on practical lessons, shows that these ‘lessons’ were at least as good as those in the average ‘quarterly’ of sixty years later.” A great diversity of lesson plans and methods resulted from the organization
\quellseite{200}
of denominational S.\,S. associations for which reason the period following upon the first uniform lesson period is often designated as the \emph{Period of Babel Systems}. Many schools continued to use the Uniform Lessons; others devised their own; some sought to adapt the lessons to the age of the pupils (e.\,g., the Unitarians had a series of eight graded texts in 1852); Episcopal schools prepared lessons on the program of the church year. John H. Vincent published in his \emph{Sunday School Teacher} a course entitled “Two Years with Jesus.” The capable albeit somewhat passionate Edward Eggleston, in his \emph{National Sunday School Teacher} published another series of lessons in 1867 and in his \emph{Manual: a Practical Guide to Sunday School Work} showed with great skill how to use the plan. He was an outspoken foe of the Uniform Lesson (the same lesson for all classes on one Sunday) and was convinced that “the needs of schools call for a variety of lessons just as a day school needs more than one reader in Arithmetic.” At the national convention of 1869, however, the lack of uniformity was keenly felt and it seemed imperative to do something toward introducing uniform standards and methods in the S.\,S. At this juncture B.\,F. Jacobs of Chicago proposed his plan of Uniform Lessons for all Sunday Schools, the helps to be prepared by the denominational boards. Because the autonomy of the denominations was sufficiently safeguarded and at the same time the ideal of uniformity was, in a measure, approached, the Convention of 1872 with an overwhelming majority adopted this plan. Eggleston’s heroic opposition did not stem the tide (he declared that it was “a movement backward; it would pull down good schools, produce a dead level uniformity”; and in 1880 he said: “The system of rigid adherence to one lesson for all the school, combined with selections now and then of subjects fit only for a theological seminary, is not in accordance with practical wisdom”). An international lesson committee was
\quellseite{201}
appointed (see p. 174). Therewith the period of the \emph{International Uniform Lessons} began. For thirteen years the new lessons dominated the entire field; all opposition was quelled. When the preparation of a separate series of lessons for younger classes was demanded in 1885, the Lesson Committee refused to issue such lessons. The London S.\,S. Union in 1890 requested “a system of graded lessons, in which the same subjects shall be considered, but with different Scripture chosen first for the primary classes, and then for more advanced classes,” but the request received as little attention. Rev. Erastus Blakeslee even began to issue, at his own risk, a graded series of inductive Bible lessons which formed a six year course. The leaders of the I.\,S.\,S.\,A. were so enamored with the idea of uniform lessons that they did not react favorably to these suggestions and criticisms although it should have been evident that beginners and adults have special needs that can not be supplied by a mere gradation of the lessons helps, but only by different lesson texts and subjects. It is quite significant that in the following years the Primary teachers were the most persistent objectors to the Uniform Lessons and the most aggressive champions of a graded system. With increasing insistence they petitioned the Lesson Committee to issue suitable lessons for little children, especially since the interest in child-psychology began to grow in our country (at a conference in 1903 one of the speakers bluntly said: “The old education put the material first, the child second; the new education puts the child before the material!”). After many heated debates a conference of I.\,S.\,S.\,A. leaders finally adopted the following resolutions: “(1) That the system of a general lesson for the whole school which has been in successful use for 35 years, is still the most practicable and effective system for the great majority of the Sunday Schools of North America. Because of its past accomplishments, its present usefulness,
\quellseite{202}
and its future possibilities, we recommend its continuance and its fullest development; (2) that the need for a graded system of lessons is expressed by so many Sunday Schools and workers that it should be adequately met by the I.\,S.\,S.\,A., and that the Lesson Committee should be instructed by the next International Convention, to be held in Louisville, Ky., June 18–23, 1908, to continue the preparation of a thoroughly graded course covering the entire range of the Sunday School.” The convention at Louisville unanimously adopted this report. But hardly had the new lessons appeared when the Southern Baptists and the Southern Presbyterians justly attacked them for the “absence of doctrine, the presence of extra-Biblical lessons, the omission of many important topics and liberal interpretation of Scriptures.” These and other criticisms resulted in some slight improvements. Graded lessons have now been issued as follows: Beginners (four and five years of age) two years; Primary (six to eight years) three; Junior (nine to twelve years) four years; Intermediate (thirteen to sixteen years); Senior (seventeen to twenty-one years) four years. From 1908, accordingly, dates the period when alongside of the International Uniform Lessons the \emph{International Graded Lessons} exist by right in the I.\,S.\,S.\,A. Of the schools that have introduced graded lessons and consistently applied modern psychology and pedagogy (often, indeed, at the expense of Bible truth) and have exerted strong influence, for good and evil, we mention the Model Sunday School of Columbia University, New York City; University Congregational School of Hyde Park, and the Hyde Park Baptist Sunday School, both of Chicago and in close touch with the University of Chicago.

\trenner

The Lutheran congregations in America were in the early period completely dependent on the catechetical work of the home Church. In his plea for help of May 31, 1693,
\quellseite{203}
Springer also asked the home Church in Sweden for 300 catechisms. The first German congregations were closely connected with the pietistic circles in the homeland. As a case in point, Boltzius of Charleston reports in 1736 how he explained the catechism in the daily evening service; we read the following entry: “Last evening we reviewed the seventh commandment and catechized upon the eighth. The consciences of some in the congregation were deeply stirred. In the exposition of the commandments to simple hearers we find it highly necessary to indicate the sins forbidden and the virtues enjoined therein not only generally, but they must be clearly specified according to the circumstances of the hearers.” Muehlenberg’s fidelity in the instruction of the neglected youth is well known. With more regularity than was possible for Muehlenberg who had to travel about extensively, Brunnholtz whose infirmity kept him closely to his parish, conducted \emph{Kinderlehre} or “Children’s Instruction.” The younger children he instructed in the Small Catechism, and the older ones and the servants in the “\emph{Order of Salvation}” and Bible History, an arrangement which clearly indicates Francke’s influence. It is to be regretted that, later, Rationalism entering from Germany and other countries made inroads upon American church life and affected also the catechetical text books. Quite undisguised it appears in F.\,G. Quitman’s \emph{Evangelical Catechism} (Hudson, 1812, or possibly 1814). Opposition to this perversion was at times rather feeble, but it never subsided altogether and kept alive the Lutheran confession in catechetical instruction.

In 1787 Professor John Caspar Velthusen of Helmstaedt prepared and edited the \emph{Lehrbuecher fuer die Jugend in Nord Carolina} (Text-books for the youth of N.\,C.) at the request of Nuessmann who had been sent to care for the Lutherans in Mecklenburg county in North Carolina. The purpose was to publish (1) a catechism, (2) a book of questions on the catechism, (3) a Bible handbook for everyone, (4) a selection of Bible stories with a brief history of religion, (5) a
\quellseite{204}
book of useful information for every day, (6) a manual of civics, (7) a manual of geography. At least the first two of these projected books appeared and are extant, viz., \emph{Helmstaedtischer Katechismus oder Christlicher Religionsunterricht nach Anleitung der heiligen Schrift}, and \emph{Fragebuch fuer Eltern und Lehrer, oder Anleitung zu Fragen und Gespraechen ueber den Katechismus mit Ruecksicht auf die Verschiedenheit der Faehigkeiten und des Alters der Jugend}. In his catechism which is a worth-while selection of Bible passages, Velthusen adheres to the doctrine of the Trinity, the eternal Sonship of Christ, the virgin birth, and the vicarious atonement, but in his statements concerning the appropriation of salvation he pays tribute, in terminology and subject matter, to the prevailing rationalism.—To this period belong also the following catechetical treatises written from widely diverging viewpoints: H.\,E. Muehlenberg, \emph{A Companion to the Catechism, or a Course of Instruction in the Christian Religion, for the Benefit of the Young}. A. Wackerhagen, \emph{Inbegriff der Glaubens- und Sittenlehre}, Philadelphia, 1804. Anon., \emph{Anfangsgruende der Religion, oder Katechismus fuer kleine Kinder}, Lebanon, 1814. G. Vorberg, \emph{Handbuch fuer die Konfirmierten}, New York; in English translation with the title \emph{Way Marks for the Confirmed}, 1815. Ph.\,F. Mayer, \emph{Instructions in the Principles and Duties of the Christian Religion, for Children and Youth}, 1816. Paul Henkel, Deutscher und englischer Katechismus. J.\,G. Lochmann, Hauptinhalt der christlichen Lehre. Id., Evangelical Catechism, Harrisburg, 1882. E.\,L. Hazelius, Materials for Catechization on Passages of Scripture, 1823.

The educational institutions available in the Lutheran Church were the parish and catechumen schools (the Pennsylvania Synod alone had as late as 1812 one hundred and sixty parish schools) and, in many congregations, instruction of the young on Sundays (Christenlehre). When the Sunday School, in its new form, was introduced in America, the Lutherans, too, soon adopted it as an additional agency of instruction. German Lutheran congregations were the pioneers in this respect. On May 7, 1804 a widow organized a S.\,S. of five or six children in St. Michael’s and Zion’s Congregations in Philadelphia; St. Paul’s Church followed in 1805. The North Carolina Synod in 1811 and again in 1812 called on the congregations to form Sunday Schools.
\quellseite{205}
In 1819 Christ Church, York, Pa., and Zion’s Church, Harrisburg, Pa., organized Sunday Schools. Adam Keller’s efforts resulted in the formation of a S.\,S. in St. John’s Church at Philadelphia in 1821 which on the opening day had an attendance of 19 teachers and 60 children in the forenoon and 100 children in the afternoon. Schools came into being at Lancaster in 1828 and at Reading in 1829. Progress was slow chiefly for the reason that the considerable number of day schools still in existence lessened the need for the Sunday School in its new form. The Pennsylvania Synod had 26 schools in 1838, in 1844 already 73. Six years after the formation of the American Sunday School Union, the General Synod organized a Lutheran Sunday School Union (1830). Influential men sponsored the S.\,S. cause, e.\,g., S.\,S. Schmucker (\emph{Plea for the Sabbath School System}, delivered at the anniversary of the Gettysburg S.\,S., 1830), S.\,A. Seiss (\emph{Thoughts on Education}, address to the Sabbath School Union, Cumberland, Md., 1850), and Benjamin Kurtz (\emph{Sermon on Sabbath Schools}). When it became evident that the parish school would not thrive in the Church as far as she used the English language, and when the inadequacy of pre-confirmation catechization was recognized, the S.\,S. became a practical necessity. “Father” F.\,C. Heyer, best remembered as America’s pioneer Lutheran missionary to India, who in his youth had been parish school teacher in Southwark, traveled very extensively as the agent of the Lutheran S.\,S. Union and established a large number of Sunday Schools in the thirties of the last century. More rapid progress was made since about 1840, the year in which such large congregations as St. Matthew’s and St. James’ in New York established schools. In German quarters it was especially Brobst who promoted this work in view of the horrifying ignorance of many children even in church homes. He issued a booklet for the instruction of beginners, and also showed
\quellseite{206}
how to open and to close the Sunday School session liturgically. In 1854 he was appointed, together with Vogelbach and B.\,M. Schmucker, to address a circular letter to the congregations of his Synod in the interest of the S.\,S. He accomplished his task so perfectly that an English translation of this letter was published as late as the eighties. Lutheran hymnals and lesson leaves introducing into Scripture and evangelical doctrine were issued, and the Catechism was adapted to S.\,S. use. As long as the instruction of confirmands was faithfully performed, Luther’s Small Catechism was used which had been translated into English long ago (of special merit are the translations made by Ph.\,F. Mayer, 1816; by C.\,F. Schaeffer, made in 1895 by order of the Pennsylvania Synod; and the Joint Committee’s translation, 1899); also new catechisms were produced. Of the expositions of Luther’s Catechism, Mann and Krotel’s work stands out as especially valuable (issued 1863 by order of the Ministerium of Pennsylvania, German and English, based on Caspari’s exposition). The \emph{Augsburg Sunday School Teacher} served as a suitable teacher’s manual; and periodicals were issued for children and for Sunday Schools. H.\,J. Schmidt (then professor in Gettysburg College and Seminary) published a \emph{History of Education} in 1842 which contained in part “a plan of culture and instruction based on Christian Principles, and designated to aid in the right education of youth, physically, intellectually, and morally”; it was the only work of the kind which the American Lutheran Church had produced till then. And it remained the only work until H. Ziegler’s \emph{Catechetics, Historical, Theoretical, and Practical} appeared in 1873. Sunday School teachers’ meetings are first mentioned in 1852; B.\,M. Schmucker conducted such meetings weekly for the purpose of studying the lesson for the following Sunday.

Of the catechetical works that appeared during this period we mention the following: John G. Morris, \emph{Catechumen’s and Communicant’s}
\quellseite{207}
\emph{Companion}, Baltimore, 1831. Id., \emph{Catechetical Exercises on Luther’s Catechism}, adapted from the German, Baltimore, 1832. S.\,W. Harkey, \emph{Lutheran Sunday School Question Book, or a Help to the Systematic Study of Sacred Scriptures}, compiled from the German, Frederick, Md., 1838. John G. Morris, \emph{Popular Exposition of the Gospels for Families, Bible Classes, and Sunday Schools}, 2 vols., Baltimore, 1840. Benjamin Kurtz, \emph{The Serial Catechism, or Progressive Instruction for Children, adapted to the growth in grace and knowledge}, Baltimore, n.\,d. Id., \emph{Easy Catechism for Young Children}, 1843. John G. Morris, \emph{Luther’s Catechism Illustrated by Additional Questions and Answers}, 1844. C.\,A. Morris, \emph{An Easy Catechism for Young Children}, n.\,d. C. Phil. Krauth, \emph{Lutheran Sunday School Hymn Book}, about 1855. S.\,A. Seiss, \emph{Training Little Ones for Christ}, Baltimore, 1853. Clemens Miller, \emph{Luther’s Catechism}, Baltimore, 1857. M. Sheeleigh, \emph{Hymns, Selected and Original, for Sunday Schools of the Evangelical Lutheran Church, Revised, with Appendix, containing Hymns for the Use of Infant Schools}, Baltimore, 1860. T.\,T. Titus, \emph{Sunday School Question Book}, n.\,d. S.\,S. Schmucker, \emph{Evangelical Lutheran Catechism, Designed for Catechumens and the Higher Classes in Sabbath Schools}, Baltimore, 1850. W.\,I. Mann and G.\,F. Krotel, \emph{Luther’s Small Catechism, Explained in Questions and Answers for the Use of the Church, School and Family}, 1863. M. Sheeleigh, \emph{Outlines of Old Testament History}, also \emph{Outlines of New Testament History}, about 1865. E.\,I. Koons, \emph{Questions on Luther’s Catechism}, 1869. G.\,F. Jager, \emph{Katechismus der christlichen Lehre in Fragen und Antworten}, Kutztown, 1833. H.\,W. Scriba, \emph{Anfangsgruende des Christentums}, translated from the French, Chambersburg, 1834. J.\,G. Schmauk, \emph{Erstes Buch fuer deutsche Schulen}, Philadelphia, 1844. F.\,P. Peixoto, \emph{Leitfaden, wonach der kleine Katechismus Luthers erklaert wird}, Sunnytown, Pa., 1845. S.\,K. Brobst, \emph{Fragebuechlein fuer die Anfangsgruende der Religion}, 1846. Id., \emph{Gebete fuer Sonntagsschulen}, 1847. Id., \emph{Fragen und Antworten aus dem Neuen Testament}, 1849. C.\,A. Morris, \emph{Liederbuechlein fuer Kinder}, gesammelt und herausgegeben von den Lehrern der deutschen Sonntagsschulen in Yorktown, 1849. F.\,C. Wyneken, \emph{Spruchbuch zum kleinen Katechismus Lutheri}, Baltimore, 1849. S.\,K. Brobst, \emph{Gesangbuch fuer die Sonntagsschulen der ev.-luth. und deutsch reformierten Kirchen}, Allentown, Pa., 1853. J.\,C. Haas, \emph{Luthers kleiner Katechismus mit Spruechen aus der Heiligen Schrift}, Philadelphia, 1856. Mann u. Krotel, \emph{Luthers Katechismen mit Fragen und Antworten}, 1863. C.\,F. Spring, \emph{Katechetische Unterweisung zur Seligkeit}, New York, 1868.

\quellseite{208}
When the lessons issued by the International Sunday School Association began to appear in 1872, they were introduced more or less universally in the English Lutheran synods of the East, especially in the General Synod where they were explained in the Augsburg Teacher and in the church papers. Not even the General Council was altogether proof against them. Quite generally catechumen instruction was retained as a supplement for the Sunday School. The text books chiefly used were Mann and Krotel, \emph{Luther’s Small Catechism}; also Seiss, Spaeth, and Jacobs, \emph{Luther’s Small Catechism}, with Scripture Texts, by authority of the General Council in America. Philadelphia, 1883; or the Stohlmann reprint of Michael Walther’s Catechism as remodeled by Luehrs (the English translation by H.\,E. Jacobs as well as the German edition were widely used). Other serviceable books were Wischan and Spaeth, \emph{Mein erstes, zweites, drittes Sonntagsschulbuch}; the Biblical History published by the Pilger Book Store of Reading; F. Greenwald, \emph{Questions on the Gospels and Epistles of the Church Year}, Lancaster, 1873. 1874; Ludwig’s abovementioned catechism with proof texts and an appendix containing The Order of Salvation; \emph{Sunday School Book of the General Council} (English 1873, German 1897, revised 1906); Wonneberger, \emph{Sonntagsschulharfe}.

The inadequacy and the imperfections of Sunday School instruction as based upon the International Uniform Lessons were keenly felt in the General Council. In the first place, the I.\,S.\,S.\,A. encourages the view that the S.\,S. is a separate, independent institution alongside the Church; furthermore, the fact that the children are baptized is not properly recognized; little, if any, real knowledge of the history of salvation is imparted; law and gospel are so often confounded; the central facts of salvation are not sufficiently clearly brought out; no instruction in the catechism being provided for, the real goal of instruction is obscured; the appreciation
\quellseite{209}
of specific catechumen instruction is suppressed rather than aroused; and finally, an organic connection of the instruction of the confirmed with the lessons learned during their pre-confirmation instruction is rendered impossible. Moreover, the Uniform Lesson System does not take sufficient account of some fundamentals of psychology and pedagogics. It was worthy of the highest praise, therefore, when the General Council in the face of severe opposition decided in 1895 to supply the urgent need of a typically Lutheran and pedagogically sound lesson system by publishing its own system. A more capable and purposeful editor than Dr. Theodore Schmauk could hardly have been found. It is true that a graded system had been advocated before (cf. the preceding sketch of the development of the S.\,S.), and some preliminary work had been done; but this was the first time that a Lutheran S.\,S. system was created and that a whole church body united in this sort of enterprise. The system is divided into Primary, Intermediate, and Senior Departments. The Primary Department is divided into Kindergarten, Children’s, and Junior Department (from the third to the ninth year); \emph{Wonderland and Sunbeams}, \emph{Workland and Sunshine}, \emph{Pictureland and Sunrays} are the lesson helps for this department. \emph{In Mother’s Arms} and \emph{At Mother’s Knees} are preliminary to the whole, serving the mothers in instructing their little ones at home. The Intermediate Department is composed of seven grades (from the tenth to the sixteenth year); teaching aids to be used are \emph{Bible Story, Bible Readings, Bible History, Bible Facts and Scenes, Bible Biography, Bible Teachings, Bible Literature}. Instruction in the Catechism with confirmation is to occur between Bible Teachings and Bible Literature, or rather run parallel with these grades. The Young People’s Department (17 to 20 years of age) and the Adult Department (21 years and above) devote themselves, with the aid of the \emph{Senior Lesson Book}
\quellseite{210}
or Commentary, to the study of Scripture proper. The General Synod is still connected with the International S.\,S. Association, using its Uniform, or Graded Lessons; but in 1911 this body authorized its S.\,S. Committee to enter into negotiations with other English speaking bodies for the purpose of creating in cooperation with them a specific Lutheran Sunday School literature. In the East there is used for the instruction of confirmands, in addition to the literature already mentioned, the \emph{Cathechism} by the author of this text book on Catechetics; Loehe’s catechism, translated by Horn; and the catechisms by Trabert and Stump. On the basis of Kaftan’s work (cf. § 19), I.\,W. Horine, in his \emph{Catechist’s Handbook}, Philadelphia, 1909, has given the catechist a valuable aid to the understanding of Luther’s Catechism.

In the third decade of the last century German, and a little later also Scandinavian, immigration into the western part of our country commenced on a large scale. The newly founded congregations and church bodies had brought some agencies of religious instruction with them from their mother countries, others they had to devise for themselves. Catechumen instruction and Christenlehre they brought with them; among the new agencies, devised in this country, is the parish school. In the public state-supported school which every child in the home country must attend, religious instruction, indeed, occupied a prominent place, but the school was a state school, not a parochial school; and if the Church exerted any influence upon religious instruction as imparted in the public school, it was only by reason of being a State Church. In their new home, however, the immigrants found Church and State—fortunately—separated by constitutional enactment; the State, therefore, could not possibly impart religious instruction in school. But for the sake of her self-preservation the Church could not dispense with religious instruction of her youth; and since education ought to be a harmonious
\quellseite{211}
unit, the spirit of Christ penetrating the sum of secular knowledge, there arose for these new congregations and church bodies the necessity of establishing independent church, or parish schools. In these schools they gathered their children and endeavored to train them by graded instruction in religious and secular branches to become efficient members of the Church and of the State. It was principally Loehe who advocated the formation of such schools; he also collected funds for the first Lutheran Normal School (Saginaw City, Mich., 1852) and sent the teachers and students for this school. The Synod of Iowa founded by Loehe’s disciples has never quite forgotten this task; but the course of developments caused the Missouri Synod to become the chief exponent and champion of the parish school ideal. Missouri’s care for the parish school is and remains the most splendid chapter in its history. In Lindemann’s \emph{Schul praxis} its principles of school management were laid down; its pedagogical organ, \emph{Schulblatt}, edited by Lindemann, has outlived all similar periodicals; and in the preparation of schoolbooks and other helps it has usually been guided by sound pedagogical principles. It is gratifying that the Missouri Synod at an early stage of its development not only realized the necessity of cultivating the language of the country in its schools, but even of founding purely English parish schools. In 1926 Missouri had in its two Normal Schools (River Forest, Ill., and Seward, Neb.) 630 students. Throughout the Missouri Synod there are 1388 parish schools with 1272 male teachers and 447 female teachers and 401 pastors teaching school. The pupils number 80,173. There remain to be mentioned the Normal School at Waverly (Iowa Synod), and at New Ulm (Wisconsin and Minnesota Synods). The Scandinavians have normal courses in connection with several of their institutions. Where conditions are adverse to the formation of complete parish schools, faithful congregations organize “partial” parish schools (cf. § 31); for instance,
\quellseite{212}
the pastor conducts a winter or confirmation school; there is frequently a summer school term of two months and Saturday School; and almost universally we find the Sunday School. The last named institution has made its way also in the West: at present every synodical body has a special Sunday School Committee which provides for the needed Sunday School literature. After years of purposeless efforts sound principles are gradually beginning to find recognition. The Sunday School systems of the Norwegians and of the Iowa and the Missouri Synods deserve specific mention.

It remains to note, briefly, the catechetical literature devised by the western Synods. Stephan and Walther introduced from Saxony the \emph{Dresden Cross Catechism} composed in 1688 by the ministerium of the Cross Church in Dresden. This was later generally replaced in the Missouri Synod by the large and small catechism of C. Dietrich (first edition 1613). The disciples of Loehe, in the Iowa Synod particularly, disseminated his catechism (first edition 1845) which had been composed with a special view to the Lutherans in America. Also the Swedes, Danes, and Norwegians made use of the catechisms of their respective homelands, e.\,g., of that of Pontoppidan. At present almost every Synod has its own synodical catechism although some trans-atlantic catechisms are still in use. Almost invariably they are bilingual. The catechism of the Missouri Synod has been edited by Schwan; that of the Wisconsin Synod was based upon the \emph{Dresden Cross Catechism}, but has been replaced by one edited by Gausewitz; that of Ohio is based upon Dietrich’s, but is about to be revised. Since 1904 also the Iowa Synod has its own catechism, edited by the author of this textbook on Catechetics. It differs from other catechisms, as to contents, in that it returns to the expository principle of Loehe (1845) and Tetelbach (1568); it excludes all supplementary material taken from dogmatics or sacred history, and it restricts
\quellseite{213}
itself to those subjects which are essential to a Christian life; as to form, in that it adopts the thetic, not erotematic, presentation. Since 1917 it is to be found, in both German and English, also in abbreviated size and in the form of questions and answers. The Augustana Synod, likewise, has its own explanation; the Norwegian Synod uses largely Dietrich’s Catechism while the United Norwegian Church employs the explanation by Sverdrup.

Among the textbooks on Biblical History those of Reading (Pilger Book Store, Reading, Penn.), the old Huebner Bible History (1714) and those issued by the Missouri Synod enjoy a large circulation. Among those in English, that published by the Augustana Synod deserves special mention; the latest, with several novel features, is by the author of this textbook. The book \emph{How I Tell the Bible Stories to My Children}, (second, revised edition, two volumes, 1926) by the same author, is designed to furnish a contribution to the art of story-telling as applied to Biblical History. Among aids for the catechist there are helps for the instruction in Biblical History by Wegener and by Simon in German, and by Rupprecht in English; for instruction in the Catechism by Stellhorn in German and Schuh in English. We mention also the writings of Schaller and Reu on instruction in the Bible (both English). In the Missouri Synod Quarterly Professor Dau has published material relative to instruction in the Catechism. Among theoretical dissertations there are to be mentioned, in addition to Lindemann’s Schul praxis, especially Grossmann’s \emph{Gemeindeschule}, Herzer’s \emph{Katechetik}, and Gerberding, \emph{The Lutheran Catechist}; also the book of the Methodist Hertzler contains some good sections. A sourcebook by the author in ten volumes is evidence that the Lutheran Church in America has taken part in historic research work in the field of Catechetics.

\trenner

\quellseite{214}
The efforts to introduce \emph{moral and religious instruction into the public schools} can not and shall not receive notice here. The publications of the Religious Education Association (headquarters in Chicago) contain considerable material in regard to this subject, as in regard to the religious training of our people in general. Some of the efforts, however, to complement public school education by religious day schools ought to be mentioned. What Charles De Garmo says in \emph{Principles of Religious Education} (p. 63): “From the standpoint of the development of religious intelligence the American (public school) system must be pronounced the most fragmentary, partial, inefficient, haphazard system in the world,” is more and more recognized as true; likewise, that the work of the Sunday School cannot be recognized as an adequate complement. G. Boville, of New York City, has briefly put the matter as follows: “The Church has enrolled about 15,000,000 children of school age in Sunday Schools, for whom are provided fifty-two half hours of instruction; there are 10,000,000 more children of school age not enrolled in any Sunday School.” In order to give the Sunday School pupils more religious instruction, and the others at least a little, this man gathered during the summer vacation the children of an East side precinct in New York City every day for six weeks in such churches as were opened to him and engaged college students to supervise and instruct them, and thus “brought together idle children of the streets, idle buildings of the churches, and idle students of the colleges.” Instruction, in part at least, was of a religious character “on broad non-sectarian lines.” In view of the success of Boville, the National \emph{Daily Vacation Bible School} Association was organized in order to arouse interest in this matter also at other places. In recent years the movement has grown tremendously. The association was made international and recently merged with the International Council of Religious
\quellseite{215}
Education. It is not saying too much that “a Vacation Bible School is now regarded as an essential department of every modern church program.” Already before this, another method had been tried in a small way. As early as 1898, H.\,R. Vaughn, pastor of a Congregational Church in Elk Mound, Wis., organized vacation religious day schools (two weeks) in conjunction with teacher training institutes which he conducted in the country and the smaller towns. He graded his schools on the basis of the public school and demanded credentials similar to those asked for by the public school boards, and with the help of these teachers sought to give the children instruction in Biblical History and in the history of missions and of the Church, and thereby to arouse in them enthusiasm for everything noble and good. Such schools — no propaganda for their extension was made — are conducted in a dozen villages and towns of Wisconsin, but also in some larger cities, e.\,g., Eau Claire, Madison, Beaver Dam, and Rockford, Ill. Professor W.\,J. Mutch of Ripon College is still the soul of the movement. In the summer of 1918 the Lutheran churches of Madison, Wis., conducted a Bible school which in some respects conformed to the movement fathered by Vaughn, but was governed by Lutheran principles. A kindred plan is prosecuted in Colorado, the so-called Greeley Plan of Bible Study. By virtue of cooperation between the State Teachers College of Colorado at Greeley and the churches, courses in religion given in the churches and approved by the College are accepted for credit. In 1907 the Lutheran pastor G.\,U. Wenner, as chairman of a committee appointed by the Federal Council of Churches of Christ in America, proposed another plan. He desired that the legislature “should authorize all children in the public schools to attend, at the option of their parents, on Wednesday afternoons, either the public schools or schools of religion in their own churches.” The salient point here
\quellseite{216}
is confessional instruction; and undoubtedly this plan marks a considerable advance, but the amount of instruction which it provides is still inadequate. — Confessional religious instruction during the week is made possible also by the \emph{Gary Plan} (as originally conceived) which W.\,E. Wirt, superintendent of the public schools of Gary, Ind., began to carry out in 1913. His desire was to care well for the child and not merely to give him formal instruction. He wanted more of the child’s time, so that the child might live most of his waking hours under the supervision of the school. “On this theory the Gary schools were conducted from 8:15 A.M. to 4:30 P.M. for six days in the week, and the school undertook to provide for the child’s play as well as his welfare in other lines. During these additional hours, Mr. Wirt thought the children should be sent in small parties to receive the benefit of any welfare agency, and among these stood the Church. Mr. Wirt informed the various churches in Gary that he stood ready to send to them for week-day instruction any and every child whose parent should sign a card desiring to have him so sent, and that the school would allot the church from one up to six hours per week of each child’s time, according as the church felt capable of using such time. This hour included the going from and returning to school. There was no attempt to control or supervise the use which the church made of this time, and no credits were given for it. Children not so sent remained in the school and were occupied in some manner not counting in their formal studies. The parent then, in signing the church card, chose for his child an elective course in religious subjects, but, in taking this elective the child lost nothing in his formal studies.” This original Gary Plan has undergone some radical changes, most serious among which is the introduction of a cooperative arrangement patterned after the ideals of the Religious Education Association. This affects, of course,
\quellseite{217}
the curriculum and robs all instruction of its confessional character. In the words of the Superintendent of the Community Church Schools, Gary’s goal is “(1) a week-day religious center beside each public school; (2) an opportunity for every child who desires it to have at least one hour of religious instruction during the week; (3) a sufficiently large faculty to provide for small classes; (4) a correlated program of activities and instruction which shall include all the agencies interested in the welfare of the child.” It is evident that not much of positive religious value will be achieved where such a plan is in operation. The growth of the Week Day School movement was quite slow at first, and for several years Gary was the only school in existence; but since 1920 increasing numbers of such schools have been started: 77 in 1920; 131 in 1921, etc.) so that there are over one thousand Week Day schools in operation. It must be said, however, that many schools have passed out of existence soon after their birth and that others diminish in size. \emph{Four types of Week Day schools} are commonly distinguished: (1) the denominational or individual church type which, by the way, is out-stripping all the other types; here the school is governed by the local church and has no connection with any other church or church school in the community; (2) the denominational-cooperating type; here several schools of different denominations cooperate in the form of a loose federation for the promotion of common interests, such as securing public school time, granting credit, advertising, but each church school has its own governing board, its own course of study, and cares for its own children and such others as are not reached by other agencies; (3) the neighborhood or city system; here the schools lose their denominational character completely: the course of study is uniform for all denominations joining in the system, the governing board is made up of representatives of all cooperating churches, and the teachers
\quellseite{218}
are chosen without reference to denominational lines; (4) the pure community type where a board of religious education chosen—not by the churches, but by the community at large, governs the course; at present very few, if any, schools of this type are in operation. It is evident that Lutherans can not take part in any Week Day program similar to types 3 or 4, we can not surrender the choice of subject matter or the selection of teachers, to any agency outside of our own church. Unquestionably Week Day instruction of the individual church type or of the denominational-cooperative type (\# 2) will help remedy the deplorable situation in which Christian education finds itself at the present time. But even though the public school authorities should be found willing to accede to our requests for public school time, one fault would still remain: instruction in the secular branches would continue untouched by the Spirit of Christ; and two views of the world, in conflict with each other, would, during these years of development, compete for mastery in the soul of the child, without the possibility of an adjustment.

